% This LaTeX document needs to be compiled with XeLaTeX.
\documentclass[10pt]{article}
\usepackage[utf8]{inputenc}
\usepackage{amsmath}
\usepackage{amsfonts}
\usepackage{amssymb}
\usepackage[version=4]{mhchem}
\usepackage{stmaryrd}
\usepackage{hyperref}
\hypersetup{colorlinks=true, linkcolor=blue, filecolor=magenta, urlcolor=cyan,}
\urlstyle{same}
\usepackage{graphicx}
\usepackage[export]{adjustbox}
\graphicspath{ {./images/} }
\usepackage{caption}
\usepackage{bbold}
\usepackage{multirow}
\usepackage{fvextra, csquotes}
% \usepackage[fallback]{xeCJK}
\usepackage{polyglossia}
\usepackage{fontspec}
\usepackage{newunicodechar}
\IfFontExistsTF{Noto Serif CJK TC}
{\setCJKmainfont{Noto Serif CJK TC}}
{\IfFontExistsTF{STSong}
  {\setCJKmainfont{STSong}}
  {\IfFontExistsTF{Droid Sans Fallback}
    {\setCJKmainfont{Droid Sans Fallback}}
    {\setCJKmainfont{SimSun}}
}}

\setmainlanguage{english}
\IfFontExistsTF{CMU Serif}
{\setmainfont{CMU Serif}}
{\IfFontExistsTF{DejaVu Sans}
  {\setmainfont{DejaVu Sans}}
  {\setmainfont{Georgia}}
}

 Finite-Sample Properties of OLS }
%New command to display footnote whose markers will always be hidden
\let\svthefootnote\thefootnote
\newcommand\blfootnotetext[1]{%
  \let\thefootnote\relax\footnote{#1}%
  \addtocounter{footnote}{-1}%
  \let\thefootnote\svthefootnote%
}
%Overriding the \footnotetext command to hide the marker if its value is `0`
\let\svfootnotetext\footnotetext
\renewcommand\footnotetext[2][?]{%
  \if\relax#1\relax%
    \ifnum\value{footnote}=0\blfootnotetext{#2}\else\svfootnotetext{#2}\fi%
  \else%
    \if?#1\ifnum\value{footnote}=0\blfootnotetext{#2}\else\svfootnotetext{#2}\fi%
    \else\svfootnotetext[#1]{#2}\fi%
  \fi
}
\newunicodechar{×}{\ifmmode\times\else{$\times$}\fi}
\newunicodechar{⋮}{\ifmmode\vdots\else{$\vdots$}\fi}
\newunicodechar{⇔}{\ifmmode\Leftrightarrow\else{$\Leftrightarrow$}\fi}
\newunicodechar{¥}{\ifmmode\text{\yen}\else\yen\fi}
\title{Chapter 1: Finite-Sample Properties of OLS}

\author{}
\date{}

\begin{document}
\maketitle

\section*{Contents}
List of Figures ..... xyii\\
Preface ..... xix\\
1 Finite-Sample Properties of OLS ..... 3\\
1.1 The Classical Linear Regression Model ..... 3\\
The Linearity Assumption ..... 4\\
Matrix Notation ..... 6\\
The Strict Exogeneity Assumption ..... 7\\
Implications of Strict Exogeneity ..... 8\\
Strict Exogeneity in Time-Series Models ..... 9\\
Other Assumptions of the Model ..... 10\\
The Classical Regression Model for Random Samples ..... 12\\
"Fixed" Regressors ..... 13\\
1.2 The Algebra of Least Squares ..... 15\\
OLS Minimizes the Sum of Squared Residuals ..... 15\\
Normal Equations ..... 16\\
Two Expressions for the OLS Estimator ..... 18\\
More Concepts and Algebra ..... 18\\
Influential Analysis (optional) ..... 21\\
A Note on the Computation of OLS Estimates ..... 23\\
1.3 Finite-Sample Properties of OLS ..... 27\\
Finite-Sample Distribution of b ..... 27\\
Finite-Sample Properties of $s^{2}$ ..... 30\\
Estimate of $\operatorname{Var}(\mathbf{b} \mid \mathbf{X})$ ..... 31\\
1.4 Hypothesis Testing under Normality ..... 33\\
Normally Distributed Error Terms ..... 33\\
Testing Hypotheses about Individual Regression Coefficients ..... 35\\
Decision Rule for the $t$-Test ..... 37\\
Confidence Interval ..... 38\\
$p$-Value ..... 38\\
Linear Hypotheses ..... 39\\
The $F$-Test ..... 40\\
A More Convenient Expression for $F$ ..... 42\\
$t$ versus $F$ ..... 43\\
An Example of a Test Statistic Whose Distribution Depends on $\mathbf{X}$ ..... 45\\
1.5 Relation to Maximum Likelihood ..... 47\\
The Maximum Likelihood Principle ..... 47\\
Conditional versus Unconditional Likelihood ..... 47\\
The Log Likelihood for the Regression Model ..... 48\\
ML via Concentrated Likelihood ..... 48\\
Cramer-Rao Bound for the Classical Regression Model ..... 49\\
The $F$-Test as a Likelihood Ratio Test ..... 52\\
Quasi-Maximum Likelihood ..... 53\\
1.6 Generalized Least Squares (GLS) ..... 54\\
Consequence of Relaxing Assumption 1.4 ..... 55\\
Efficient Estimation with Known V ..... 55\\
A Special Case: Weighted Least Squares (WLS) ..... 58\\
Limiting Nature of GLS ..... 58\\
1.7 Application: Returns to Scale in Electricity Supply ..... 60\\
The Electricity Supply Industry ..... 60\\
The Data ..... 60\\
Why Do We Need Econometrics? ..... 61\\
The Cobb-Douglas Technology ..... 62\\
How Do We Know Things Are Cobb-Douglas? ..... 63\\
Are the OLS Assumptions Satisfied? ..... 64\\
Restricted Least Squares ..... 65\\
Testing the Homogeneity of the Cost Function ..... 65\\
Detour: A Cautionary Note on $R^{2}$ ..... 67\\
Testing Constant Returns to Scale ..... 67\\
Importance of Plotting Residuals ..... 68\\
Subsequent Developments ..... 68\\
Problem Set ..... 71\\
Answers to Selected Questions ..... 84\\
2 Large-Sample Theory ..... 88\\
2.1 Review of Limit Theorems for Sequences of Random Variables ..... 88\\
Various Modes of Convergence ..... 89\\
Three Useful Results ..... 92\\
Viewing Estimators as Sequences of Random Variables ..... 94\\
Laws of Large Numbers and Central Limit Theorems ..... 95\\
2.2 Fundamental Concepts in Time-Series Analysis ..... 97\\
Need for Ergodic Stationarity ..... 97\\
Various Classes of Stochastic Processes ..... 98\\
Different Formulation of Lack of Serial Dependence ..... 106\\
The CLT for Ergodic Stationary Martingale Differences Sequences ..... 106\\
2.3 Large-Sample Distribution of the OLS Estimator ..... 109\\
The Model ..... 109\\
Asymptotic Distribution of the OLS Estimator ..... 113\\
$s^{2}$ Is Consistent ..... 115\\
2.4 Hypothesis Testing ..... 117\\
Testing Linear Hypotheses ..... 117\\
The Test Is Consistent ..... 119\\
Asymptotic Power ..... 120\\
Testing Nonlinear Hypotheses ..... 121\\
2.5 Estimating $\mathrm{E}\left(\varepsilon_{i}^{2} \mathbf{x}_{i} \mathbf{x}_{i}^{\prime}\right)$ Consistently ..... 123\\
Using Residuals for the Errors ..... 123\\
Data Matrix Representation of S ..... 125\\
Finite-Sample Considerations ..... 125\\
2.6 Implications of Conditional Homoskedasticity ..... 126\\
Conditional versus Unconditional Homoskedasticity ..... 126\\
Reduction to Finite-Sample Formulas ..... 127\\
Large-Sample Distribution of $t$ and $F$ Statistics ..... 128\\
Variations of Asymptotic Tests under Conditional Homoskedasticity ..... 129\\
2.7 Testing Conditional Homoskedasticity ..... 131\\
2.8 Estimation with Parameterized Conditional Heteroskedasticity (optional) ..... 133\\
The Functional Form ..... 133\\
WLS with Known $\boldsymbol{\alpha}$ ..... 134\\
Regression of $e_{i}^{2}$ on $\mathbf{z}_{i}$ Provides a Consistent Estimate of $\boldsymbol{\alpha}$ ..... 135\\
WLS with Estimated $\boldsymbol{\alpha}$ ..... 136\\
OLS versus WLS ..... 137\\
2.9 Least Squares Projection ..... 137\\
Optimally Predicting the Value of the Dependent Variable ..... 138\\
Best Linear Predictor ..... 139\\
OLS Consistently Estimates the Projection Coefficients ..... 140\\
2.10 Testing for Serial Correlation ..... 141\\
Box-Pierce and Ljung-Box ..... 142\\
Sample Autocorrelations Calculated from Residuals ..... 144\\
Testing with Predetermined, but Not Strictly Exogenous, Regressors ..... 146\\
An Auxiliary Regression-Based Test ..... 147\\
2.11 Application: Rational Expectations Econometrics ..... 150\\
The Efficient Market Hypotheses ..... 150\\
Testable Implications ..... 152\\
Testing for Serial Correlation ..... 153\\
Is the Nominal Interest Rate the Optimal Predictor? ..... 156\\
$R_{t}$ Is Not Strictly Exogenous ..... 158\\
Subsequent Developments ..... 159\\
2.12 Time Regressions ..... 160\\
The Asymptotic Distribution of the OLS Estimator ..... 161\\
Hypothesis Testing for Time Regressions ..... 163\\
Appendix 2.A: Asymptotics with Fixed Regressors ..... 164\\
Appendix 2.B: Proof of Proposition 2.10 ..... 165\\
Problem Set ..... 168\\
Answers to Selected Questions ..... 183\\
3 Single-Equation GMM ..... 186\\
3.1 Endogeneity Bias: Working's Example ..... 187\\
A Simultaneous Equations Model of Market Equilibrium ..... 187\\
Endogeneity Bias ..... 188\\
Observable Supply Shifters ..... 189\\
3.2 More Examples ..... 193\\
A Simple Macroeconometric Model ..... 193\\
Errors-in-Variables ..... 194\\
Production Function ..... 196\\
3.3 The General Formulation ..... 198\\
Regressors and Instruments ..... 198\\
Identification ..... 200\\
Order Condition for Identification ..... 202\\
The Assumption for Asymptotic Normality ..... 202\\
3.4 Generalized Method of Moments Defined ..... 204\\
Method of Moments ..... 205\\
Generalized Method of Moments ..... 206\\
Sampling Error ..... 207\\
3.5 Large-Sample Properties of GMM ..... 208\\
Asymptotic Distribution of the GMM Estimator ..... 209\\
Estimation of Error Variance ..... 210\\
Hypothesis Testing ..... 211\\
Estimation of S ..... 212\\
Efficient GMM Estimator ..... 212\\
Asymptotic Power ..... 214\\
Small-Sample Properties ..... 215\\
3.6 Testing Overidentifying Restrictions ..... 217\\
Testing Subsets of Orthogonality Conditions ..... 218\\
3.7 Hypothesis Testing by the Likelihood-Ratio Principle ..... 222\\
The $L R$ Statistic for the Regression Model ..... 223\\
Variable Addition Test (optional) ..... 224\\
3.8 Implications of Conditional Homoskedasticity ..... 225\\
Efficient GMM Becomes 2SLS ..... 226\\
$J$ Becomes Sargan's Statistic ..... 227\\
Small-Sample Properties of 2SLS ..... 229\\
Alternative Derivations of 2SLS ..... 229\\
When Regressors Are Predetermined ..... 231\\
Testing a Subset of Orthogonality Conditions ..... 232\\
Testing Conditional Homoskedasticity ..... 234\\
Testing for Serial Correlation ..... 234\\
3.9 Application: Returns from Schooling ..... 236\\
The NLS-Y Data ..... 236\\
The Semi-Log Wage Equation ..... 237\\
Omitted Variable Bias ..... 238\\
IQ as the Measure of Ability ..... 239\\
Errors-in-Variables ..... 239\\
2SLS to Correct for the Bias ..... 242\\
Subsequent Developments ..... 243\\
Problem Set ..... 244\\
Answers to Selected Questions ..... 254\\
4 Multiple-Equation GMM ..... 258\\
4.1 The Multiple-Equation Model ..... 259\\
Linearity ..... 259\\
Stationarity and Ergodicity ..... 260\\
Orthogonality Conditions ..... 261\\
Identification ..... 262\\
The Assumption for Asymptotic Normality ..... 264\\
Connection to the "Complete" System of Simultaneous Equations ..... 265\\
4.2 Multiple-Equation GMM Defined ..... 265\\
4.3 Large-Sample Theory ..... 268\\
4.4 Single-Equation versus Multiple-Equation Estimation ..... 271\\
When Are They "Equivalent"? ..... 272\\
Joint Estimation Can Be Hazardous ..... 273\\
4.5 Special Cases of Multiple-Equation GMM: FIVE, 3SLS, and SUR ..... 274\\
Conditional Homoskedasticity ..... 274\\
Full-Information Instrumental Variables Efficient (FIVE) ..... 275\\
Three-Stage Least Squares (3SLS) ..... 276\\
Seemingly Unrelated Regressions (SUR) ..... 279\\
SUR versus OLS ..... 281\\
4.6 Common Coefficients ..... 286\\
The Model with Common Coefficients ..... 286\\
The GMM Estimator ..... 287\\
Imposing Conditional Homoskedasticity ..... 288\\
Pooled OLS ..... 290\\
Beautifying the Formulas ..... 292\\
The Restriction That Isn't ..... 293\\
4.7 Application: Interrelated Factor Demands ..... 296\\
The Translog Cost Function ..... 296\\
Factor Shares ..... 297\\
Substitution Elasticities ..... 298\\
Properties of Cost Functions ..... 299\\
Stochastic Specifications ..... 300\\
The Nature of Restrictions ..... 301\\
Multivariate Regression Subject to Cross-Equation Restrictions ..... 302\\
Which Equation to Delete? ..... 304\\
Results ..... 305\\
Problem Set ..... 308\\
Answers to Selected Questions ..... 320\\
5 Panel Data ..... 323\\
5.1 The Error-Components Model ..... 324\\
Error Components ..... 324\\
Group Means ..... 327\\
A Reparameterization ..... 327\\
5.2 The Fixed-Effects Estimator ..... 330\\
The Formula ..... 330\\
Large-Sample Properties ..... 331\\
Digression: When $\boldsymbol{\eta}_{i}$ Is Spherical ..... 333\\
Random Effects versus Fixed Effects ..... 334\\
Relaxing Conditional Homoskedasticity ..... 335\\
5.3 Unbalanced Panels (optional) ..... 337\\
"Zeroing Out" Missing Observations ..... 338\\
Zeroing Out versus Compression ..... 339\\
No Selectivity Bias ..... 340\\
5.4 Application: International Differences in Growth Rates ..... 342\\
Derivation of the Estimation Equation ..... 342\\
Appending the Error Term ..... 343\\
Treatment of $\alpha_{i}$ ..... 344\\
Consistent Estimation of Speed of Convergence ..... 345\\
Appendix 5.A: Distribution of Hausman Statistic ..... 346\\
Problem Set ..... 349\\
Answers to Selected Questions ..... 363\\
6 Serial Correlation ..... 365\\
6.1 Modeling Serial Correlation: Linear Processes ..... 365\\
MA( $q$ ) ..... 366\\
MA( $\infty$ ) as a Mean Square Limit ..... 366\\
Filters ..... 369\\
Inverting Lag Polynomials ..... 372\\
6.2 ARMA Processes ..... 375\\
AR(1) and Its MA( $\infty$ ) Representation ..... 376\\
Autocovariances of AR(1) ..... 378\\
AR( $p$ ) and Its MA( $\infty$ ) Representation ..... 378\\
$\operatorname{ARMA}(p, q)$ ..... 380\\
ARMA $(p, q)$ with Common Roots ..... 382\\
Invertibility ..... 383\\
Autocovariance-Generating Function and the Spectrum ..... 383\\
6.3 Vector Processes ..... 387\\
6.4 Estimating Autoregressions ..... 392\\
Estimation of AR(1) ..... 392\\
Estimation of AR(p) ..... 393\\
Choice of Lag Length ..... 394\\
Estimation of VARs ..... 397\\
Estimation of ARMA( $p, q$ ) ..... 398\\
6.5 Asymptotics for Sample Means of Serially Correlated Processes ..... 400\\
LLN for Covariance-Stationary Processes ..... 401\\
Two Central Limit Theorems ..... 402\\
Multivariate Extension ..... 404\\
6.6 Incorporating Serial Correlation in GMM ..... 406\\
The Model and Asymptotic Results ..... 406\\
Estimating S When Autocovariances Vanish after Finite Lags ..... 407\\
Using Kernels to Estimate S ..... 408\\
VARHAC ..... 410\\
6.7 Estimation under Conditional Homoskedasticity (Optional) ..... 413\\
Kernel-Based Estimation of $\mathbf{S}$ under Conditional Homoskedasticity ..... 413\\
Data Matrix Representation of Estimated Long-Run Variance ..... 414\\
Relation to GLS ..... 415\\
6.8 Application: Forward Exchange Rates as Optimal Predictors ..... 418\\
The Market Efficiency Hypothesis ..... 419\\
Testing Whether the Unconditional Mean Is Zero ..... 420\\
Regression Tests ..... 423\\
Problem Set ..... 428\\
Answers to Selected Questions ..... 441\\
7 Extremum Estimators ..... 445\\
7.1 Extremum Estimators ..... 446\\
"Measurability" of $\hat{\boldsymbol{\theta}}$ ..... 446\\
Two Classes of Extremum Estimators ..... 447\\
Maximum Likelihood (ML) ..... 448\\
Conditional Maximum Likelihood ..... 450\\
Invariance of ML ..... 452\\
Nonlinear Least Squares (NLS) ..... 453\\
Linear and Nonlinear GMM ..... 454\\
7.2 Consistency ..... 456\\
Two Consistency Theorems for Extremum Estimators ..... 456\\
Consistency of M-Estimators ..... 458\\
Concavity after Reparameterization ..... 461\\
Identification in NLS and ML ..... 462\\
Consistency of GMM ..... 467\\
7.3 Asymptotic Normality ..... 469\\
Asymptotic Normality of M-Estimators ..... 470\\
Consistent Asymptotic Variance Estimation ..... 473\\
Asymptotic Normality of Conditional ML ..... 474\\
Two Examples ..... 476\\
Asymptotic Normality of GMM ..... 478\\
GMM versus ML ..... 481\\
Expressing the Sampling Error in a Common Format ..... 483\\
7.4 Hypothesis Testing ..... 487\\
The Null Hypothesis ..... 487\\
The Working Assumptions ..... 489\\
The Wald Statistic ..... 489\\
The Lagrange Multiplier (LM) Statistic ..... 491\\
The Likelihood Ratio (LR) Statistic ..... 493\\
Summary of the Trinity ..... 494\\
7.5 Numerical Optimization ..... 497\\
Newton-Raphson ..... 497\\
Gauss-Newton ..... 498\\
Writing Newton-Raphson and Gauss-Newton in a Common Format ..... 498\\
Equations Nonlinear in Parameters Only ..... 499\\
Problem Set ..... 501\\
Answers to Selected Questions ..... 505\\
8 Examples of Maximum Likelihood ..... 507\\
8.1 Qualitative Response (QR) Models ..... 507\\
Score and Hessian for Observation $t$ ..... 508\\
Consistency ..... 509\\
Asymptotic Normality ..... 510\\
8.2 Truncated Regression Models ..... 511\\
The Model ..... 511\\
Truncated Distributions ..... 512\\
The Likelihood Function ..... 513\\
Reparameterizing the Likelihood Function ..... 514\\
Verifying Consistency and Asymptotic Normality ..... 515\\
Recovering Original Parameters ..... 517\\
8.3 Censored Regression (Tobit) Models ..... 518\\
Tobit Likelihood Function ..... 518\\
Reparameterization ..... 519\\
8.4 Multivariate Regressions ..... 521\\
The Multivariate Regression Model Restated ..... 522\\
The Likelihood Function ..... 523\\
Maximizing the Likelihood Function ..... 524\\
Consistency and Asymptotic Normality ..... 525\\
8.5 FIML ..... 526\\
The Multiple-Equation Model with Common Instruments Restated ..... 526\\
The Complete System of Simultaneous Equations ..... 529\\
Relationship between ( $\boldsymbol{\Gamma}_{0}, \mathbf{B}_{0}$ ) and $\boldsymbol{\delta}_{0}$ ..... 530\\
The FIML Likelihood Function ..... 531\\
The FIML Concentrated Likelihood Function ..... 532\\
Testing Overidentifying Restrictions ..... 533\\
Properties of the FIML Estimator ..... 533\\
ML Estimation of the SUR Model ..... 535\\
8.6 LIML ..... 538\\
LIML Defined ..... 538\\
Computation of LIML ..... 540\\
LIML versus 2SLS ..... 542\\
8.7 Serially Correlated Observations ..... 543\\
Two Questions ..... 543\\
Unconditional ML for Dependent Observations ..... 545\\
ML Estimation of AR(1) Processes ..... 546\\
Conditional ML Estimation of AR(1) Processes ..... 547\\
Conditional ML Estimation of AR( $p$ ) and VAR( $p$ ) Processes ..... 549\\
Problem Set ..... 551\\
9 Unit-Root Econometrics ..... 557\\
9.1 Modeling Trends ..... 557\\
Integrated Processes ..... 558\\
Why Is It Important to Know if the Process Is I(1)? ..... 560\\
Which Should Be Taken as the Null, I(0) or I(1)? ..... 562\\
Other Approaches to Modeling Trends ..... 563\\
9.2 Tools for Unit-Root Econometrics ..... 563\\
Linear I(0) Processes ..... 563\\
Approximating I(1) by a Random Walk ..... 564\\
Relation to ARMA Models ..... 566\\
The Wiener Process ..... 567\\
A Useful Lemma ..... 570\\
9.3 Dickey-Fuller Tests ..... 573\\
The AR(1) Model ..... 573\\
Deriving the Limiting Distribution under the I(1) Null ..... 574\\
Incorporating the Intercept ..... 577\\
Incorporating Time Trend ..... 581\\
9.4 Augmented Dickey-Fuller Tests ..... 585\\
The Augmented Autoregression ..... 585\\
Limiting Distribution of the OLS Estimator ..... 586\\
Deriving Test Statistics ..... 590\\
Testing Hypotheses about $\zeta$ ..... 591\\
What to Do When $p$ Is Unknown? ..... 592\\
A Suggestion for the Choice of $p_{\text {max }}(T)$ ..... 594\\
Including the Intercept in the Regression ..... 595\\
Incorporating Time Trend ..... 597\\
Summary of the DF and ADF Tests and Other Unit-Root Tests ..... 599\\
9.5 Which Unit-Root Test to Use? ..... 601\\
Local-to-Unity Asymptotics ..... 602\\
Small-Sample Properties ..... 602\\
9.6 Application: Purchasing Power Parity ..... 603\\
The Embarrassing Resiliency of the Random Walk Model? ..... 604\\
Problem Set ..... 605\\
Answers to Selected Questions ..... 619\\
10 Cointegration ..... 623\\
10.1 Cointegrated Systems ..... 624\\
Linear Vector $\mathrm{I}(0)$ and $\mathrm{I}(1)$ Processes ..... 624\\
The Beveridge-Nelson Decomposition ..... 627\\
Cointegration Defined ..... 629\\
10.2 Alternative Representations of Cointegrated Systems ..... 633\\
Phillips's Triangular Representation ..... 633\\
VAR and Cointegration ..... 636\\
The Vector Error-Correction Model (VECM) ..... 638\\
Johansen's ML Procedure ..... 640\\
10.3 Testing the Null of No Cointegration ..... 643\\
Spurious Regressions ..... 643\\
The Residual-Based Test for Cointegration ..... 644\\
Testing the Null of Cointegration ..... 649\\
10.4 Inference on Cointegrating Vectors ..... 650\\
The SOLS Estimator ..... 650\\
The Bivariate Example ..... 652\\
Continuing with the Bivariate Example ..... 653\\
Allowing for Serial Correlation ..... 654\\
General Case ..... 657\\
Other Estimators and Finite-Sample Properties ..... 658\\
10.5 Application: The Demand for Money in the United States ..... 659\\
The Data ..... 660\\
$(m-p, y, R)$ as a Cointegrated System ..... 660\\
DOLS ..... 662\\
Unstable Money Demand? ..... 663\\
Problem Set ..... 665\\
Appendix A: Partitioned Matrices and Kronecker Products ..... 670\\
Addition and Multiplication of Partitioned Matrices ..... 671\\
Inverting Partitioned Matrices ..... 672\\
Index ..... 675

\section*{List of Figures}
1.1 Hypothetical, True, and Estimated Values ..... 16\\
1.2 Equipment Investment and Growth ..... 23

\begin{enumerate}
  \item 3 t Distribution ..... 38\\
1.4 F Distribution ..... 43\\
$1.5 t$-versus $F$-Tests ..... 45\\
1.6 Output Determination ..... 62\\
1.7 Plot of Residuals against Log Output ..... 69\\
2.1 Plots of Serially Uncorrelated Time Series: (a) An i.i.d. $N(0,1)$ sequence. (b) $\mathrm{ARCH}(1)$ with shocks taken from panel (a) ..... 107\\
2.2 What's Observed When ..... 151\\
2.3 Inflation and Interest Rates ..... 154\\
2.4 Real Interest Rates ..... 154\\
3.1 Supply and Demand Curves ..... 190\\
4.1 OLS and GMM ..... 283\\
6.1 Forecast Error, Yen/Dollar ..... 421\\
6.2 Comelogram of $30-f$, Yen/Dollar ..... 422\\
6.3 Yen/Dollar Spot Rate, Jan. 1975-Dec. 1989 ..... 424\\
6.4 Plot of s30 against $f$, Yen/Dollar ..... 425\\
6.5 Plot of $s 30-s$ against $f-s$, Yen/Dollar ..... 427\\
6.6 Correlogram of $s_{t+1}-s_{t}$, Yen/Dollar ..... 443\\
6.7 Comelogram of Three-Month Ex-Post Real Rate ..... 443\\
8.1 Effect of Truncation ..... 512\\
9.1 Log U.S. Real GDP, Value for 1869 Set to 0 ..... 558\\
9.2 Illustration of the Functional Central Limit Theorem ..... 569\\
9.3 Finite-Sample Distribution of OLS Estimate of AR(1) Coefficient, $T=100$ ..... 577\\
9.4 Dollar/Sterling Real Exchange Rate, 1791-1990, 1914 Value Set to 0 ..... 605\\
10.1 (a) Log Income and Consumption ..... 624\\
10.1 (b) Log Income Minus Log Consumption ..... 649\\
10.2 U.S. Real NNP and Real M1, in Logs ..... 661\\
10.3 Log M1 Velocity (left scale) and Short-Term Commercial Paper Rate (right scale) ..... 661\\
10.4 Plot of $m-p-y$ against $R$ ..... 664
\end{enumerate}

\section*{Preface}
This book is designed to serve as the textbook for a first-year graduate course in econometrics. It has two distinguishing features. First, it covers a full range of techniques with the estimation method called the Generalized Method of Moments (GMM) as the organizing principle. I believe this unified approach is the most efficient way to cover the first-year materials in an accessible yet rigorous manner. Second, most chapters include a section examining in detail original applied articles from such diverse fields in economics as industrial organization, labor, finance, international, and macroeconomics. So the reader will know how to use the techniques covered in the chapter and under what conditions they are applicable.

Over the last several years, the lecture notes on which this book is based have been used at the University of Pennsylvania, Columbia University, Princeton University, the University of Tokyo, Boston College, Harvard University, and Ohio State University. Students seem to like the book a lot. My own experience from teaching out of the book is that students think the book is better than the instructor.

\section*{Prerequisites}
The reader of this book is assumed to have a working knowledge of the basics of calculus, probability theory, and linear algebra. An understanding of the following concepts is taken for granted: functions of several variables, partial derivatives, integrals, random variables, joint distributions, independence, unconditional and conditional expectations, variances and covariances of vector random variables, normal distributions, chi-square distributions, matrix multiplication, inverses of matrices, the rank of a matrix, determinants, and positive definite matrices. Any relevant concepts above this level will be introduced as the discussion progresses. Results on partitioned matrices and Kronecker products are collected in the appendix. Prior exposure to undergraduate econometrics is not required.

\section*{Organization of the Book}
To understand how the book is organized, it is useful to distinguish between a model and an estimation procedure. The basic premise of econometrics is that economic data (such as postwar U.S. GDP) are random variables. A model is a\\
family of probability distributions that could possibly have generated the economic data. An estimation procedure is a data-based protocol for choosing from the model a particular distribution that is likely to have generated the data. Most estimation procedures in econometrics are a specialization of the GMM estimation principle. For example, when GMM is applied to a model called the classical linear regression model, the resulting estimation procedure is Ordinary Least Squares (OLS), the most basic estimation procedure in econometrics. This viewpoint is the organizing principle in the first six chapters of the book, where most of the standard estimation procedures are presented.

The book could have presented GMM in the first chapter, but that would deprive the reader of the option to follow a series of topics specific to OLS without getting distracted by GMM. For this reason I chose to use the first two chapters to present the finite-sample and large-sample theory of OLS. GMM is presented in Chapter 3 as a generalization of OLS.

A major expositional innovation of the book is to treat multiple-equation estimation procedures - such as Seemingly Unrelated Regressions (SUR), Three-Stage Least Squares (3SLS), the Random-Effects method, covered in Chapter 4, and the Fixed-Effects method covered in Chapter 5-as special cases of the singleequation GMM of Chapter 3. This makes it possible to derive the statistical properties of those advanced techniques merely by suitably specializing the results about single-equation GMM developed in Chapter 3. Chapter 6 completes the book's discussion of GMM by indicating how serial dependence in the error term can be incorporated in GMM.

For some models in econometrics, Maximum Likelihood (ML) is the more natural estimation principle than GMM. ML is covered in Chapters 7 and 8. To make clear the relationship between GMM and ML, the book's discussion of ML starts out in Chapter 7 with an estimation principle called Extremum Estimators, which includes both ML and GMM as special cases. Applications of ML to various models are covered in Chapter 8.

The book also includes an extensive treatment of time-series analysis. Basic time-series topics are covered in Section 2.2 and in the first half of Chapter 6. That is enough of a prerequisite for the important recent advances in nonstationary time-series analysis, which are covered in the last two chapters of the book.

\section*{Designing a Course Out of the Book}
Several different courses can be designed based on the book.

\begin{itemize}
  \item Assuming that the course meets twice for an hour and a half, eight weeks should be enough to cover core theory, which is Chapters 1-4 and 6 (excluding those\\
sections on specific economic applications), Chapter 7 (with proofs and examples skipped), and Chapter 8.
  \item A twelve-week semester course can cover, in addition to the core theory, Chapter 5 and the economic applications included in Chapters 1-6.
  \item A short (say, six-week) course specializing in GMM estimation in cross-section and panel data would cover Chapters 1-5 (excluding Sections 2.10-2.12 but including Section 2.2). Chapters 7 and 8 (excluding Section 8.7) would add the ML component to the course.
  \item A short time-series course covering recent developments with economic applications would consist of Chapter 1 (excluding Section 1.7), Chapter 2, parts of Chapters 6 and 7 (Sections 6.1-6.5, 6.8, and 7.1), and Chapters 9 and 10. The applications sections in Chapters 2, 6,9, and 10 can be dropped if the course focuses on just theory.
\end{itemize}

\section*{Review Questions and Analytical Exercises}
The book includes a number of short questions for review at the end of each chapter, with lots of hints (and even answers). They can be used to check whether the reader actually understood the materials of the section. On the second reading, if not the first, readers should try to answer them before proceeding. Answers to selected review questions are available from the book's website, \href{http://pup.princeton.edu/}{http://pup.princeton.edu/} titles/6946.html.

There are several analytical exercises at the end of each chapter that ask the reader to prove results left unproved in the text or supplementary results that are useful for their own sake. Unless otherwise indicated, analytical exercises can be skipped without loss of continuity.

\section*{Empirical Exercises}
Each chapter usually has one big empirical exercise. It asks for a replication of the empirical results of the original article discussed in the applications section of the chapter and an estimation of the various extensions of the article's model using the estimation procedures presented in the chapter. The dataset for estimation, which usually is the same as the one used by the original article, can be downloaded from the book's website mentioned above.

To implement the estimation procedure on the dataset, readers should run a statistical package on a computer. There are many statistical packages that are widely used in econometrics. They include GAUSS (\href{http://www.aptech.com}{www.aptech.com}), MATLAB (\href{http://www.mathworks.com}{www.mathworks.com}), Eviews (\href{http://www.eviews.com}{www.eviews.com}), LIMDEP (\href{http://www.limdep.com}{www.limdep.com}),

RATS (\href{http://www.estima.com}{www.estima.com}), SAS (\href{http://www.sas.com}{www.sas.com}), Stata (\href{http://www.stata.com}{www.stata.com}), and TSP (\href{http://www.tsp.com}{www.tsp.com}). GAUSS and MATLAB are different from the rest in that they are matrix-based languages rather than a collection of procedures. Consider, for example, carrying out OLS with GAUSS or MATLAB. After loading the dataset into the computer's workspace, it takes several lines reflecting the matrix operations of OLS to calculate the OLS estimate and associated statistics (such as $R^{2}$ ). With the other packages, which are procedure-driven and sometimes called "canned packages," those several lines can be replaced by a one-line command invoking the OLS procedure. For example, TSP's OLS command is OLSQ.

There are advantages and disadvantages with canned packages. Obviously, it takes far fewer lines to accomplish the same thing, so one can spend less time on programming. On the other hand, procedures in a canned package, which accept data and spit out point estimates and associated statistics, are essentially a black box. Sometimes it is not clear from the documentation of the package how certain statistics are calculated. Although those canned packages mentioned above regularly incorporate new developments in econometrics, the estimation procedure desired may not be currently supported by the package, in which case it will be necessary to write one's own procedures in GAUSS or MATLAB. But it may be a blessing in disguise; actually writing down the underlying matrix operations provides you with an excellent chance to understand the estimation procedure.

With only a few exceptions, all the calculations needed for the empirical exercises of the book can be carried out with any of the canned packages mentioned above. My recommendation, therefore, is for economics Ph.D. students who are planning to write an applied thesis using state-of-the-art estimation procedures or a theoretical thesis proposing new ones, to use GAUSS or MATLAB. Otherwise, students should use any of the canned packages mentioned above.

\section*{Mathematical Notation}
There is no single mathematical notation used by everyone in econometrics. The book's notation follows the most standard, if not universal, practice. Vectors are treated as column vectors and written in bold lowercase letters. Matrices are in bold uppercase letters. The transpose of the matrix $\mathbf{A}$ is denoted by $\mathbf{A}^{\prime}$. Scalar variables are (mostly lowercase) letters in italics.

\section*{Acknowledgments}
I acknowledge with gratitude help from the following individuals and institutions. Mark Watson, Dale Jorgenson, Bo Honore, Serena Ng, Masao Ogaki, and Jushan Bai were kind and courageous enough to use early versions of the book as the\\
textbook for their econometrics courses. Comments made by them and their students have been incorporated in this final version. Yuzo Honda read the manuscript and offered helpful suggestions. Naoto Kunitomo, Whitney Newey, Serena Ng, Pierre Perron, Jim Stock, Katsuto Tanaka, Mark Watson, Hal White, and Yoshihiro Yajima took time out to answer my questions and enquiries. Two graduate students at University of Tokyo, Mari Sakudo and Naoki Shimoi, read the entire manuscript to weed out typos. Their effort was underwritten by a grant-in-aid from the Zengin Foundation for Studies on Economics and Finance. Peter Dougherty, my editor at Princeton University Press, provided enthusiasm and just the right amount of pressure. Stephanie Hogue was a versatile enough $\mathrm{IAT}_{\mathrm{E}} \mathrm{X}$ expert to accommodate my formatting whims. Ellen Foos supervised production of the book. Last but not least, Jessica Helfand agreed, probably out of friendship, to do the cover design.

For more than five years all my free time went into writing this book. Now, having completed the book, I feel like someone who has just been released from prison. My research suffered, but hopefully the profession has not noticed.

\begin{abstract}
The Ordinary Least Squares (OLS) estimator is the most basic estimation procedure in econometrics. This chapter covers the finite- or small-sample properties of the OLS estimator, that is, the statistical properties of the OLS estimator that are valid for any given sample size. The materials covered in this chapter are entirely standard. The exposition here differs from that of most other textbooks in its emphasis on the role played by the assumption that the regressors are "strictly exogenous."\\
In the final section, we apply the finite-sample theory to the estimation of the cost function using cross-section data on individual firms. The question posed in Nerlove's (1963) study is of great practical importance: are there increasing returns to scale in electricity supply? If yes, microeconomics tells us that the industry should be regulated. Besides providing you with a hands-on experience of using the techniques to test interesting hypotheses, Nerlove's paper has a careful discussion of why the OLS is an appropriate estimation procedure in this particular application.
\end{abstract}

\subsection*{1.1 The Classical Linear Regression Model}
In this section we present the assumptions that comprise the classical linear regression model. In the model, the variable in question (called the dependent variable, the regressand, or more generically the left-hand [-side] variable) is related to several other variables (called the regressors, the explanatory variables, or the right-hand [-side] variables). Suppose we observe $n$ values for those variables. Let $y_{i}$ be the $i$-th observation of the dependent variable in question and let ( $x_{i 1}, x_{i 2}, \ldots, x_{i K}$ ) be the $i$-th observation of the $K$ regressors. The sample or data is a collection of those $n$ observations.

The data in economics cannot be generated by experiments (except in experimental economics), so both the dependent and independent variables have to be treated as random variables, variables whose values are subject to chance. A model\\
is a set of restrictions on the joint distribution of the dependent and independent variables. That is, a model is a set of joint distributions satisfying a set of assumptions. The classical regression model is a set of joint distributions satisfying Assumptions 1.1-1.4 stated below.

\section*{The Linearity Assumption}
The first assumption is that the relationship between the dependent variable and the regressors is linear.

\section*{Assumption 1.1 (linearity):}

\begin{equation*}
y_{i}=\beta_{1} x_{i 1}+\beta_{2} x_{i 2}+\cdots+\beta_{K} x_{i K}+\varepsilon_{i} \quad(i=1,2, \ldots, n), \tag{1.1.1}
\end{equation*}


where $\beta$ 's are unknown parameters to be estimated, and $\varepsilon_{i}$ is the unobserved error term with certain properties to be specified below.

The part of the right-hand side involving the regressors, $\beta_{1} x_{i 1}+\beta_{2} x_{i 2}+\cdots+\beta_{K} x_{i K}$, is called the regression or the regression function, and the coefficients ( $\beta$ 's) are called the regression coefficients. They represent the marginal and separate effects of the regressors. For example, $\beta_{2}$ represents the change in the dependent variable when the second regressor increases by one unit while other regressors are held constant. In the language of calculus, this can be expressed as $\partial y_{i} / \partial x_{i 2}=\beta_{2}$. The linearity implies that the marginal effect does not depend on the level of regressors. The error term represents the part of the dependent variable left unexplained by the regressors.

Example 1.1 (consumption function): The simple consumption function familiar from introductory economics is


\begin{equation*}
C O N_{i}=\beta_{1}+\beta_{2} Y D_{i}+\varepsilon_{i} \tag{1.1.2}
\end{equation*}


where $C O N$ is consumption and $Y D$ is disposable income. If the data are annual aggregate time-series, $C O N_{i}$ and $Y D_{i}$ are aggregate consumption and disposable income for year $i$. If the data come from a survey of individual households, $C O N_{i}$ is consumption by the $i$-th household in the cross-section sample of $n$ households. The consumption function can be written as (1.1.1) by setting $y_{i}=C O N_{i}, x_{i 1}=1$ (a constant), and $x_{i 2}=Y D_{i}$. The error term $\varepsilon_{i}$ represents other variables besides disposable income that influence consumption. They include those variables - such as financial assets - that\\
might be observable but the researcher decided not to include as regressors, as well as those variables - such as the "mood" of the consumer - that are hard to measure. When the equation has only one nonconstant regressor, as here, it is called the simple regression model.

The linearity assumption is not as restrictive as it might first seem, because the dependent variable and the regressors can be transformations of the variables in question. Consider

Example 1.2 (wage equation): A simplified version of the wage equation routinely estimated in labor economics is


\begin{equation*}
\log \left(W A G E_{i}\right)=\beta_{1}+\beta_{2} S_{i}+\beta_{3} T E N U R E_{i}+\beta_{4} E X P R_{i}+\varepsilon_{i}, \tag{1.1.3}
\end{equation*}


where $W A G E=$ the wage rate for the individual, $S=$ education in years, TENURE = years on the current job, and $E X P R=$ experience in the labor force (i.e., total number of years to date on all the jobs held currently or previously by the individual). The wage equation fits the generic format (1.1.1) with $y_{i}=\log \left(W A G E_{i}\right)$. The equation is said to be in the semi-log form because only the dependent variable is in logs. The equation is derived from the following nonlinear relationship between the level of the wage rate and the regressors:


\begin{equation*}
W A G E_{i}=\exp \left(\beta_{1}\right) \exp \left(\beta_{2} S_{i}\right) \exp \left(\beta_{3} T E N U R E_{i}\right) \exp \left(\beta_{4} E X P R_{i}\right) \exp \left(\varepsilon_{i}\right) . \tag{1.1.4}
\end{equation*}


By taking logs of both sides of (1.1.4) and noting that $\log [\exp (x)]=x$, one obtains (1.1.3). The coefficients in the semi-log form have the interpretation of percentage changes, not changes in levels. For example, a value of 0.05 for $\beta_{2}$ implies that an additional year of education has the effect of raising the wage rate by 5 percent. The difference in the interpretation comes about because the dependent variable is the log wage rate, not the wage rate itself, and the change in logs equals the percentage change in levels.

Certain other forms of nonlinearities can also be accommodated. Suppose, for example, the marginal effect of education tapers off as the level of education gets higher. This can be captured by including in the wage equation the squared term $S^{2}$ as an additional regressor in the wage equation. If the coefficient of the squared\\
term is $\beta_{5}$, the marginal effect of education is

$$
\beta_{2}+2 \beta_{5} S(=\partial \log (W A G E) / \partial S)
$$

If $\beta_{5}$ is negative, the marginal effect of education declines with the level of education.

There are, of course, cases of genuine nonlinearity. For example, the relationship (1.1.4) could not have been made linear if the error term entered additively rather than multiplicatively:

$$
W A G E_{i}=\exp \left(\beta_{1}\right) \exp \left(\beta_{2} S_{i}\right) \exp \left(\beta_{3} T E N U R E_{i}\right) \exp \left(\beta_{4} E X P R_{i}\right)+\varepsilon_{i}
$$

Estimation of nonlinear regression equations such as this will be discussed in Chapter 7.

\section*{Matrix Notation}
Before stating other assumptions of the classical model, we introduce the vector and matrix notation. The notation will prove useful for stating other assumptions precisely and also for deriving the OLS estimator of $\boldsymbol{\beta}$. Define $K$-dimensional (column) vectors $\mathbf{x}_{i}$ and $\boldsymbol{\beta}$ as

\[
\underset{(K \times 1)}{\mathbf{x}_{i}}=\left[\begin{array}{c}
x_{i 1}  \tag{1.1.5}\\
x_{i 2} \\
\vdots \\
x_{i K}
\end{array}\right], \underset{(K \times 1)}{\boldsymbol{\beta}}=\left[\begin{array}{c}
\beta_{1} \\
\beta_{2} \\
\vdots \\
\beta_{K}
\end{array}\right] .
\]

By the definition of vector inner products, $\mathbf{x}_{i}^{\prime} \boldsymbol{\beta}=\beta_{1} x_{i 1}+\beta_{2} x_{i 2}+\cdots+\beta_{K} x_{i K}$. So the equations in Assumption 1.1 can be written as


\begin{equation*}
y_{i}=\mathbf{x}_{i}^{\prime} \boldsymbol{\beta}+\varepsilon_{i} \quad(i=1,2, \ldots, n) . \tag{1.1.1́}
\end{equation*}


Also define

\[
\underset{(n \times 1)}{\mathbf{y}}=\left[\begin{array}{c}
y_{1}  \tag{1.1.6}\\
\vdots \\
y_{n}
\end{array}\right], \underset{(n \times 1)}{\boldsymbol{\varepsilon}}=\left[\begin{array}{c}
\varepsilon_{1} \\
\vdots \\
\varepsilon_{n}
\end{array}\right], \underset{(n \times K)}{\mathbf{X}}=\left[\begin{array}{c}
\mathbf{x}_{1}^{\prime} \\
\vdots \\
\mathbf{x}_{n}^{\prime}
\end{array}\right]=\left[\begin{array}{ccc}
x_{11} & \cdots & x_{1 K} \\
\vdots & \cdots & \vdots \\
x_{n 1} & \cdots & x_{n K}
\end{array}\right] .
\]

In the vectors and matrices in (1.1.6), there are as many rows as there are observations, with the rows corresponding to the observations. For this reason $\mathbf{y}$ and $\mathbf{X}$ are sometimes called the data vector and the data matrix. Since the number of\\
columns of $\mathbf{X}$ equals the number of rows of $\boldsymbol{\beta}, \mathbf{X}$ and $\boldsymbol{\beta}$ are conformable and $\mathbf{X} \boldsymbol{\beta}$ is an $n \times 1$ vector. Its $i$-th element is $\mathbf{x}_{i}^{\prime} \boldsymbol{\beta}$. Therefore, Assumption 1.1 can be written compactly as

$$
\underset{(n \times 1)}{\mathbf{y}}=\underbrace{\underset{(n \times K)}{\mathbf{X}} \underset{(K \times 1)}{\boldsymbol{\beta}}}_{(n \times 1)}+\underset{(n \times 1)}{\boldsymbol{\varepsilon}}
$$

\section*{The Strict Exogeneity Assumption}
The next assumption of the classical regression model is

\section*{Assumption 1.2 (strict exogeneity):}

\begin{equation*}
\mathrm{E}\left(\varepsilon_{i} \mid \mathbf{X}\right)=0 \quad(i=1,2, \ldots, n) \tag{1.1.7}
\end{equation*}


Here, the expectation (mean) is conditional on the regressors for all observations. This point may be made more apparent by writing the assumption without using the data matrix as

$$
\mathrm{E}\left(\varepsilon_{i} \mid \mathbf{x}_{1}, \ldots, \mathbf{x}_{n}\right)=0 \quad(i=1,2, \ldots, n)
$$

To state the assumption differently, take, for any given observation $i$, the joint distribution of the $n K+1$ random variables, $f\left(\varepsilon_{i}, \mathbf{x}_{1}, \ldots, \mathbf{x}_{n}\right)$, and consider the conditional distribution, $f\left(\varepsilon_{i} \mid \mathbf{x}_{1}, \ldots, \mathbf{x}_{n}\right)$. The conditional mean $\mathrm{E}\left(\varepsilon_{i} \mid \mathbf{x}_{1}, \ldots, \mathbf{x}_{n}\right)$ is in general a nonlinear function of $\left(\mathbf{x}_{1}, \ldots, \mathbf{x}_{n}\right)$. The strict exogeneity assumption says that this function is a constant of value zero. ${ }^{1}$

Assuming this constant to be zero is not restrictive if the regressors include a constant, because the equation can be rewritten so that the conditional mean of the error term is zero. To see this, suppose that $\mathrm{E}\left(\varepsilon_{i} \mid \mathbf{X}\right)$ is $\mu$ and $x_{i 1}=1$. The equation can be written as

$$
\begin{aligned}
y_{i} & =\beta_{1}+\beta_{2} x_{i 2}+\cdots+\beta_{K} x_{i K}+\varepsilon_{i} \\
& =\left(\beta_{1}+\mu\right)+\beta_{2} x_{i 2}+\cdots+\beta_{K} x_{i K}+\left(\varepsilon_{i}-\mu\right)
\end{aligned}
$$

If we redefine $\beta_{1}$ to be $\beta_{1}+\mu$ and $\varepsilon_{i}$ to be $\varepsilon_{i}-\mu$, the conditional mean of the new error term is zero. In virtually all applications, the regressors include a constant term.

\footnotetext{${ }^{1}$ Some authors define the term "strict exogeneity" somewhat differently. For example, in Koopmans and Hood (1953) and Engle, Hendry, and Richards (1983), the regressors are strictly exogenous if $\mathbf{x}_{i}$ is independent of $\varepsilon_{j}$ for all $i, j$. This definition is stronger than, but not inconsistent with, our definition of strict exogeneity.
}Example 1.3 (continuation of Example 1.1): For the simple regression model of Example 1.1, the strict exogeneity assumption can be written as

$$
\mathrm{E}\left(\varepsilon_{i} \mid Y D_{1}, Y D_{2}, \ldots, Y D_{n}\right)=0 .
$$

Since $\mathbf{x}_{i}=\left(1, Y D_{i}\right)^{\prime}$, you might wish to write the strict exogeneity assumption as

$$
\mathrm{E}\left(\varepsilon_{i} \mid 1, Y D_{1}, 1, Y D_{2}, \ldots, 1, Y D_{n}\right)=0 .
$$

But since a constant provides no information, the expectation conditional on

$$
\left(1, Y D_{1}, 1, Y D_{2}, \ldots, 1, Y D_{n}\right)
$$

is the same as the expectation conditional on

$$
\left(Y D_{1}, Y D_{2}, \ldots, Y D_{n}\right) .
$$

\section*{Implications of Strict Exogeneity}
The strict exogeneity assumption has several implications.

\begin{itemize}
  \item The unconditional mean of the error term is zero, i.e.,
\end{itemize}


\begin{equation*}
E\left(\varepsilon_{i}\right)=0 \quad(i=1,2, \ldots, n) . \tag{1.1.8}
\end{equation*}


This is because, by the Law of Total Expectations from basic probability theory, ${ }^{2} \mathrm{E}\left[\mathrm{E}\left(\varepsilon_{i} \mid \mathbf{X}\right)\right]=\mathrm{E}\left(\varepsilon_{i}\right)$.

\begin{itemize}
  \item If the cross moment $\mathrm{E}(x y)$ of two random variables $x$ and $y$ is zero, then we say that $x$ is orthogonal to $y$ (or $y$ is orthogonal to $x$ ). Under strict exogeneity, the regressors are orthogonal to the error term for all observations, i.e.,
\end{itemize}

$$
\mathrm{E}\left(x_{j k} \varepsilon_{i}\right)=0 \quad(i, j=1, \ldots, n ; k=1, \ldots, K)
$$

or

\[
\mathrm{E}\left(\mathbf{x}_{j} \cdot \varepsilon_{i}\right)=\left[\begin{array}{c}
\mathrm{E}\left(x_{j 1} \varepsilon_{i}\right)  \tag{1.1.9}\\
\mathrm{E}\left(x_{j 2} \varepsilon_{i}\right) \\
\vdots \\
\mathrm{E}\left(x_{j K} \varepsilon_{i}\right)
\end{array}\right]=\underset{(K \times 1)}{\mathbf{0}} \quad(\text { for all } i, j) .
\]

\footnotetext{${ }^{2}$ The Law of Total Expectations states that $\mathrm{E}[\mathrm{E}(y \mid \mathbf{x})]=\mathrm{E}(y)$.
}The proof is a good illustration of the use of properties of conditional expectations and goes as follows.

Proof. Since $x_{j k}$ is an element of $\mathbf{X}$, strict exogeneity implies


\begin{equation*}
\mathrm{E}\left(\varepsilon_{i} \mid x_{j k}\right)=\mathrm{E}\left[\mathrm{E}\left(\varepsilon_{i} \mid \mathbf{X}\right) \mid x_{j k}\right]=0 \tag{1.1.10}
\end{equation*}


by the Law of Iterated Expectations from probability theory. ${ }^{3}$ It follows from this that

$$
\begin{aligned}
\mathrm{E}\left(x_{j k} \varepsilon_{i}\right) & =\mathrm{E}\left[\mathrm{E}\left(x_{j k} \varepsilon_{i} \mid x_{j k}\right)\right] \quad \text { (by the Law of Total Expectations) } \\
& =\mathrm{E}\left[x_{j k} \mathrm{E}\left(\varepsilon_{i} \mid x_{j k}\right)\right] \quad \text { (by the linearity of conditional expectations }{ }^{4} \text { ) } \\
& =0
\end{aligned}
$$ \(\square\)

The point here is that strict exogeneity requires the regressors be orthogonal not only to the error term from the same observation (i.e., $\mathrm{E}\left(x_{i k} \varepsilon_{i}\right)=0$ for all $k$ ), but also to the error term from the other observations (i.e., $\mathrm{E}\left(x_{j k} \varepsilon_{i}\right)=0$ for all $k$ and for $j \neq i$ ).

\begin{itemize}
  \item Because the mean of the error term is zero, the orthogonality conditions (1.1.9) are equivalent to zero-correlation conditions. This is because
\end{itemize}

$$
\begin{aligned}
\operatorname{Cov}\left(\varepsilon_{i}, x_{j k}\right) & =\mathrm{E}\left(x_{j k} \varepsilon_{i}\right)-\mathrm{E}\left(x_{j k}\right) \mathrm{E}\left(\varepsilon_{i}\right) \quad \text { (by definition of covariance) } \\
& =\mathrm{E}\left(x_{j k} \varepsilon_{i}\right) \quad\left(\text { since } \mathrm{E}\left(\varepsilon_{i}\right)=0,\right. \text { see (1.1.8)) } \\
& =0 \quad(\text { by the orthogonality conditions (1.1.9)). }
\end{aligned}
$$

In particular, for $i=j, \operatorname{Cov}\left(x_{i k}, \varepsilon_{i}\right)=0$. Therefore, strict exogeneity implies the requirement (familiar to those who have studied econometrics before) that the regressors be contemporaneously uncorrelated with the error term.

\section*{Strict Exogeneity in Time-Series Models}
For time-series models where $i$ is time, the implication (1.1.9) of strict exogeneity can be rephrased as: the regressors are orthogonal to the past, current, and future error terms (or equivalently, the error term is orthogonal to the past, current, and future regressors). But for most time-series models, this condition (and a fortiori strict exogeneity) is not satisfied, so the finite-sample theory based on strict exogeneity to be developed in this section is rarely applicable in time-series con-

\footnotetext{${ }^{3}$ The Law of Iterated Expectations states that $\mathrm{E}[\mathrm{E}(y \mid \mathbf{x}, \mathbf{z}) \mid \mathbf{x}]=\mathrm{E}(y \mid \mathbf{x})$.\\
${ }^{4}$ The linearity of conditional expectations states that $\mathrm{E}[f(\mathbf{x}) y \mid \mathbf{x}]=f(\mathbf{x}) \mathrm{E}(y \mid \mathbf{x})$.
}
texts. However, as will be shown in the next chapter, the estimator possesses good large-sample properties without strict exogeneity.

The clearest example of a failure of strict exogeneity is a model where the regressor includes the lagged dependent variable. Consider the simplest such model:


\begin{equation*}
y_{i}=\beta y_{i-1}+\varepsilon_{i} \quad(i=1,2, \ldots, n) . \tag{1.1.11}
\end{equation*}


This is called the first-order autoregressive model (AR(1)). (We will study this model more fully in Chapter 6.) Suppose, consistent with the spirit of the strict exogeneity assumption, that the regressor for observation $i, y_{i-1}$, is orthogonal to the error term for $i$ so $\mathrm{E}\left(y_{i-1} \varepsilon_{i}\right)=0$. Then

$$
\begin{aligned}
\mathrm{E}\left(y_{i} \varepsilon_{i}\right) & =\mathrm{E}\left[\left(\beta y_{i-1}+\varepsilon_{i}\right) \varepsilon_{i}\right] \quad(\text { by }(1.1 .11)) \\
& =\beta \mathrm{E}\left(y_{i-1} \varepsilon_{i}\right)+\mathrm{E}\left(\varepsilon_{i}^{2}\right) \\
& =\mathrm{E}\left(\varepsilon_{i}^{2}\right) \quad\left(\text { since } \mathrm{E}\left(y_{i-1} \varepsilon_{i}\right)=0 \text { by hypothesis }\right) .
\end{aligned}
$$

Therefore, unless the error term is always zero, $\mathrm{E}\left(y_{i} \varepsilon_{i}\right)$ is not zero. But $y_{i}$ is the regressor for observation $i+1$. Thus, the regressor is not orthogonal to the past error term, which is a violation of strict exogeneity.

\section*{Other Assumptions of the Model}
The remaining assumptions comprising the classical regression model are the following.

Assumption 1.3 (no multicollinearity): The rank of the $n \times K$ data matrix, $\mathbf{X}$, is $K$ with probability 1.

\section*{Assumption 1.4 (spherical error variance):}

\begin{equation*}
\text { (homoskedasticity) } \mathrm{E}\left(\varepsilon_{i}^{2} \mid \mathbf{X}\right)=\sigma^{2}>0 \quad(i=1,2, \ldots, n),^{5} \tag{1.1.12}
\end{equation*}


(no correlation between observations)


\begin{equation*}
\mathrm{E}\left(\varepsilon_{i} \varepsilon_{j} \mid \mathbf{X}\right)=0 \quad(i, j=1,2, \ldots, n ; i \neq j) \tag{1.1.13}
\end{equation*}


\footnotetext{${ }^{5}$ When a symbol (which here is $\sigma^{2}$ ) is given to a moment (which here is the second moment $\mathrm{E}\left(\varepsilon_{i}^{2} \mid \mathbf{X}\right)$ ), by implication the moment is assumed to exist and is finite. We will follow this convention for the rest of this book.
}To understand Assumption 1.3, recall from matrix algebra that the rank of a matrix equals the number of linearly independent columns of the matrix. The assumption says that none of the $K$ columns of the data matrix $\mathbf{X}$ can be expressed as a linear combination of the other columns of $\mathbf{X}$. That is, $\mathbf{X}$ is of full column rank. Since the $K$ columns cannot be linearly independent if their dimension is less than $K$, the assumption implies that $n \geq K$, i.e., there must be at least as many observations as there are regressors. The regressors are said to be (perfectly) multicollinear if the assumption is not satisfied. It is easy to see in specific applications when the regressors are multicollinear and what problems arise.

Example 1.4 (continuation of Example 1.2): If no individuals in the sample ever changed jobs, then $\operatorname{TENURE}_{i}=E X P R_{i}$ for all $i$, in violation of the no multicollinearity assumption. There is evidently no way to distinguish the tenure effect on the wage rate from the experience effect. If we substitute this equality into the wage equation to eliminate $T E N U R E_{i}$, the wage equation becomes

$$
\log \left(W A G E_{i}\right)=\beta_{1}+\beta_{2} S_{i}+\left(\beta_{3}+\beta_{4}\right) E X P R_{i}+\varepsilon_{i}
$$

which shows that only the sum $\beta_{3}+\beta_{4}$, but not $\beta_{3}$ and $\beta_{4}$ separately, can be estimated.

The homoskedasticity assumption (1.1.12) says that the conditional second moment, which in general is a nonlinear function of $\mathbf{X}$, is a constant. Thanks to strict exogeneity, this condition can be stated equivalently in more familiar terms. Consider the conditional variance $\operatorname{Var}\left(\varepsilon_{i} \mid \mathbf{X}\right)$. It equals the same constant because

$$
\begin{aligned}
\operatorname{Var}\left(\varepsilon_{i} \mid \mathbf{X}\right) & \equiv \mathrm{E}\left(\varepsilon_{i}^{2} \mid \mathbf{X}\right)-\mathrm{E}\left(\varepsilon_{i} \mid \mathbf{X}\right)^{2} \quad \text { (by definition of conditional variance) } \\
& =\mathrm{E}\left(\varepsilon_{i}^{2} \mid \mathbf{X}\right) \quad \text { (since } \mathrm{E}\left(\varepsilon_{i} \mid \mathbf{X}\right)=0 \text { by strict exogeneity). }
\end{aligned}
$$

Similarly, (1.1.13) is equivalent to the requirement that

$$
\operatorname{Cov}\left(\varepsilon_{i}, \varepsilon_{j} \mid \mathbf{X}\right)=0 \quad(i, j=1,2, \ldots, n ; i \neq j)
$$

That is, in the joint distribution of $\left(\varepsilon_{i}, \varepsilon_{j}\right)$ conditional on $\mathbf{X}$, the covariance is zero. In the context of time-series models, (1.1.13) states that there is no serial correlation in the error term.

Since the $(i, j)$ element of the $n \times n$ matrix $\boldsymbol{\varepsilon} \boldsymbol{\varepsilon}^{\prime}$ is $\varepsilon_{i} \varepsilon_{j}$, Assumption 1.4 can be written compactly as


\begin{equation*}
\mathrm{E}\left(\boldsymbol{\varepsilon} \boldsymbol{\varepsilon}^{\prime} \mid \mathbf{X}\right)=\sigma^{2} \mathbf{I}_{n} . \tag{1.1.14}
\end{equation*}


The discussion of the previous paragraph shows that the assumption can also be written as

$$
\operatorname{Var}(\boldsymbol{\varepsilon} \mid \mathbf{X})=\sigma^{2} \mathbf{I}_{n}
$$

However, (1.1.14) is the preferred expression, because the more convenient measure of variability is second moments (such as $\left.\mathrm{E}\left(\varepsilon_{i}^{2} \mid \mathbf{X}\right)\right)$ rather than variances. This point will become clearer when we deal with the large sample theory in the next chapter. Assumption 1.4 is sometimes called the spherical error variance assumption because the $n \times n$ matrix of second moments (which are also variances and covariances) is proportional to the identity matrix $\mathbf{I}_{n}$. This assumption will be relaxed later in this chapter.

\section*{The Classical Regression Model for Random Samples}
The sample ( $\mathbf{y}, \mathbf{X}$ ) is a random sample if $\left\{y_{i}, \mathbf{x}_{i}\right\}$ is i.i.d. (independently and identically distributed) across observations. Since by Assumption $1.1 \varepsilon_{i}$ is a function of ( $y_{i}, \mathbf{x}_{i}$ ) and since ( $y_{i}, \mathbf{x}_{i}$ ) is independent of ( $y_{j}, \mathbf{x}_{j}$ ) for $j \neq i$, ( $\varepsilon_{i}, \mathbf{x}_{i}$ ) is independent of $\mathbf{x}_{j}$ for $j \neq i$. So


\begin{align*}
\mathrm{E}\left(\varepsilon_{i} \mid \mathbf{X}\right) & =\mathrm{E}\left(\varepsilon_{i} \mid \mathbf{x}_{i}\right) \\
\mathrm{E}\left(\varepsilon_{i}^{2} \mid \mathbf{X}\right) & =\mathrm{E}\left(\varepsilon_{i}^{2} \mid \mathbf{x}_{i}\right) \\
\text { and } \mathrm{E}\left(\varepsilon_{i} \varepsilon_{j} \mid \mathbf{X}\right) & \left.=\mathrm{E}\left(\varepsilon_{i} \mid \mathbf{x}_{i}\right) \mathrm{E}\left(\varepsilon_{j} \mid \mathbf{x}_{j}\right) \quad \text { (for } i \neq j\right) \tag{1.1.15}
\end{align*}


(Proving the last equality in (1.1.15) is a review question.) Therefore, Assumptions 1.2 and 1.4 reduce to

\[
\begin{array}{ll}
\text { Assumption 1.2: } & \mathrm{E}\left(\varepsilon_{i} \mid \mathbf{x}_{i}\right)=0 \quad(i=1,2, \ldots, n), \\
\text { Assumption 1.4: } & \mathrm{E}\left(\varepsilon_{i}^{2} \mid \mathbf{x}_{i}\right)=\sigma^{2}>0 \quad(i=1,2, \ldots, n) . \tag{1.1.17}
\end{array}
\]

The implication of the identical distribution aspect of a random sample is that the joint distribution of $\left(\varepsilon_{i}, \mathbf{x}_{i}\right)$ does not depend on $i$. So the unconditional second moment $\mathrm{E}\left(\varepsilon_{i}^{2}\right)$ is constant across $i$ (this is referred to as unconditional homoskedasticity) and the functional form of the conditional second moment $\mathrm{E}\left(\varepsilon_{i}^{2} \mid \mathbf{x}_{i}\right)$ is the same across $i$. However, Assumption 1.4 -that the value of the conditional\\
second moment is the same across $i$ - does not follow. Therefore, Assumption 1.4 remains restrictive for the case of a random sample; without it, the conditional second moment $\mathrm{E}\left(\varepsilon_{i}^{2} \mid \mathbf{x}_{i}\right)$ can differ across $i$ through its possible dependence on $\mathbf{x}_{i}$. To emphasize the distinction, the restrictions on the conditional second moments, (1.1.12) and (1.1.17), are referred to as conditional homoskedasticity.

\section*{"Fixed" Regressors}
We have presented the classical linear regression model, treating the regressors as random. This is in contrast to the treatment in most textbooks, where $\mathbf{X}$ is assumed to be "fixed" or deterministic. If $\mathbf{X}$ is fixed, then there is no need to distinguish between the conditional distribution of the error term, $f\left(\varepsilon_{i} \mid \mathbf{x}_{1}, \ldots, \mathbf{x}_{n}\right)$, and the unconditional distribution, $f\left(\varepsilon_{i}\right)$, so that Assumptions 1.2 and 1.4 can be written as

\[
\begin{array}{ll}
\text { Assumption 1.2: } & \mathrm{E}\left(\varepsilon_{i}\right)=0(i=1, \ldots, n), \\
\text { Assumption 1.4: } & \mathrm{E}\left(\varepsilon_{i}^{2}\right)=\sigma^{2}(i=1, \ldots, n) ; \\
& \mathrm{E}\left(\varepsilon_{i} \varepsilon_{j}\right)=0(i, j=1, \ldots, n ; i \neq j) . \tag{1.1.19}
\end{array}
\]

Although it is clearly inappropriate for a nonexperimental science like econometrics, the assumption of fixed regressors remains popular because the regression model with fixed $\mathbf{X}$ can be interpreted as a set of statements conditional on $\mathbf{X}$, allowing us to dispense with "| $\mathbf{X}$ " from the statements such as Assumptions 1.2 and 1.4 of the model.

However, the economy in the notation comes at a price. It is very easy to miss the point that the error term is being assumed to be uncorrelated with current, past, and future regressors. Also, the distinction between the unconditional and conditional homoskedasticity gets lost if the regressors are deterministic. Throughout this book, the regressors are treated as random, and, unless otherwise noted, statements conditional on $\mathbf{X}$ are made explicit by inserting "| $\mathbf{X}$."

\section*{QUESTIONS FOR REVIEW}
\begin{enumerate}
  \item (Change in units in the semi-log form) In the wage equation, (1.1.3), of Example 1.2, if WAGE is measured in cents rather than in dollars, what difference does it make to the equation? Hint: $\log (x y)=\log (x)+\log (y)$.
  \item Prove the last equality in (1.1.15). Hint: $\mathrm{E}\left(\varepsilon_{i} \varepsilon_{j} \mid \mathbf{X}\right)=\mathrm{E}\left[\varepsilon_{j} \mathrm{E}\left(\varepsilon_{i} \mid \mathbf{X}, \varepsilon_{j}\right) \mid \mathbf{X}\right]$. ( $\varepsilon_{i}, \mathbf{x}_{i}$ ) is independent of ( $\varepsilon_{j}, \mathbf{x}_{1}, \ldots, \mathbf{x}_{i-1}, \mathbf{x}_{i+1}, \ldots, \mathbf{x}_{n}$ ) for $i \neq j$.
  \item (Combining linearity and strict exogeneity) Show that Assumptions 1.1 and 1.2 imply
\end{enumerate}


\begin{equation*}
\mathrm{E}\left(y_{i} \mid \mathbf{X}\right)=\mathbf{x}_{i}^{\prime} \boldsymbol{\beta} \quad(i=1,2, \ldots, n) . \tag{1.1.20}
\end{equation*}


Conversely, show that this assumption implies that there exist error terms that satisfy those two assumptions.\\
4. (Normally distributed random sample) Consider a random sample on consumption and disposable income, $\left(C O N_{i}, Y D_{i}\right)(i=1,2, \ldots, n)$. Suppose the joint distribution of $\left(C O N_{i}, Y D_{i}\right)$ (which is the same across $i$ because of the random sample assumption) is normal. Clearly, Assumption 1.3 is satisfied; the rank of $\mathbf{X}$ would be less than $K$ only by pure accident. Show that the other assumptions, Assumptions 1.1, 1.2, and 1.4, are satisfied. Hint: If two random variables, $y$ and $x$, are jointly normally distributed, then the conditional expectation is linear in $x$, i.e.,

$$
\mathrm{E}(y \mid x)=\beta_{1}+\beta_{2} x,
$$

and the conditional variance, $\operatorname{Var}(y \mid x)$, does not depend on $x$. Here, the fact that the distribution is the same across $i$ is important; if the distribution differed across $i$, $\beta_{1}$ and $\beta_{2}$ could vary across $i$.\\
5. (Multicollinearity for the simple regression model) Show that Assumption 1.3 for the simple regression model is that the nonconstant regressor ( $x_{i 2}$ ) is really nonconstant (i.e., $x_{i 2} \neq x_{j 2}$ for some pairs of ( $i, j$ ), $i \neq j$, with probability one).\\
6. (An exercise in conditional and unconditional expectations) Show that Assumptions 1.2 and 1.4 imply


\begin{align*}
\operatorname{Var}\left(\varepsilon_{i}\right)=\sigma^{2} & (i=1,2, \ldots, n) \\
\text { and } \operatorname{Cov}\left(\varepsilon_{i}, \varepsilon_{j}\right)=0 & (i \neq j ; i, j=1,2, \ldots n) . \tag{*}
\end{align*}


Hint: Strict exogeneity implies $\mathrm{E}\left(\varepsilon_{i}\right)=0$. So (*) is equivalent to

$$
\begin{gathered}
\mathrm{E}\left(\varepsilon_{i}^{2}\right)=\sigma^{2} \quad(i=1,2, \ldots, n) \\
\text { and } \quad \mathrm{E}\left(\varepsilon_{i} \varepsilon_{j}\right)=0 \quad(i \neq j ; i, j=1,2, \ldots, n) .
\end{gathered}
$$

\subsection*{1.2 The Algebra of Least Squares}
This section describes the computational procedure for obtaining the OLS estimate, $\mathbf{b}$, of the unknown coefficient vector $\boldsymbol{\beta}$ and introduces a few concepts that derive from $\mathbf{b}$.

\section*{OLS Minimizes the Sum of Squared Residuals}
Although we do not observe the error term, we can calculate the value implied by a hypothetical value, $\widetilde{\boldsymbol{\beta}}$, of $\boldsymbol{\beta}$ as

$$
y_{i}-\mathbf{x}_{i}^{\prime} \tilde{\boldsymbol{\beta}}
$$

This is called the residual for observation $i$. From this, form the sum of squared residuals (SSR):

$$
S S R(\tilde{\boldsymbol{\beta}}) \equiv \sum_{i=1}^{n}\left(y_{i}-\mathbf{x}_{i}^{\prime} \tilde{\boldsymbol{\beta}}\right)^{2}=(\mathbf{y}-\mathbf{X} \tilde{\boldsymbol{\beta}})^{\prime}(\mathbf{y}-\mathbf{X} \tilde{\boldsymbol{\beta}})
$$

This sum is also called the error sum of squares (ESS) or the residual sum of squares (RSS). It is a function of $\boldsymbol{\widetilde { \boldsymbol { \beta } }}$ because the residual depends on it. The OLS estimate, $\mathbf{b}$, of $\boldsymbol{\beta}$ is the $\widetilde{\boldsymbol{\beta}}$ that minimizes this function:


\begin{equation*}
\mathbf{b} \equiv \underset{\widetilde{\boldsymbol{\beta}}}{\operatorname{argmin}} S S R(\widetilde{\boldsymbol{\beta}}) \tag{1.2.1}
\end{equation*}


The relationship among $\boldsymbol{\beta}$ (the unknown coefficient vector), $\mathbf{b}$ (the OLS estimate of it), and $\boldsymbol{\tilde { \boldsymbol { \beta } }}$ (a hypothetical value of $\boldsymbol{\beta}$ ) is illustrated in Figure 1.1 for $K=1$. Because $S S R(\tilde{\boldsymbol{\beta}})$ is quadratic in $\tilde{\boldsymbol{\beta}}$, its graph has the U shape. The value of $\widetilde{\boldsymbol{\beta}}$ corresponding to the bottom is $\mathbf{b}$, the OLS estimate. Since it depends on the sample ( $\mathbf{y}, \mathbf{X}$ ), the OLS estimate $\mathbf{b}$ is in general different from the true value $\boldsymbol{\beta}$; if $\mathbf{b}$ equals $\boldsymbol{\beta}$, it is by sheer accident.

By having squared residuals in the objective function, this method imposes a heavy penalty on large residuals; the OLS estimate is chosen to prevent large residuals for a few observations at the expense of tolerating relatively small residuals for many other observations. We will see in the next section that this particular criterion brings about some desirable properties for the estimate.

\begin{figure}[h]
\begin{center}
  \includegraphics[alt={},max width=\textwidth]{14fb0795-6d81-4e76-9788-792298f1ddae-033_699_1109_369_494}
\captionsetup{labelformat=empty}
\caption{Figure 1.1: Hypothetical, True, and Estimated Values}
\end{center}
\end{figure}

\section*{Normal Equations}
A sure-fire way of solving the minimization problem is to derive the first-order conditions by setting the partial derivatives equal to zero. To this end we seek a $K$-dimensional vector of partial derivatives, $\partial S S R(\tilde{\boldsymbol{\beta}}) / \partial \tilde{\boldsymbol{\beta}}$. ${ }^{6}$ The task is facilitated by writing $\operatorname{SSR}(\widetilde{\boldsymbol{\beta}})$ as

$$
\begin{aligned}
\operatorname{SSR}(\tilde{\boldsymbol{\beta}}) & =(\mathbf{y}-\mathbf{X} \tilde{\boldsymbol{\beta}})^{\prime}(\mathbf{y}-\mathbf{X} \tilde{\boldsymbol{\beta}}) \quad\left(\text { since the } i \text {-th element of } \mathbf{y}-\mathbf{X} \tilde{\boldsymbol{\beta}} \text { is } y_{i}-\mathbf{x}_{i}^{\prime} \tilde{\boldsymbol{\beta}}\right) \\
& =\left(\mathbf{y}^{\prime}-\tilde{\boldsymbol{\beta}}^{\prime} \mathbf{X}^{\prime}\right)(\mathbf{y}-\mathbf{X} \tilde{\boldsymbol{\beta}}) \quad\left(\text { since }(\mathbf{X} \tilde{\boldsymbol{\beta}})^{\prime}=\tilde{\boldsymbol{\beta}}^{\prime} \mathbf{X}^{\prime}\right) \\
& =\mathbf{y}^{\prime} \mathbf{y}-\tilde{\boldsymbol{\beta}}^{\prime} \mathbf{X}^{\prime} \mathbf{y}-\mathbf{y}^{\prime} \mathbf{X} \tilde{\boldsymbol{\beta}}+\tilde{\boldsymbol{\beta}}^{\prime} \mathbf{X}^{\prime} \mathbf{X} \tilde{\boldsymbol{\beta}} \\
& =\mathbf{y}^{\prime} \mathbf{y}-2 \mathbf{y}^{\prime} \mathbf{X} \tilde{\boldsymbol{\beta}}+\tilde{\boldsymbol{\beta}}^{\prime} \mathbf{X}^{\prime} \mathbf{X} \tilde{\boldsymbol{\beta}}
\end{aligned}
$$

(since the scalar $\widetilde{\boldsymbol{\beta}}^{\prime} \mathbf{X}^{\prime} \mathbf{y}$ equals its transpose $\mathbf{y}^{\prime} \mathbf{X} \widetilde{\boldsymbol{\beta}}$ )


\begin{equation*}
\equiv \mathbf{y}^{\prime} \mathbf{y}-2 \mathbf{a}^{\prime} \tilde{\boldsymbol{\beta}}+\tilde{\boldsymbol{\beta}}^{\prime} \mathbf{A} \tilde{\boldsymbol{\beta}} \quad \text { with } \mathbf{a} \equiv \mathbf{X}^{\prime} \mathbf{y} \text { and } \mathbf{A} \equiv \mathbf{X}^{\prime} \mathbf{X} . \tag{1.2.2}
\end{equation*}


The term $\mathbf{y}^{\prime} \mathbf{y}$ does not depend on $\tilde{\boldsymbol{\beta}}$ and so can be ignored in the differentiation of $S S R(\widetilde{\boldsymbol{\beta}})$. Recalling from matrix algebra that

$$
\frac{\partial\left(\mathbf{a}^{\prime} \tilde{\boldsymbol{\beta}}\right)}{\partial \tilde{\boldsymbol{\beta}}}=\mathbf{a} \text { and } \frac{\partial\left(\tilde{\boldsymbol{\beta}}^{\prime} \mathbf{A} \tilde{\boldsymbol{\beta}}\right)}{\partial \tilde{\boldsymbol{\beta}}}=2 \mathbf{A} \tilde{\boldsymbol{\beta}} \quad \text { for } \mathbf{A} \text { symmetric, }
$$

\footnotetext{${ }^{6}$ If $h: \mathbb{R}^{K} \rightarrow \mathbb{R}$ is a scalar-valued function of a $K$-dimensional vector $\mathbf{x}$, the derivative of $h$ with respect to $\mathbf{x}$ is a $K$-dimensional vector whose $k$-th element is $\partial h(\mathbf{x}) / \partial x_{k}$ where $x_{k}$ is the $k$-th element of $\mathbf{x}$. (This $K$-dimensional vector is called the gradient.) Here, the $\mathbf{x}$ is $\widetilde{\boldsymbol{\beta}}$ and the function $h(\mathbf{x})$ is $\operatorname{SSR}(\widetilde{\boldsymbol{\beta}})$.
}
the $K$-dimensional vector of partial derivatives is
$$
\frac{\partial S S R(\tilde{\boldsymbol{\beta}})}{\partial \widetilde{\boldsymbol{\beta}}}=-2 \mathbf{a}+2 \mathbf{A} \tilde{\boldsymbol{\beta}}
$$

The first-order conditions are obtained by setting this equal to zero. Recalling from (1.2.2) that $\mathbf{a}$ here is $\mathbf{X}^{\prime} \mathbf{y}$ and $\mathbf{A}$ is $\mathbf{X}^{\prime} \mathbf{X}$ and rearranging, we can write the first-order conditions as


\begin{equation*}
\underset{(K \times K)(K \times 1)}{\mathbf{X}^{\prime} \mathbf{X}}=\mathbf{X}^{\prime} \mathbf{y} \tag{1.2.3}
\end{equation*}


Here, we have replaced $\widetilde{\boldsymbol{\beta}}$ by $\mathbf{b}$ because the OLS estimate $\mathbf{b}$ is the $\widetilde{\boldsymbol{\beta}}$ that satisfies the first-order conditions. These $K$ equations are called the normal equations.

The vector of residuals evaluated at $\widetilde{\boldsymbol{\beta}}=\mathbf{b}$,


\begin{equation*}
\underset{(n \times 1)}{\mathbf{e}} \equiv \mathbf{y}-\mathbf{X} \mathbf{b} \tag{1.2.4}
\end{equation*}


is called the vector of OLS residuals. Its $i$-th element is $e_{i} \equiv y_{i}-\mathbf{x}_{i}^{\prime} \mathbf{b}$. Rearranging (1.2.3) gives


\begin{gather*}
\mathbf{X}^{\prime}(\mathbf{y}-\mathbf{X b})=\mathbf{0} \text { or } \mathbf{X}^{\prime} \mathbf{e}=\mathbf{0} \text { or } \\
\frac{1}{n} \sum_{i=1}^{n} \mathbf{x}_{i} \cdot e_{i}=\mathbf{0} \text { or } \frac{1}{n} \sum_{i=1}^{n} \mathbf{x}_{i} \cdot\left(y_{i}-\mathbf{x}_{i}^{\prime} \mathbf{b}\right)=\mathbf{0} \tag{1.2.3'}
\end{gather*}


which shows that the normal equations can be interpreted as the sample analogue of the orthogonality conditions $\mathrm{E}\left(\mathbf{x}_{i} \cdot \varepsilon_{i}\right)=\mathbf{0}$. This point will be pursued more fully in subsequent chapters.

To be sure, the first-order conditions are just a necessary condition for minimization, and we have to check the second-order condition to make sure that $\mathbf{b}$ achieves the minimum, not the maximum. Those who are familiar with the Hessian of a function of several variables ${ }^{7}$ can immediately recognize that the second-order condition is satisfied because (as noted below) $\mathbf{X}^{\prime} \mathbf{X}$ is positive definite. There is, however, a more direct way to show that $\mathbf{b}$ indeed achieves the minimum. It utilizes the "add-and-subtract" strategy, which is effective when the objective function is quadratic, as here. Application of the strategy to the algebra of least squares is left to you as an analytical exercise.

\footnotetext{${ }^{7}$ The Hessian of $h(\mathbf{x})$ is a square matrix whose $(k, \ell)$ element is $\partial^{2} h(\mathbf{x}) / \partial x_{k} \partial x_{\ell}$.
}\section*{Two Expressions for the OLS Estimator}
Thus, we have obtained a system of $K$ linear simultaneous equations in $K$ unknowns in $\mathbf{b}$. By Assumption 1.3 (no multicollinearity), the coefficient matrix $\mathbf{X}^{\prime} \mathbf{X}$ is positive definite (see review question 1 below for a proof) and hence nonsingular. So the normal equations can be solved uniquely for $\mathbf{b}$ by premultiplying both sides of (1.2.3) by $\left(\mathbf{X}^{\prime} \mathbf{X}\right)^{-1}$ :


\begin{equation*}
\mathbf{b}=\left(\mathbf{X}^{\prime} \mathbf{X}\right)^{-1} \mathbf{X}^{\prime} \mathbf{y} . \tag{1.2.5}
\end{equation*}


Viewed as a function of the sample ( $\mathbf{y}, \mathbf{X}$ ), (1.2.5) is sometimes called the OLS estimator. For any given sample ( $\mathbf{y}, \mathbf{X}$ ), the value of this function is the OLS estimate. In this book, as in most other textbooks, the two terms will be used almost interchangeably.

Since $\left(\mathbf{X}^{\prime} \mathbf{X}\right)^{-1} \mathbf{X}^{\prime} \mathbf{y}=\left(\mathbf{X}^{\prime} \mathbf{X} / n\right)^{-1} \mathbf{X}^{\prime} \mathbf{y} / n$, the OLS estimator can also be rewritten as


\begin{equation*}
\mathbf{b}=\mathbf{S}_{\mathbf{x x}}^{-1} \mathbf{s}_{\mathbf{x y}}, \tag{1.2.5́}
\end{equation*}


where


\begin{align*}
& \mathbf{S}_{\mathbf{x} \mathbf{x}}=\frac{1}{n} \mathbf{X}^{\prime} \mathbf{X}=\frac{1}{n} \sum_{i=1}^{n} \mathbf{x}_{i} \mathbf{x}_{i}^{\prime} \quad\left(\text { sample average of } \mathbf{x}_{i} \mathbf{x}_{i}^{\prime}\right),  \tag{1.2.6a}\\
& \mathbf{s}_{\mathbf{x y}}=\frac{1}{n} \mathbf{X}^{\prime} \mathbf{y}=\frac{1}{n} \sum_{i=1}^{n} \mathbf{x}_{i} \cdot y_{i} \quad\left(\text { sample average of } \mathbf{x}_{i} \cdot y_{i}\right) . \tag{1.2.6b}
\end{align*}


The data matrix form (1.2.5) is more convenient for developing the finite-sample results, while the sample average form (1.2.5') is the form to be utilized for largesample theory.

\section*{More Concepts and Algebra}
Having derived the OLS estimator of the coefficient vector, we can define a few related concepts.

\begin{itemize}
  \item The fitted value for observation $i$ is defined as $\hat{y}_{i} \equiv \mathbf{x}_{i}^{\prime} \mathbf{b}$. The vector of fitted value, $\hat{\mathbf{y}}$, equals Xb. Thus, the vector of OLS residuals can be written as $\mathbf{e}= \mathbf{y}-\hat{\mathbf{y}}$.
  \item The projection matrix $\mathbf{P}$ and the annihilator $\mathbf{M}$ are defined as
\end{itemize}


\begin{gather*}
\underset{(n \times n)}{\mathbf{P}} \equiv \mathbf{X}\left(\mathbf{X}^{\prime} \mathbf{X}\right)^{-1} \mathbf{X}^{\prime}  \tag{1.2.7}\\
\underset{(n \times n)}{\mathbf{M}} \equiv \mathbf{I}_{n}-\mathbf{P} \tag{1.2.8}
\end{gather*}


They have the following nifty properties (proving them is a review question):


\begin{align*}
& \text { Both } \mathbf{P} \text { and } \mathbf{M} \text { are symmetric and idempotent, }{ }^{8}  \tag{1.2.9}\\
& \mathbf{P X}=\mathbf{X} \text { (hence the term projection matrix), }  \tag{1.2.10}\\
& \quad \mathbf{M X}=\mathbf{0} \text { (hence the term annihilator). } \tag{1.2.11}
\end{align*}


Since $\mathbf{e}$ is the residual vector at $\widetilde{\boldsymbol{\beta}}=\mathbf{b}$, the sum of squared OLS residuals, SSR, equals $\mathbf{e}^{\prime} \mathbf{e}$. It can further be written as


\begin{equation*}
S S R=\mathbf{e}^{\prime} \mathbf{e}=\boldsymbol{\varepsilon}^{\prime} \mathbf{M} \boldsymbol{\varepsilon} \tag{1.2.12}
\end{equation*}


(Proving this is a review question.) This expression, relating SSR to the true error term $\boldsymbol{\varepsilon}$, will be useful later on.

\begin{itemize}
  \item The OLS estimate of $\sigma^{2}$ (the variance of the error term), denoted $s^{2}$, is the sum of squared residuals divided by $n-K$ :
\end{itemize}


\begin{equation*}
s^{2} \equiv \frac{S S R}{n-K}=\frac{\mathbf{e}^{\prime} \mathbf{e}}{n-K} \tag{1.2.13}
\end{equation*}


(The definition presumes that $n>K$; otherwise $s^{2}$ is not well-defined.) As will be shown in Proposition 1.2 below, dividing the sum of squared residuals by $n-K$ (called the degrees of freedom) rather than by $n$ (the sample size) makes this estimate unbiased for $\sigma^{2}$. The intuitive reason is that $K$ parameters ( $\boldsymbol{\beta}$ ) have to be estimated before obtaining the residual vector e used to calculate $s^{2}$. More specifically, $\mathbf{e}$ has to satisfy the $K$ normal equations (1.2.3'), which limits the variability of the residual.

\begin{itemize}
  \item The square root of $s^{2}, s$, is called the standard error of the regression (SER) or standard error of the equation (SEE). It is an estimate of the standard deviation of the error term.
  \item The sampling error is defined as $\mathbf{b}-\boldsymbol{\beta}$. It too can be related to $\boldsymbol{\varepsilon}$ as follows.
\end{itemize}


\begin{align*}
\mathbf{b}-\boldsymbol{\beta} & =\left(\mathbf{X}^{\prime} \mathbf{X}\right)^{-1} \mathbf{X}^{\prime} \mathbf{y}-\boldsymbol{\beta} \quad(\text { by }(1.2 .5)) \\
& =\left(\mathbf{X}^{\prime} \mathbf{X}\right)^{-1} \mathbf{X}^{\prime}(\mathbf{X} \boldsymbol{\beta}+\boldsymbol{\varepsilon})-\boldsymbol{\beta} \quad(\text { since } \mathbf{y}=\mathbf{X} \boldsymbol{\beta}+\boldsymbol{\varepsilon} \text { by Assumption 1.1 }) \\
& =\left(\mathbf{X}^{\prime} \mathbf{X}\right)^{-1}\left(\mathbf{X}^{\prime} \mathbf{X}\right) \boldsymbol{\beta}+\left(\mathbf{X}^{\prime} \mathbf{X}\right)^{-1} \mathbf{X}^{\prime} \boldsymbol{\varepsilon}-\boldsymbol{\beta} \\
& =\boldsymbol{\beta}+\left(\mathbf{X}^{\prime} \mathbf{X}\right)^{-1} \mathbf{X}^{\prime} \boldsymbol{\varepsilon}-\boldsymbol{\beta}=\left(\mathbf{X}^{\prime} \mathbf{X}\right)^{-1} \mathbf{X}^{\prime} \boldsymbol{\varepsilon} . \tag{1.2.14}
\end{align*}


\footnotetext{${ }^{8}$ A square matrix $\mathbf{A}$ is said to be idempotent if $\mathbf{A}=\mathbf{A}^{2}$.
}\begin{itemize}
  \item Uncentered $\boldsymbol{R}^{\mathbf{2}}$. One measure of the variability of the dependent variable is the sum of squares, $\sum y_{i}^{2}=\mathbf{y}^{\prime} \mathbf{y}$. Because the OLS residual is chosen to satisfy the normal equations, we have the following decomposition of $\mathbf{y}^{\prime} \mathbf{y}$ :
\end{itemize}


\begin{align*}
\mathbf{y}^{\prime} \mathbf{y} & =(\hat{\mathbf{y}}+\mathbf{e})^{\prime}(\hat{\mathbf{y}}+\mathbf{e}) \quad(\text { since } \mathbf{e}=\mathbf{y}-\hat{\mathbf{y}}) \\
& =\hat{\mathbf{y}}^{\prime} \hat{\mathbf{y}}+2 \hat{\mathbf{y}}^{\prime} \mathbf{e}+\mathbf{e}^{\prime} \mathbf{e} \\
& =\hat{\mathbf{y}}^{\prime} \hat{\mathbf{y}}+2 \mathbf{b}^{\prime} \mathbf{X}^{\prime} \mathbf{e}+\mathbf{e}^{\prime} \mathbf{e} \quad(\text { since } \hat{\mathbf{y}} \equiv \mathbf{X b}) \\
& =\hat{\mathbf{y}}^{\prime} \hat{\mathbf{y}}+\mathbf{e}^{\prime} \mathbf{e} \quad\left(\text { since } \mathbf{X}^{\prime} \mathbf{e}=\mathbf{0} \text { by the normal equations; see }\left(1.2 .3^{\prime}\right)\right) \tag{1.2.15}
\end{align*}


The uncentered $\boldsymbol{R}^{\mathbf{2}}$ is defined as


\begin{equation*}
R_{u c}^{2} \equiv 1-\frac{\mathbf{e}^{\prime} \mathbf{e}}{\mathbf{y}^{\prime} \mathbf{y}} \tag{1.2.16}
\end{equation*}


Because of the decomposition (1.2.15), this equals

$$
\frac{\hat{\mathbf{y}}^{\prime} \hat{\mathbf{y}}}{\mathbf{y}^{\prime} \mathbf{y}}
$$

Since both $\hat{\mathbf{y}}^{\prime} \hat{\mathbf{y}}$ and $\mathbf{e}^{\prime} \mathbf{e}$ are nonnegative, $0 \leq R_{u c}^{2} \leq 1$. Thus, the uncentered $R^{2}$ has the interpretation of the fraction of the variation of the dependent variable that is attributable to the variation in the explanatory variables. The closer the fitted value tracks the dependent variable, the closer is the uncentered $R^{2}$ to one.

\begin{itemize}
  \item (Centered) $\boldsymbol{R}^{\mathbf{2}}$, the coefficient of determination. If the only regressor is a constant (so that $K=1$ and $x_{i 1}=1$ ), then it is easy to see from (1.2.5) that $\mathbf{b}$ equals $\bar{y}$, the sample mean of the dependent variable, which means that $\hat{y}_{i}=\bar{y}$ for all $i, \hat{\mathbf{y}}^{\prime} \hat{\mathbf{y}}$ in (1.2.15) equals $n \bar{y}^{2}$, and $\mathbf{e}^{\prime} \mathbf{e}$ equals $\sum_{i}\left(y_{i}-\bar{y}\right)^{2}$. If the regressors also include nonconstant variables, then it can be shown (the proof is left as an analytical exercise) that $\sum_{i}\left(y_{i}-\bar{y}\right)^{2}$ is decomposed as
\end{itemize}


\begin{equation*}
\sum_{i=1}^{n}\left(y_{i}-\bar{y}\right)^{2}=\sum_{i=1}^{n}\left(\hat{y}_{i}-\bar{y}\right)^{2}+\sum_{i=1}^{n} e_{i}^{2} \text { with } \bar{y} \equiv \frac{1}{n} \sum_{i=1}^{n} y_{i} . \tag{1.2.17}
\end{equation*}


The coefficient of determination, $R^{2}$, is defined as


\begin{equation*}
R^{2} \equiv 1-\frac{\sum_{i=1}^{n} e_{i}^{2}}{\sum_{i=1}^{n}\left(y_{i}-\bar{y}\right)^{2}} \tag{1.2.18}
\end{equation*}


Because of the decomposition (1.2.17), this $R^{2}$ equals

$$
\frac{\sum_{i=1}^{n}\left(\hat{y}_{i}-\bar{y}\right)^{2}}{\sum_{i=1}^{n}\left(y_{i}-\bar{y}\right)^{2}} .
$$

Therefore, provided that the regressors include a constant so that the decomposition (1.2.17) is valid, $0 \leq R^{2} \leq 1$. Thus, this $R^{2}$ as defined in (1.2.18) is a measure of the explanatory power of the nonconstant regressors.

If the regressors do not include a constant but (as some regression software packages do) you nevertheless calculate $R^{2}$ by the formula (1.2.18), then the $R^{2}$ can be negative. This is because, without the benefit of an intercept, the regression could do worse than the sample mean in terms of tracking the dependent variable. On the other hand, some other regression packages (notably STATA) switch to the formula (1.2.16) for the $R^{2}$ when a constant is not included, in order to avoid negative values for the $R^{2}$. This is a mixed blessing. Suppose that the regressors do not include a constant but that a linear combination of the regressors equals a constant. This occurs if, for example, the intercept is replaced by seasonal dummies. ${ }^{9}$ The regression is essentially the same when one of the regressors in the linear combination is replaced by a constant. Indeed, one should obtain the same vector of fitted values. But if the formula for the $R^{2}$ is (1.2.16) for regressions without a constant and (1.2.18) for those with a constant, the calculated $R^{2}$ declines (see Review Question 7 below) after the replacement by a constant.

\section*{Influential Analysis (optional)}
Since the method of least squares seeks to prevent a few large residuals at the expense of incurring many relatively small residuals, only a few observations can be extremely influential in the sense that dropping them from the sample changes some elements of $\mathbf{b}$ substantially. There is a systematic way to find those influential observations. ${ }^{10}$ Let $\mathbf{b}^{(i)}$ be the OLS estimate of $\boldsymbol{\beta}$ that would be obtained if OLS were used on a sample from which the $i$-th observation was omitted. The key equation is


\begin{equation*}
\mathbf{b}^{(i)}-\mathbf{b}=-\left(\frac{1}{1-p_{i}}\right)\left(\mathbf{X}^{\prime} \mathbf{X}\right)^{-1} \mathbf{x}_{i} \cdot e_{i}, \tag{1.2.19}
\end{equation*}


\footnotetext{${ }^{9}$ Dummy variables will be introduced in the empirical exercise for this chapter.\\
${ }^{10}$ See Krasker, Kuh, and Welsch (1983) for more details.
}
where $\mathbf{x}_{i}$ as before is the $i$-th row of $\mathbf{X}, e_{i}$ is the OLS residual for observation $i$, and $p_{i}$ is defined as

\begin{equation*}
p_{i} \equiv \mathbf{x}_{i}^{\prime}\left(\mathbf{X}^{\prime} \mathbf{X}\right)^{-1} \mathbf{x}_{i} \tag{1.2.20}
\end{equation*}

which is the $i$-th diagonal element of the projection matrix $\mathbf{P}$. (Proving (1.2.19) would be a good exercise in matrix algebra, but we will not do it here.) It is easy to show (see Review Question 7 of Section 1.3) that

\begin{equation*}
0 \leq p_{i} \leq 1 \text { and } \sum_{i=1}^{n} p_{i}=K \tag{1.2.21}
\end{equation*}


So $p_{i}$ equals $K / n$ on average.\\
To illustrate the use of (1.2.19) in a specific example, consider the relationship between equipment investment and economic growth for the world's poorest countries between 1960 and 1985. Figure 1.2 plots the average annual GDP-per-worker growth between 1960 and 1985 against the ratio of equipment investment to GDP over the same period for thirteen countries whose GDP per worker in 1965 was less than 10 percent of that of the United States. ${ }^{11}$ It is clear visually from the plot that the position of the estimated regression line would depend very much on the single outlier (Botswana). Indeed, if Botswana is dropped from the sample, the estimated slope coefficient drops from 0.37 to 0.058 . In the present case of simple regression, it is easy to spot outliers by visually inspecting the plot such as Figure 1.2. This strategy would not work if there were more than one nonconstant regressor. Analysis based on formula (1.2.19) is not restricted to simple regressions. Table 1.1 displays the data along with the OLS residuals, the values of $p_{i}$, and (1.2.19) for each observation. Botswana's $p_{i}$ of 0.7196 is well above the average of 0.154 $(=K / n=2 / 13)$ and is highly influential, as the last two columns of the table indicate. Note that we could not have detected the influential observation by looking at the residuals, which is not surprising because the algebra of least squares is designed to avoid large residuals at the expense of many small residuals for other observations.

What should be done with influential observations? It depends. If the influential observations satisfy the regression model, they provide valuable information about the regression function unavailable from the rest of the sample and should definitely be kept in the sample. But more probable is that the influential observations are atypical of the rest of the sample because they do not satisfy the model.

\footnotetext{${ }^{11}$ The data are from the Penn World Table, reprinted in DeLong and Summers (1991). To their credit, their analysis is based on the whole sample of sixty-one countries.
}\begin{figure}[h]
\begin{center}
  \includegraphics[alt={},max width=\textwidth]{14fb0795-6d81-4e76-9788-792298f1ddae-040_622_1079_369_498}
\captionsetup{labelformat=empty}
\caption{Figure 1.2: Equipment Investment and Growth}
\end{center}
\end{figure}

In this case they should definitely be dropped from the sample. For the example just examined, there was a worldwide growth in the demand for diamonds, Botswana's main export, and production of diamonds requires heavy investment in drilling equipment. If the reason to expect an association between growth and equipment investment is the beneficial effect on productivity of the introduction of new technologies through equipment, then Botswana, whose high GDP growth is demand-driven, should be dropped from the sample.

\section*{A Note on the Computation of OLS Estimates ${ }^{12}$}
So far, we have focused on the conceptual aspects of the algebra of least squares. But for applied researchers who actually calculate OLS estimates using digital computers, it is important to be aware of a certain aspect of digital computing in order to avoid the risk of obtaining unreliable estimates without knowing it. The source of a potential problem is that the computer approximates real numbers by so-called floating-point numbers. When an arithmetic operation involves both very large numbers and very small numbers, floating-point calculation can produce inaccurate results. This is relevant in the computation of OLS estimates when the regressors greatly differ in magnitude. For example, one of the regressors may be the interest rate stated as a fraction, and another may be U.S. GDP in dollars. The matrix $\mathbf{X}^{\prime} \mathbf{X}$ will then contain both very small and very large numbers, and the arithmetic operation of inverting this matrix by the digital computer will produce unreliable results.

\footnotetext{${ }^{12}$ A fuller treatment of this topic can be found in Section 1.5 of Davidson and MacKinnon (1993).
}\begin{table}[h]
\begin{center}
\captionsetup{labelformat=empty}
\caption{Table 1.1: Influential Analysis}
\begin{tabular}{|l|l|l|l|l|l|l|}
\hline
Country & GDP/worker growth & Equipment/ GDP & Residual & $p_{i}$ & (1.2.19) for $\beta_{1}$ & (1.2.19) for $\beta_{2}$ \\
\hline
Botswana & 0.0676 & 0.1310 & 0.0119 & 0.7196 & 0.0104 & -0.3124 \\
\hline
Cameroon & 0.0458 & 0.0415 & 0.0233 & 0.0773 & -0.0021 & 0.0045 \\
\hline
Ethiopia & 0.0094 & 0.0212 & -0.0056 & 0.1193 & 0.0010 & -0.0119 \\
\hline
India & 0.0115 & 0.0278 & -0.0059 & 0.0980 & 0.0009 & -0.0087 \\
\hline
Indonesia & 0.0345 & 0.0221 & 0.0192 & 0.1160 & -0.0034 & 0.0394 \\
\hline
Ivory Coast & 0.0278 & 0.0243 & 0.0117 & 0.1084 & -0.0019 & 0.0213 \\
\hline
Kenya & 0.0146 & 0.0462 & -0.0096 & 0.0775 & 0.0007 & 0.0023 \\
\hline
Madagascar & -0.0102 & 0.0219 & -0.0254 & 0.1167 & 0.0045 & -0.0527 \\
\hline
Malawi & 0.0153 & 0.0361 & -0.0052 & 0.0817 & 0.0006 & -0.0036 \\
\hline
Mali & 0.0044 & 0.0433 & -0.0188 & 0.0769 & 0.0016 & -0.0006 \\
\hline
Pakistan & 0.0295 & 0.0263 & 0.0126 & 0.1022 & -0.0020 & 0.0205 \\
\hline
Tanzania & 0.0184 & 0.0860 & -0.0206 & 0.2281 & -0.0021 & 0.0952 \\
\hline
Thailand & 0.0341 & 0.0395 & 0.0123 & 0.0784 & -0.0012 & 0.0047 \\
\hline
\end{tabular}
\end{center}
\end{table}

A simple solution to this problem is to choose the units of measurement so that the regressors are similar in magnitude. For example, state the interest rate in percents and U.S. GDP in trillion dollars. This sort of care would prevent the problem most of the time. A more systematic transformation of the $\mathbf{X}$ matrix is to subtract the sample means of all regressors and divide by the sample standard deviations before forming $\mathbf{X}^{\prime} \mathbf{X}$ (and adjust the OLS estimates to undo the transformation). Most OLS programs (such as TSP) take a more sophisticated transformation of the $\mathbf{X}$ matrix (called the QR decomposition) to produce accurate results.

\section*{QUESTIONS FOR REVIEW}
\begin{enumerate}
  \item Prove that $\mathbf{X}^{\prime} \mathbf{X}$ is positive definite if $\mathbf{X}$ is of full column rank. Hint: What needs to be shown is that $\mathbf{c}^{\prime} \mathbf{X}^{\prime} \mathbf{X} \mathbf{c}>0$ for $\mathbf{c} \neq \mathbf{0}$. Define $\mathbf{z} \equiv \mathbf{X} \mathbf{c}$. Then $\mathbf{c}^{\prime} \mathbf{X}^{\prime} \mathbf{X} \mathbf{c}=\mathbf{z}^{\prime} \mathbf{z}=\sum_{k=1}^{K} z_{i}^{2}$. If $\mathbf{X}$ is of full column rank, then $\mathbf{z} \neq \mathbf{0}$ for any $\mathbf{c} \neq \mathbf{0}$.
  \item Verify that $\mathbf{X}^{\prime} \mathbf{X} / n=\frac{1}{n} \sum_{i} \mathbf{x}_{i} \mathbf{x}_{i}^{\prime}$ and $\mathbf{X}^{\prime} \mathbf{y} / n=\frac{1}{n} \sum_{i} \mathbf{x}_{i} \cdot y_{i}$ as in (1.2.6). Hint: The $(k, \ell)$ element of $\mathbf{X}^{\prime} \mathbf{X}$ is $\sum_{i} x_{i k} x_{i \ell}$.
  \item (OLS estimator for the simple regression model) In the simple regression model, $K=2$ and $x_{i 1}=1$. Show that
\end{enumerate}

$$
\mathbf{S}_{\mathbf{x x}}=\left[\begin{array}{cc}
1 & \bar{x}_{2} \\
\bar{x}_{2} & \frac{1}{n} \sum_{i=1}^{n} x_{i 2}^{2}
\end{array}\right], \quad \mathbf{s}_{\mathbf{x y}}=\left[\begin{array}{c}
\bar{y} \\
\frac{1}{n} \sum_{i=1}^{n} x_{i 2} y_{i}
\end{array}\right]
$$

where

$$
\bar{y} \equiv \frac{1}{n} \sum_{i=1}^{n} y_{i} \quad \text { and } \quad \bar{x}_{2} \equiv \frac{1}{n} \sum_{i=1}^{n} x_{i 2} .
$$

Show that

$$
b_{2}=\frac{\frac{1}{n} \sum_{i=1}^{n}\left(x_{i 2}-\bar{x}_{2}\right)\left(y_{i}-\bar{y}\right)}{\frac{1}{n} \sum_{i=1}^{n}\left(x_{i 2}-\bar{x}_{2}\right)^{2}} \text { and } b_{1}=\bar{y}-\bar{x}_{2} b_{2} .
$$

(You may recognize the denominator of the expression for $b_{2}$ as the sample variance of the nonconstant regressor and the numerator as the sample covariance between the nonconstant regressor and the dependent variable.) Hint:

$$
\frac{1}{n} \sum_{i=1}^{n} x_{i 2}^{2}-\left(\bar{x}_{2}\right)^{2}=\frac{1}{n} \sum_{i=1}^{n}\left(x_{i 2}-\bar{x}_{2}\right)^{2}
$$

and

$$
\frac{1}{n} \sum_{i=1}^{n} x_{i 2} y_{i}-\bar{x}_{2} \bar{y}=\frac{1}{n} \sum_{i=1}^{n}\left(x_{i 2}-\bar{x}_{2}\right)\left(y_{i}-\bar{y}\right) .
$$

You can take (1.2.5') and use the brute force of matrix inversion. Alternatively, write down the two normal equations. The first normal equation is $b_{1}=\bar{y}-\bar{x}_{2} b_{2}$. Substitute this into the second normal equation to eliminate $b_{1}$ and then solve for $b_{2}$.\\
4. Prove (1.2.9)-(1.2.11). Hint: They should easily follow from the definition of $\mathbf{P}$ and $\mathbf{M}$.\\
5. (Matrix algebra of fitted values and residuals) Show the following:\\
(a) $\hat{\mathbf{y}}=\mathbf{P y}, \mathbf{e}=\mathbf{M y}=\mathbf{M} \boldsymbol{\varepsilon}$. Hint: Use (1.2.5).\\
(b) (1.2.12), namely, $S S R=\boldsymbol{\varepsilon}^{\prime} \mathbf{M} \boldsymbol{\varepsilon}$.\\
6. (Change in units and $R^{2}$ ) Does a change in the unit of measurement for the dependent variable change $R^{2}$ ? A change in the unit of measurement for the regressors? Hint: Check whether the change affects the denominator and the numerator in the definition for $R^{2}$.\\
7. (Relation between $R_{u c}^{2}$ and $R^{2}$ ) Show that

$$
1-R^{2}=\left(1+\frac{n \cdot \bar{y}^{2}}{\sum_{i=1}^{n}\left(y_{i}-\bar{y}\right)^{2}}\right)\left(1-R_{u c}^{2}\right)
$$

Hint: Use (1.2.16), (1.2.18), and the identity $\sum_{i}\left(y_{i}-\bar{y}\right)^{2}=\sum_{i} y_{i}^{2}-n \cdot \bar{y}^{2}$.\\
8. Show that

$$
R_{u c}^{2}=\frac{\mathbf{y}^{\prime} \mathbf{P y}}{\mathbf{y}^{\prime} \mathbf{y}} .
$$

\begin{enumerate}
  \setcounter{enumi}{8}
  \item (Computation of the statistics) Verify that $\mathbf{b}, S S R, s^{2}$, and $R^{2}$ can be calculated from the following sample averages: $\mathbf{S}_{\mathbf{x x}}, \mathbf{s}_{\mathbf{x y}}, \mathbf{y}^{\prime} \mathbf{y} / n$, and $\bar{y}$. (If the regressors include a constant, then $\bar{y}$ is the element of $\mathbf{s}_{\mathbf{x y}}$ corresponding to the constant.) Therefore, those sample averages need to be computed just once in order to obtain the regression coefficients and related statistics.
\end{enumerate}

\subsection*{1.3 Finite-Sample Properties of OLS}
Having derived the OLS estimator, we now examine its finite-sample properties, namely, the characteristics of the distribution of the estimator that are valid for any given sample size $n$.

\section*{Finite-Sample Distribution of $\mathbf{b}$}
\section*{Proposition 1.1 (finite-sample properties of the OLS estimator of $\boldsymbol{\beta}$ ):}
(a) (unbiasedness) Under Assumptions 1.1-1.3, $\mathrm{E}(\mathbf{b} \mid \mathbf{X})=\boldsymbol{\beta}$.\\
(b) (expression for the variance) Under Assumptions 1.1-1.4, $\operatorname{Var}(\mathbf{b} \mid \mathbf{X})=\sigma^{2}$. $\left(\mathbf{X}^{\prime} \mathbf{X}\right)^{-1}$.\\
(c) (Gauss-Markov Theorem) Under Assumptions 1.1-1.4, the OLS estimator is efficient in the class of linear unbiased estimators. That is, for any unbiased estimator $\widehat{\boldsymbol{\beta}}$ that is linear in $\mathbf{y}, \operatorname{Var}(\widehat{\boldsymbol{\beta}} \mid \mathbf{X}) \geq \operatorname{Var}(\mathbf{b} \mid \mathbf{X})$ in the matrix sense. ${ }^{13}$\\
(d) Under Assumptions 1.1-1.4, $\operatorname{Cov}(\mathbf{b}, \mathbf{e} \mid \mathbf{X})=\mathbf{0}$, where $\mathbf{e} \equiv \mathbf{y}-\mathbf{X} \mathbf{b}$.

Before plunging into the proof, let us be clear about what this proposition means.

\begin{itemize}
  \item The matrix inequality in part (c) says that the $K \times K$ matrix $\operatorname{Var}(\widehat{\boldsymbol{\beta}} \mid \mathbf{X})- \operatorname{Var}(\mathbf{b} \mid \mathbf{X})$ is positive semidefinite, so
\end{itemize}

$$
\mathbf{a}^{\prime}[\operatorname{Var}(\widehat{\boldsymbol{\beta}} \mid \mathbf{X})-\operatorname{Var}(\mathbf{b} \mid \mathbf{X})] \mathbf{a} \geq 0 \quad \text { or } \quad \mathbf{a}^{\prime} \operatorname{Var}(\widehat{\boldsymbol{\beta}} \mid \mathbf{X}) \mathbf{a} \geq \mathbf{a}^{\prime} \operatorname{Var}(\mathbf{b} \mid \mathbf{X}) \mathbf{a}
$$

for any $K$-dimensional vector $\mathbf{a}$. In particular, consider a special vector whose elements are all 0 except for the $k$-th element, which is 1 . For this particular $\mathbf{a}$, the quadratic form $\mathbf{a}^{\prime} \mathbf{A a}$ picks up the $(k, k)$ element of $\mathbf{A}$. But the $(k, k)$ element of $\operatorname{Var}(\widehat{\boldsymbol{\beta}} \mid \mathbf{X})$, for example, is $\operatorname{Var}\left(\widehat{\beta}_{k} \mid \mathbf{X}\right)$ where $\widehat{\beta}_{k}$ is the $k$-th element of $\widehat{\boldsymbol{\beta}}$. Thus the matrix inequality in (c) implies


\begin{equation*}
\operatorname{Var}\left(\widehat{\beta}_{k} \mid \mathbf{X}\right) \geq \operatorname{Var}\left(b_{k} \mid \mathbf{X}\right) \quad(k=1,2, \ldots, K) \tag{1.3.1}
\end{equation*}


That is, for any regression coefficient, the variance of the OLS estimator is no larger than that of any other linear unbiased estimator.

\footnotetext{${ }^{13}$ Let $\mathbf{A}$ and $\mathbf{B}$ be two square matrices of the same size. We say that $\mathbf{A} \geq \mathbf{B}$ if $\mathbf{A}-\mathbf{B}$ is positive semidefinite. A $K \times K$ matrix $\mathbf{C}$ is said to be positive semidefinite (or nonnegative definite) if $\mathbf{x}^{\prime} \mathbf{C x} \geq 0$ for all $K$-dimensional vectors $\mathbf{x}$.
}\begin{itemize}
  \item As clear from (1.2.5), the OLS estimator is linear in $\mathbf{y}$. There are many other estimators of $\boldsymbol{\beta}$ that are linear and unbiased (you will be asked to provide one in a review question below). The Gauss-Markov Theorem says that the OLS estimator is efficient in the sense that its conditional variance matrix $\operatorname{Var}(\mathbf{b} \mid \mathbf{X})$ is smallest among linear unbiased estimators. For this reason the OLS estimator is called the Best Linear Unbiased Estimator (BLUE).
  \item The OLS estimator $\mathbf{b}$ is a function of the sample $(\mathbf{y}, \mathbf{X})$. Since $(\mathbf{y}, \mathbf{X})$ are random, so is $\mathbf{b}$. Now imagine that we fix $\mathbf{X}$ at some given value, calculate $\mathbf{b}$ for all samples corresponding to all possible realizations of $\mathbf{y}$, and take the average of b (the Monte Carlo exercise to this chapter will ask you to do this). This average is the (population) conditional mean $\mathrm{E}(\mathbf{b} \mid \mathbf{X})$. Part (a) (unbiasedness) says that this average equals the true value $\boldsymbol{\beta}$.
  \item There is another notion of unbiasedness that is weaker than the unbiasedness of part (a). By the Law of Total Expectations, $\mathrm{E}[\mathrm{E}(\mathbf{b} \mid \mathbf{X})]=\mathrm{E}(\mathbf{b})$. So (a) implies
\end{itemize}


\begin{equation*}
\mathrm{E}(\mathbf{b})=\boldsymbol{\beta} . \tag{1.3.2}
\end{equation*}


This says: if we calculated $\mathbf{b}$ for all possible different samples, differing not only in $\mathbf{y}$ but also in $\mathbf{X}$, the average would be the true value. This unconditional statement is probably more relevant in economics because samples do differ in both $\mathbf{y}$ and $\mathbf{X}$. The import of the conditional statement (a) is that it implies the unconditional statement (1.3.2), which is more relevant.

\begin{itemize}
  \item The same holds for the conditional statement (c) about the variance. A review question below asks you to show that statements (a) and (b) imply
\end{itemize}


\begin{equation*}
\operatorname{Var}(\widehat{\boldsymbol{\beta}}) \geq \operatorname{Var}(\mathbf{b}) \tag{1.3.3}
\end{equation*}


where $\widehat{\boldsymbol{\beta}}$ is any linear unbiased estimator (so that $\mathrm{E}(\widehat{\boldsymbol{\beta}} \mid \mathbf{X})=\boldsymbol{\beta}$ ).\\
We will now go through the proof of this important result. The proof may look lengthy; if so, it is only because it records every step, however easy. In the first reading, you can skip the proof of part (c). Proof of (d) is a review question.

\section*{Proof.}
(a) (Proof that $\mathrm{E}(\mathbf{b} \mid \mathbf{X})=\boldsymbol{\beta}) \mathrm{E}(\mathbf{b}-\boldsymbol{\beta} \mid \mathbf{X})=\mathbf{0}$ whenever $\mathrm{E}(\mathbf{b} \mid \mathbf{X})=\boldsymbol{\beta}$. So we prove the former. By the expression for the sampling error (1.2.14),\\
$\mathbf{b}-\boldsymbol{\beta}=\mathbf{A} \boldsymbol{\varepsilon}$ where $\mathbf{A}$ here is $\left(\mathbf{X}^{\prime} \mathbf{X}\right)^{-1} \mathbf{X}^{\prime}$. So

$$
\mathrm{E}(\mathbf{b}-\boldsymbol{\beta} \mid \mathbf{X})=\mathrm{E}(\mathbf{A} \boldsymbol{\varepsilon} \mid \mathbf{X})=\mathbf{A} \mathrm{E}(\boldsymbol{\varepsilon} \mid \mathbf{X})
$$

Here, the second equality holds by the linearity of conditional expectations; A is a function of $\mathbf{X}$ and so can be treated as if nonrandom. Since $\mathrm{E}(\boldsymbol{\varepsilon} \mid \mathbf{X})=\mathbf{0}$, the last expression is zero.\\
(b) (Proof that $\left.\operatorname{Var}(\mathbf{b} \mid \mathbf{X})=\sigma^{2} \cdot\left(\mathbf{X}^{\prime} \mathbf{X}\right)^{-1}\right)$

$$
\begin{aligned}
\operatorname{Var}(\mathbf{b} \mid \mathbf{X}) & =\operatorname{Var}(\mathbf{b}-\boldsymbol{\beta} \mid \mathbf{X}) \quad(\text { since } \boldsymbol{\beta} \text { is not random }) \\
& =\operatorname{Var}(\mathbf{A} \boldsymbol{\varepsilon} \mid \mathbf{X}) \quad\left(\text { by (1.2.14) and } \mathbf{A} \equiv\left(\mathbf{X}^{\prime} \mathbf{X}\right)^{-1} \mathbf{X}^{\prime}\right) \\
& =\mathbf{A} \operatorname{Var}(\boldsymbol{\varepsilon} \mid \mathbf{X}) \mathbf{A}^{\prime} \quad(\text { since } \mathbf{A} \text { is a function of } \mathbf{X}) \\
& =\mathbf{A} E\left(\boldsymbol{\varepsilon} \boldsymbol{\varepsilon}^{\prime} \mid \mathbf{X}\right) \mathbf{A}^{\prime} \quad(\text { by Assumption 1.2 }) \\
& =\mathbf{A}\left(\sigma^{2} \mathbf{I}_{n}\right) \mathbf{A}^{\prime} \quad(\text { by Assumption 1.4, see (1.1.14) }) \\
& =\sigma^{2} \mathbf{A} \mathbf{A}^{\prime} \\
& =\sigma^{2} \cdot\left(\mathbf{X}^{\prime} \mathbf{X}\right)^{-1} \quad\left(\text { since } \mathbf{A} \mathbf{A}^{\prime}=\left(\mathbf{X}^{\prime} \mathbf{X}\right)^{-1} \mathbf{X}^{\prime} \mathbf{X}\left(\mathbf{X}^{\prime} \mathbf{X}\right)^{-1}=\left(\mathbf{X}^{\prime} \mathbf{X}\right)^{-1}\right) .
\end{aligned}
$$

(c) (Gauss-Markov) Since $\widehat{\boldsymbol{\beta}}$ is linear in $\mathbf{y}$, it can be written as $\widehat{\boldsymbol{\beta}}=\mathbf{C y}$ for some matrix $\mathbf{C}$, which possibly is a function of $\mathbf{X}$. Let $\mathbf{D} \equiv \mathbf{C}-\mathbf{A}$ or $\mathbf{C}=\mathbf{D}+\mathbf{A}$ where $\mathbf{A} \equiv\left(\mathbf{X}^{\prime} \mathbf{X}\right)^{-1} \mathbf{X}^{\prime}$. Then

$$
\begin{aligned}
\widehat{\boldsymbol{\beta}} & =(\mathbf{D}+\mathbf{A}) \mathbf{y} \\
& =\mathbf{D} \mathbf{y}+\mathbf{A} \mathbf{y} \\
& =\mathbf{D}(\mathbf{X} \boldsymbol{\beta}+\boldsymbol{\varepsilon})+\mathbf{b} \quad\left(\text { since } \mathbf{y}=\mathbf{X} \boldsymbol{\beta}+\boldsymbol{\varepsilon} \text { and } \mathbf{A} \mathbf{y}=\left(\mathbf{X}^{\prime} \mathbf{X}\right)^{-1} \mathbf{X}^{\prime} \mathbf{y}=\mathbf{b}\right) \\
& =\mathbf{D} \mathbf{X} \boldsymbol{\beta}+\mathbf{D} \boldsymbol{\varepsilon}+\mathbf{b}
\end{aligned}
$$

Taking the conditional expectation of both sides, we obtain

$$
\mathrm{E}(\widehat{\boldsymbol{\beta}} \mid \mathbf{X})=\mathbf{D} \mathbf{X} \boldsymbol{\beta}+\mathrm{E}(\mathbf{D} \boldsymbol{\varepsilon} \mid \mathbf{X})+\mathrm{E}(\mathbf{b} \mid \mathbf{X})
$$

Since both $\mathbf{b}$ and $\widehat{\boldsymbol{\beta}}$ are unbiased and since $\mathrm{E}(\mathbf{D} \boldsymbol{\varepsilon} \mid \mathbf{X})=\mathbf{D} \mathrm{E}(\boldsymbol{\varepsilon} \mid \mathbf{X})=\mathbf{0}$, it follows that $\mathbf{D X} \boldsymbol{\beta}=\mathbf{0}$. For this to be true for any given $\boldsymbol{\beta}$, it is necessary that $\mathbf{D X}=\mathbf{0}$. So $\widehat{\boldsymbol{\beta}}=\mathbf{D} \boldsymbol{\varepsilon}+\mathbf{b}$ and

$$
\begin{aligned}
\widehat{\boldsymbol{\beta}}-\boldsymbol{\beta} & =\mathbf{D} \boldsymbol{\varepsilon}+(\mathbf{b}-\boldsymbol{\beta}) \\
& =(\mathbf{D}+\mathbf{A}) \boldsymbol{\varepsilon} \quad(\text { by }(1.2 .14))
\end{aligned}
$$

So

$$
\begin{aligned}
\operatorname{Var}(\widehat{\boldsymbol{\beta}} \mid \mathbf{X})= & \operatorname{Var}(\widehat{\boldsymbol{\beta}}-\boldsymbol{\beta} \mid \mathbf{X}) \\
= & \operatorname{Var}[(\mathbf{D}+\mathbf{A}) \boldsymbol{\varepsilon} \mid \mathbf{X}] \\
= & (\mathbf{D}+\mathbf{A}) \operatorname{Var}(\boldsymbol{\varepsilon} \mid \mathbf{X})\left(\mathbf{D}^{\prime}+\mathbf{A}^{\prime}\right) \\
& (\text { since both } \mathbf{D} \text { and } \mathbf{A} \text { are functions of } \mathbf{X}) \\
= & \sigma^{2} \cdot(\mathbf{D}+\mathbf{A})\left(\mathbf{D}^{\prime}+\mathbf{A}^{\prime}\right) \quad\left(\text { since } \operatorname{Var}(\boldsymbol{\varepsilon} \mid \mathbf{X})=\sigma^{2} \mathbf{I}_{n}\right) \\
= & \sigma^{2} \cdot\left(\mathbf{D} \mathbf{D}^{\prime}+\mathbf{A D}^{\prime}+\mathbf{D A}^{\prime}+\mathbf{A} \mathbf{A}^{\prime}\right)
\end{aligned}
$$

But $\mathbf{D A}^{\prime}=\mathbf{D X}\left(\mathbf{X}^{\prime} \mathbf{X}\right)^{-1}=\mathbf{0}$ since $\mathbf{D X}=\mathbf{0}$. Also, $\mathbf{A A}^{\prime}=\left(\mathbf{X}^{\prime} \mathbf{X}\right)^{-1}$ as shown in (b). So

$$
\begin{aligned}
\operatorname{Var}(\widehat{\boldsymbol{\beta}} \mid \mathbf{X}) & =\sigma^{2} \cdot\left[\mathbf{D \mathbf { D } ^ { \prime }}+\left(\mathbf{X}^{\prime} \mathbf{X}\right)^{-1}\right] \\
& \geq \sigma^{2} \cdot\left(\mathbf{X}^{\prime} \mathbf{X}\right)^{-1} \quad\left(\text { since } \mathbf{D D}^{\prime} \text { is positive semidefinite }\right) \\
& =\operatorname{Var}(\mathbf{b} \mid \mathbf{X}) \quad(\text { by }(\mathrm{b}))
\end{aligned}
$$

It should be emphasized that the strict exogeneity assumption (Assumption 1.2) is critical for proving unbiasedness. Anything short of strict exogeneity will not do. For example, it is not enough to assume that $\mathrm{E}\left(\varepsilon_{i} \mid \mathbf{x}_{i}\right)=0$ for all $i$ or that $\mathrm{E}\left(\mathbf{x}_{i} \cdot \varepsilon_{i}\right)=\mathbf{0}$ for all $i$. We noted in Section 1.1 that most time-series models do not satisfy strict exogeneity even if they satisfy weaker conditions such as the orthogonality condition $\mathrm{E}\left(\mathbf{x}_{i} \cdot \varepsilon_{i}\right)=\mathbf{0}$. It follows that for those models the OLS estimator is not unbiased.

\section*{Finite-Sample Properties of $\boldsymbol{s}^{\mathbf{2}}$}
We defined the OLS estimator of $\sigma^{2}$ in (1.2.13). It, too, is unbiased.

Proposition 1.2 (Unbiasedness of $\boldsymbol{s}^{\mathbf{2}}$ ): Under Assumptions 1.1-1.4, $\mathrm{E}\left(s^{\mathbf{2}} \mid \mathbf{X}\right)= \sigma^{2}$ (and hence $\mathrm{E}\left(s^{2}\right)=\sigma^{2}$ ), provided $n>K$ (so that $s^{2}$ is well-defined).

We can prove this proposition easily by the use of the trace operator. ${ }^{14}$

Proof. Since $s^{2}=\mathbf{e}^{\prime} \mathbf{e} /(n-K)$, the proof amounts to showing that $\mathrm{E}\left(\mathbf{e}^{\prime} \mathbf{e} \mid \mathbf{X}\right)= (n-K) \sigma^{2}$. As shown in (1.2.12), $\mathbf{e}^{\prime} \mathbf{e}=\boldsymbol{\varepsilon}^{\prime} \mathbf{M} \boldsymbol{\varepsilon}$ where $\mathbf{M}$ is the annihilator. The proof consists of proving two properties: (1) $\mathrm{E}\left(\boldsymbol{\varepsilon}^{\prime} \mathbf{M} \boldsymbol{\varepsilon} \mid \mathbf{X}\right)=\sigma^{2} \cdot$ trace(M), and (2) $\operatorname{trace}(\mathbf{M})=n-K$.

\footnotetext{${ }^{14}$ The trace of a square matrix $\mathbf{A}$ is the sum of the diagonal elements of $\mathbf{A}$ : trace $(\mathbf{A})=\sum_{i} a_{i i}$.
}
(1) (Proof that $\mathrm{E}\left(\boldsymbol{\varepsilon}^{\prime} \mathbf{M} \boldsymbol{\varepsilon} \mid \mathbf{X}\right)=\sigma^{2} \cdot$ trace(M)) Since $\boldsymbol{\varepsilon}^{\prime} \mathbf{M} \boldsymbol{\varepsilon}=\sum_{i=1}^{n} \sum_{j=1}^{n} m_{i j} \varepsilon_{i} \varepsilon_{j}$ (this is just writing out the quadratic form $\boldsymbol{\varepsilon}^{\prime} \mathbf{M} \boldsymbol{\varepsilon}$ ), we have
$$
\begin{aligned}
\mathrm{E}\left(\boldsymbol{\varepsilon}^{\prime} \mathbf{M} \boldsymbol{\varepsilon} \mid \mathbf{X}\right) & =\sum_{i=1}^{n} \sum_{j=1}^{n} m_{i j} \mathrm{E}\left(\varepsilon_{i} \varepsilon_{j} \mid \mathbf{X}\right) \underset{\left.\mathrm{E}\left(m_{i j} \varepsilon_{i} \varepsilon_{j} \mid \mathbf{X}\right)=m_{i j} \mathrm{E}\left(\varepsilon_{i} \varepsilon_{j} \mid \mathbf{X}\right)\right)}{\text { 's are functions of } \mathbf{X}} \\
& =\sum_{i=1}^{n} m_{i i} \sigma^{2} \quad\left(\text { since } \mathrm{E}\left(\varepsilon_{i} \varepsilon_{j} \mid \mathbf{X}\right)=0 \text { for } i \neq j\right. \text { by Assumption 1.4) } \\
& =\sigma^{2} \sum_{i=1}^{n} m_{i i} \\
& =\sigma^{2} \cdot \operatorname{trace}(\mathbf{M})
\end{aligned}
$$
(2) (Proof that trace $(\mathbf{M})=n-K$ )
$$
\begin{aligned}
\operatorname{trace}(\mathbf{M}) & =\operatorname{trace}\left(\mathbf{I}_{n}-\mathbf{P}\right) \quad\left(\text { since } \mathbf{M} \equiv \mathbf{I}_{n}-\mathbf{P} ;\right. \text { see (1.2.8)) } \\
& =\operatorname{trace}\left(\mathbf{I}_{n}\right)-\operatorname{trace}(\mathbf{P}) \quad(\text { fact: the trace operator is linear }) \\
& =n-\operatorname{trace}(\mathbf{P})
\end{aligned}
$$
and
$$
\begin{aligned}
\operatorname{trace}(\mathbf{P}) & =\operatorname{trace}\left[\mathbf{X}\left(\mathbf{X}^{\prime} \mathbf{X}\right)^{-1} \mathbf{X}^{\prime}\right] \quad\left(\text { since } \mathbf{P} \equiv \mathbf{X}\left(\mathbf{X}^{\prime} \mathbf{X}\right)^{-1} \mathbf{X}^{\prime}\right. \text {; see (1.2.7)) } \\
& =\operatorname{trace}\left[\left(\mathbf{X}^{\prime} \mathbf{X}\right)^{-1} \mathbf{X}^{\prime} \mathbf{X}\right] \quad(\text { fact: } \operatorname{trace}(\mathbf{A B})=\operatorname{trace}(\mathbf{B A})) \\
& =\operatorname{trace}\left(\mathbf{I}_{K}\right)=K
\end{aligned}
$$

So trace $(\mathbf{M})=n-K$.

\section*{Estimate of $\operatorname{Var}(\mathbf{b} \mid \mathbf{X})$}
If $s^{2}$ is the estimate of $\sigma^{2}$, a natural estimate of $\operatorname{Var}(\mathbf{b} \mid \mathbf{X})=\sigma^{2} \cdot\left(\mathbf{X}^{\prime} \mathbf{X}\right)^{-1}$ is


\begin{equation*}
\widehat{\operatorname{Var}(\mathbf{b} \mid \mathbf{X}}) \equiv s^{2} \cdot\left(\mathbf{X}^{\prime} \mathbf{X}\right)^{-1} \tag{1.3.4}
\end{equation*}


This is one of the statistics included in the computer printout of any OLS software package.

\section*{QUESTIONS FOR REVIEW}
\begin{enumerate}
  \item (Role of the no-multicollinearity assumption) In Propositions 1.1 and 1.2, where did we use Assumption 1.3 that $\operatorname{rank}(\mathbf{X})=K$ ? Hint: We need the no-multicollinearity condition to make sure $\mathbf{X}^{\prime} \mathbf{X}$ is invertible.
  \item (Example of a linear estimator) For the consumption function example in Example 1.1, propose a linear and unbiased estimator of $\beta_{2}$ that is different from the OLS estimator. Hint: How about $\widehat{\beta}_{2}=\left(C O N_{2}-C O N_{1}\right) /\left(Y D_{2}-Y D_{1}\right)$ ? Is it linear in $\left(C O N_{1}, \ldots, C O N_{n}\right)$ ? Is it unbiased in the sense that $\mathrm{E}\left(\widehat{\beta}_{2}\right.$ | $\left.Y D_{1}, \ldots, Y D_{n}\right)=\beta_{2} ?$
  \item (What Gauss-Markov does not mean) Under Assumptions 1.1-1.4, does there exist a linear, but not necessarily unbiased, estimator of $\boldsymbol{\beta}$ that has a variance smaller than that of the OLS estimator? If so, how small can the variance be?\\
Hint: If an estimator of $\boldsymbol{\beta}$ is a constant, then the estimator is trivially linear in $\mathbf{y}$.
  \item (Gauss-Markov for Unconditional Variance)\\
(a) Prove: $\operatorname{Var}(\widehat{\boldsymbol{\beta}})=\mathrm{E}[\operatorname{Var}(\widehat{\boldsymbol{\beta}} \mid \mathbf{X})]+\operatorname{Var}[\mathrm{E}(\widehat{\boldsymbol{\beta}} \mid \mathbf{X})]$. Hint: By definition,
\end{enumerate}

$$
\operatorname{Var}(\widehat{\boldsymbol{\beta}} \mid \mathbf{X}) \equiv \mathrm{E}\left[(\widehat{\boldsymbol{\beta}}-\mathrm{E}(\widehat{\boldsymbol{\beta}} \mid \mathbf{X}))(\widehat{\boldsymbol{\beta}}-\mathrm{E}(\widehat{\boldsymbol{\beta}} \mid \mathbf{X}))^{\prime} \mid \mathbf{X}\right]
$$

and

$$
\operatorname{Var}[\mathrm{E}(\widehat{\boldsymbol{\beta}} \mid \mathbf{X})] \equiv \mathrm{E}\left\{[\mathrm{E}(\widehat{\boldsymbol{\beta}} \mid \mathbf{X})-\mathrm{E}(\widehat{\boldsymbol{\beta}})][\mathrm{E}(\widehat{\boldsymbol{\beta}} \mid \mathbf{X})-\mathrm{E}(\widehat{\boldsymbol{\beta}})]^{\prime}\right\}
$$

Use the add-and-subtract strategy: take $\widehat{\boldsymbol{\beta}}-\mathrm{E}(\widehat{\boldsymbol{\beta}} \mid \mathbf{X})$ and add and subtract $\mathrm{E}(\widehat{\boldsymbol{\beta}})$.\\
(b) Prove (1.3.3). Hint: If $\operatorname{Var}(\widehat{\boldsymbol{\beta}} \mid \mathbf{X}) \geq \operatorname{Var}(\mathbf{b} \mid \mathbf{X})$, then $\mathrm{E}[\operatorname{Var}(\widehat{\boldsymbol{\beta}} \mid \mathbf{X})] \geq \mathrm{E}[\operatorname{Var}(\mathbf{b} \mid \mathbf{X})]$\\
5. Propose an unbiased estimator of $\sigma^{2}$ if you had data on $\boldsymbol{\varepsilon}$. Hint: How about $\boldsymbol{\varepsilon}^{\prime} \boldsymbol{\varepsilon} / n$ ? Is it unbiased?\\
6. Prove part (d) of Proposition 1.1. Hint: By definition,

$$
\operatorname{Cov}(\mathbf{b}, \mathbf{e} \mid \mathbf{X}) \equiv \mathrm{E}\left\{[\mathbf{b}-\mathrm{E}(\mathbf{b} \mid \mathbf{X})][\mathbf{e}-\mathrm{E}(\mathbf{e} \mid \mathbf{X})]^{\prime} \mid \mathbf{X}\right\}
$$

Since $\mathrm{E}(\mathbf{b} \mid \mathbf{X})=\boldsymbol{\beta}$, we have $\mathbf{b}-\mathrm{E}(\mathbf{b} \mid \mathbf{X})=\mathbf{A} \boldsymbol{\varepsilon}$ where $\mathbf{A}$ here is $\left(\mathbf{X}^{\prime} \mathbf{X}\right)^{-1} \mathbf{X}^{\prime}$. Use $\mathbf{M} \boldsymbol{\varepsilon}=\mathbf{e}$ (see Review Question 5 to Section 1.2) to show that $\mathbf{e}-\mathrm{E}(\mathbf{e} \mid \mathbf{X})=\mathbf{M} \boldsymbol{\varepsilon} \cdot \mathrm{E}\left(\mathbf{A} \boldsymbol{\varepsilon} \boldsymbol{\varepsilon}^{\prime} \mathbf{M} \mid \mathbf{X}\right)=\mathbf{A} \mathrm{E}\left(\boldsymbol{\varepsilon} \boldsymbol{\varepsilon}^{\prime} \mid \mathbf{X}\right) \mathbf{M}$ since both $\mathbf{A}$ and $\mathbf{M}$ are functions of $\mathbf{X}$. Finally, use $\mathbf{M X}=\mathbf{0}$ (see (1.2.11)).\\
7. Prove (1.2.21). Hint: Since $\mathbf{P}$ is positive semidefinite, its diagonal elements are nonnegative. Note that $\sum_{i=1}^{n} p_{i}=\operatorname{trace}(\mathbf{P})$.

\subsection*{1.4 Hypothesis Testing under Normality}
Very often, the economic theory that motivated the regression equation also specifies the values that the regression coefficients should take. Suppose that the underlying theory implies the restriction that $\beta_{2}$ equals 1 . Although Proposition 1.1 guarantees that, on average, $b_{2}$ (the OLS estimate of $\beta_{2}$ ) equals 1 if the restriction is true, $b_{2}$ may not be exactly equal to 1 for a particular sample at hand. Obviously, we cannot conclude that the restriction is false just because the estimate $b_{2}$ differs from 1 . In order for us to decide whether the sampling error $b_{2}-1$ is "too large" for the restriction to be true, we need to construct from the sampling error some test statistic whose probability distribution is known given the truth of the hypothesis. It might appear that doing so requires one to specify the joint distribution of $(\mathbf{X}, \boldsymbol{\varepsilon})$ because, as is clear from (1.2.14), the sampling error is a function of ( $\mathbf{X}, \boldsymbol{\varepsilon}$ ). A surprising fact about the theory of hypothesis testing to be presented in this section is that the distribution can be derived without specifying the joint distribution when the conditional distribution of $\boldsymbol{\varepsilon}$ conditional on $\mathbf{X}$ is normal; there is no need to specify the distribution of $\mathbf{X}$.

In the language of hypothesis testing, the restriction to be tested (such as " $\beta_{2}=1$ ") is called the null hypothesis (or simply the null). It is a restriction on the maintained hypothesis, a set of assumptions which, combined with the null, produces some test statistic with a known distribution. For the present case of testing hypothesis about regression coefficients, only the normality assumption about the conditional distribution of $\boldsymbol{\varepsilon}$ needs to be added to the classical regression model (Assumptions 1.1-1.4) to form the maintained hypothesis (as just noted, there is no need to specify the joint distribution of $(\mathbf{X}, \boldsymbol{\varepsilon}))$. Sometimes the maintained hypothesis is somewhat loosely referred to as "the model." We say that the model is correctly specified if the maintained hypothesis is true. Although too large a value of the test statistic is interpreted as a failure of the null, the interpretation is valid only as long as the model is correctly specified. It is possible that the test statistic does not have the supposed distribution when the null is true but the model is false.

\section*{Normally Distributed Error Terms}
In many applications, the error term consists of many miscellaneous factors not captured by the regressors. The Central Limit Theorem suggests that the error term has a normal distribution. In other applications, the error term is due to errors in measuring the dependent variable. It is known that very often measure-\\
ment errors are normally distributed (in fact, the normal distribution was originally developed for measurement errors). It is therefore worth entertaining the normality assumption:

Assumption 1.5 (normality of the error term): The distribution of $\boldsymbol{\varepsilon}$ conditional on $\mathbf{X}$ is jointly normal.

Recall from probability theory that the normal distribution has several convenient features:

\begin{itemize}
  \item The distribution depends only on the mean and the variance. Thus, once the mean and the variance are known, you can write down the density function. If the distribution conditional on $\mathbf{X}$ is normal, the mean and the variance can depend on $\mathbf{X}$. It follows that, if the distribution conditional on $\mathbf{X}$ is normal and if neither the conditional mean nor the conditional variance depends on $\mathbf{X}$, then the marginal (i.e., unconditional) distribution is the same normal distribution.
  \item In general, if two random variables are independent, then they are uncorrelated, but the converse is not true. However, if two random variables are joint normal, the converse is also true, so that independence and a lack of correlation are equivalent. This carries over to conditional distributions: if two random variables are joint normal and uncorrelated conditional on $\mathbf{X}$, then they are independent conditional on $\mathbf{X}$.
  \item A linear function of random variables that are jointly normally distributed is itself normally distributed. This also carries over to conditional distributions. If the distribution of $\boldsymbol{\varepsilon}$ conditional on $\mathbf{X}$ is normal, then $\mathbf{A} \boldsymbol{\varepsilon}$, where the elements of matrix $\mathbf{A}$ are functions of $\mathbf{X}$, is normal conditional on $\mathbf{X}$.
\end{itemize}

It is thanks to these features of normality that Assumption 1.5 delivers the following properties to be exploited in the derivation of test statistics:

\begin{itemize}
  \item The mean and the variance of the distribution of $\boldsymbol{\varepsilon}$ conditional on $\mathbf{X}$ are already specified in Assumptions 1.2 and 1.4. Therefore, Assumption 1.5 together with Assumptions 1.2 and 1.4 implies that the distribution of $\boldsymbol{\varepsilon}$ conditional on $\mathbf{X}$ is $N\left(\mathbf{0}, \sigma^{2} \mathbf{I}_{n}\right)$ :
\end{itemize}


\begin{equation*}
\boldsymbol{\varepsilon} \mid \mathbf{X} \sim N\left(\mathbf{0}, \sigma^{2} \mathbf{I}_{n}\right) . \tag{1.4.1}
\end{equation*}


Thus, the distribution of $\boldsymbol{\varepsilon}$ conditional on $\mathbf{X}$ does not depend on $\mathbf{X}$. It then follows that $\boldsymbol{\varepsilon}$ and $\mathbf{X}$ are independent. Therefore, in particular, the marginal or unconditional distribution of $\boldsymbol{\varepsilon}$ is $N\left(\mathbf{0}, \sigma^{2} \mathbf{I}_{n}\right)$.

\begin{itemize}
  \item We know from (1.2.14) that the sampling error $\mathbf{b}-\boldsymbol{\beta}$ is linear in $\boldsymbol{\varepsilon}$ given $\mathbf{X}$. Since $\boldsymbol{\varepsilon}$ is normal given $\mathbf{X}$, so is the sampling error. Its mean and variance are given by parts (a) and (b) of Proposition 1.1. Thus, under Assumptions 1.1-1.5,
\end{itemize}


\begin{equation*}
(\mathbf{b}-\boldsymbol{\beta}) \mid \mathbf{X} \sim N\left(\mathbf{0}, \sigma^{2} \cdot\left(\mathbf{X}^{\prime} \mathbf{X}\right)^{-1}\right) . \tag{1.4.2}
\end{equation*}


\section*{Testing Hypotheses about Individual Regression Coefficients}
The type of hypothesis we first consider is about the $k$-th coefficient

$$
\mathrm{H}_{0}: \beta_{k}=\bar{\beta}_{k} .
$$

Here, $\bar{\beta}_{k}$ is some known value specified by the null hypothesis. We wish to test this null against the alternative hypothesis $\mathrm{H}_{1}: \beta_{k} \neq \bar{\beta}_{k}$, at a significance level of $\alpha$. Looking at the $k$-th component of (1.4.2) and imposing the restriction of the null, we obtain

$$
\left(b_{k}-\bar{\beta}_{k}\right) \mid \mathbf{X} \sim N\left(0, \sigma^{2} \cdot\left(\left(\mathbf{X}^{\prime} \mathbf{X}\right)^{-1}\right)_{k k}\right)
$$

where $\left(\left(\mathbf{X}^{\prime} \mathbf{X}\right)^{-1}\right)_{k k}$ is the $(k, k)$ element of $\left(\mathbf{X}^{\prime} \mathbf{X}\right)^{-1}$. So if we define the ratio $z_{k}$ by dividing $b_{k}-\bar{\beta}_{k}$ by its standard deviation


\begin{equation*}
z_{k} \equiv \frac{b_{k}-\bar{\beta}_{k}}{\sqrt{\sigma^{2} \cdot\left(\left(\mathbf{X}^{\prime} \mathbf{X}\right)^{-1}\right)_{k k}}}, \tag{1.4.3}
\end{equation*}


then the distribution of $z_{k}$ is $N(0,1)$ (the standard normal distribution).\\
Suppose for a second that $\sigma^{2}$ is known. Then the statistic $z_{k}$ has some desirable properties as a test statistic. First, its value can be calculated from the sample. Second, its distribution conditional on $\mathbf{X}$ does not depend on $\mathbf{X}$ (which should not be confused with the fact that the value of $z_{k}$ depends on $\mathbf{X}$ ). So $z_{k}$ and $\mathbf{X}$ are independently distributed, and, regardless of the value of $\mathbf{X}$, the distribution of $z_{k}$ is the same as its unconditional distribution. This is convenient because different samples differ not only in $\mathbf{y}$ but also in $\mathbf{X}$. Third, the distribution is known. In particular, it does not depend on unknown parameters (such as $\boldsymbol{\beta}$ ). (If the distribution of a statistic depends on unknown parameters, those parameters are called nuisance parameters.) Using this statistic, we can determine whether or not the sampling error $b_{k}-\bar{\beta}_{k}$ is too large: it is too large if the test statistic takes on a value that is surprising for a realization from the distribution.

If we do not know the true value of $\sigma^{2}$, a natural idea is to replace the nuisance parameter $\sigma^{2}$ by its OLS estimate $s^{2}$. The statistic after the substitution of $s^{2}$ for\\
$\sigma^{2}$ is called the $\boldsymbol{t}$-ratio or the $\boldsymbol{t}$-value. The denominator of this statistic is called the standard error of the OLS estimate of $\beta_{k}$ and is sometimes written as $\operatorname{SE}\left(b_{k}\right)$ :


\begin{equation*}
S E\left(b_{k}\right) \equiv \sqrt{s^{2} \cdot\left(\left(\mathbf{X}^{\prime} \mathbf{X}\right)^{-1}\right)_{k k}}=\sqrt{(k, k) \text { element of } \widehat{\operatorname{Var}(\mathbf{b} \mid \mathbf{X})} \text { in (1.3.4). }} \tag{1.4.4}
\end{equation*}


Since $s^{2}$, being a function of the sample, is a random variable, this substitution changes the distribution of the statistic, but fortunately the changed distribution, too, is known and depends on neither nuisance parameters nor $\mathbf{X}$.

Proposition 1.3 (distribution of the $\boldsymbol{t}$-ratio): Suppose Assumptions 1.1-1.5 hold. Under the null hypothesis $\mathrm{H}_{0}: \beta_{k}=\bar{\beta}_{k}$, the $t$-ratio defined as


\begin{equation*}
t_{k} \equiv \frac{b_{k}-\bar{\beta}_{k}}{S E\left(b_{k}\right)} \equiv \frac{b_{k}-\bar{\beta}_{k}}{\sqrt{s^{2} \cdot\left(\left(\mathbf{X}^{\prime} \mathbf{X}\right)^{-1}\right)_{k k}}} \tag{1.4.5}
\end{equation*}


is distributed as $t(n-K)$ (the $t$ distribution with $n-K$ degrees of freedom).

Proof. We can write

$$
\begin{aligned}
t_{k} & =\frac{b_{k}-\bar{\beta}_{k}}{\sqrt{\sigma^{2} \cdot\left(\left(\mathbf{X}^{\prime} \mathbf{X}\right)^{-1}\right)_{k k}}} \cdot \sqrt{\frac{\sigma^{2}}{s^{2}}}=\frac{z_{k}}{\sqrt{s^{2} / \sigma^{2}}} \\
& =\frac{z_{k}}{\sqrt{\frac{\mathbf{e}^{\prime} \mathbf{e} /(n-K)}{\sigma^{2}}}}=\frac{z_{k}}{\sqrt{\frac{q}{n-K}}}
\end{aligned}
$$

where $q \equiv \mathbf{e}^{\prime} \mathbf{e} / \sigma^{2}$ to reflect the substitution of $s^{2}$ for $\sigma^{2}$. We have already shown that $z_{k}$ is $N(0,1)$. We will show:\\
(1) $q \mid \mathbf{X} \sim \chi^{2}(n-K)$,\\
(2) two random variables $z_{k}$ and $q$ are independent conditional on $\mathbf{X}$.

Then, by the definition of the $t$ distribution, the ratio of $z_{k}$ to $\sqrt{q /(n-K)}$ is distributed as $t$ with $n-K$ degrees of freedom, ${ }^{15}$ and we are done.\\
(1) Since $\mathbf{e}^{\prime} \mathbf{e}=\boldsymbol{\varepsilon}^{\prime} \mathbf{M} \boldsymbol{\varepsilon}$ from (1.2.12), we have

$$
q=\frac{\mathbf{e}^{\prime} \mathbf{e}}{\sigma^{2}}=\frac{\boldsymbol{\varepsilon}^{\prime}}{\sigma} \mathbf{M} \frac{\boldsymbol{\varepsilon}}{\sigma}
$$

\footnotetext{${ }^{15}$ Fact: If $x \sim N(0,1), y \sim \chi^{2}(m)$ and if $x$ and $y$ are independent, then the ratio $x / \sqrt{y / m}$ has the $t$ distribution with $m$ degrees of freedom.
}The middle matrix $\mathbf{M}$, being the annihilator, is idempotent. Also, $\boldsymbol{\varepsilon} / \sigma \mid \mathbf{X} \sim N\left(\mathbf{0}, \mathbf{I}_{n}\right)$ by (1.4.1). Therefore, this quadratic form is distributed as $\chi^{2}$ with degrees of freedom equal to $\operatorname{rank}(\mathbf{M}) .{ }^{16} \operatorname{But} \operatorname{rank}(\mathbf{M})=\operatorname{trace}(\mathbf{M})$, because $\mathbf{M}$ is idempotent. ${ }^{17}$ We have already shown in the proof of Proposition 1.2 that $\operatorname{trace}(\mathbf{M})=n-K$. So $q \mid \mathbf{X} \sim \chi^{2}(n-K)$.\\
(2) Both $\mathbf{b}$ and $\mathbf{e}$ are linear functions of $\boldsymbol{\varepsilon}$ (by (1.2.14) and the fact that $\mathbf{e}=\mathbf{M} \boldsymbol{\varepsilon}$ ), so they are jointly normal conditional on $\mathbf{X}$. Also, they are uncorrelated conditional on $\mathbf{X}$ (see part (d) of Proposition 1.1). So $\mathbf{b}$ and $\mathbf{e}$ are independently distributed conditional on $\mathbf{X}$. But $z_{k}$ is a function of $\mathbf{b}$ and $q$ is a function of $\mathbf{e}$. So $z_{k}$ and $q$ are independently distributed conditional on $\mathbf{X}$. ${ }^{18}$

\section*{Decision Rule for the $\boldsymbol{t}$-Test}
The test of the null hypothesis based on the $t$-ratio is called the $t$-test and proceeds as follows:

Step 1: Given the hypothesized value, $\bar{\beta}_{k}$, of $\beta_{k}$, form the $t$-ratio as in (1.4.5). Too large a deviation of $t_{k}$ from 0 is a sign of the failure of the null hypothesis. The next step specifies how large is too large.\\
Step 2: Go to the $t$-table (most statistics and econometrics textbooks include the $t$ table) and look up the entry for $n-K$ degrees of freedom. Find the critical value, $t_{\alpha / 2}(n-K)$, such that the area in the $t$ distribution to the right of $t_{\alpha / 2}(n-K)$ is $\alpha / 2$, as illustrated in Figure 1.3. (If $n-K=30$ and $\alpha=5 \%$, for example, $t_{\alpha / 2}(n-K)=2.042$.) Then, since the $t$ distribution is symmetric around 0 ,

$$
\operatorname{Prob}\left(-t_{\alpha / 2}(n-K)<t<t_{\alpha / 2}(n-K)\right)=1-\alpha .
$$

Step 3: Accept $\mathrm{H}_{0}$ if $-t_{\alpha / 2}(n-K)<t_{k}<t_{\alpha / 2}(n-K)$ (that is, if $\left|t_{k}\right|<t_{\alpha / 2}(n- K)$ ), where $t_{k}$ is the $t$-ratio from Step 1. Reject $\mathrm{H}_{0}$ otherwise. Since $t_{k} \sim t(n-K)$ under $\mathrm{H}_{0}$, the probability of rejecting $\mathrm{H}_{0}$ when $\mathrm{H}_{0}$ is true is $\alpha$. So the size (significance level) of the test is indeed $\alpha$.

A convenient feature of the $t$-test is that the critical value does not depend on $\mathbf{X}$; there is no need to calculate critical values for each sample.

\footnotetext{${ }^{16}$ Fact: If $\mathbf{x} \sim N\left(\mathbf{0}, \mathbf{I}_{n}\right)$ and $\mathbf{A}$ is idempotent, then $\mathbf{x}^{\prime} \mathbf{A} \mathbf{x}$ has a chi-squared distribution with degrees of freedom equal to the rank of $\mathbf{A}$.\\
${ }^{17}$ Fact: If $\mathbf{A}$ is idempotent, then $\operatorname{rank}(\mathbf{A})=\operatorname{trace}(\mathbf{A})$.\\
${ }^{18}$ Fact: If $\mathbf{x}$ and $\mathbf{y}$ are independently distributed, then so are $f(\mathbf{x})$ and $g(\mathbf{y})$.
}\begin{figure}[h]
\begin{center}
  \includegraphics[alt={},max width=\textwidth]{14fb0795-6d81-4e76-9788-792298f1ddae-055_643_1125_369_485}
\captionsetup{labelformat=empty}
\caption{Figure 1.3: $t$ Distribution}
\end{center}
\end{figure}

\section*{Confidence Interval}
Step 3 can also be stated in terms of $b_{k}$ and $S E\left(b_{k}\right)$. Since $t_{k}$ is as in (1.4.5), you accept $\mathrm{H}_{0}$ whenever

$$
-t_{\alpha / 2}(n-K)<\frac{b_{k}-\bar{\beta}_{k}}{S E\left(b_{k}\right)}<t_{\alpha / 2}(n-K)
$$

or

$$
b_{k}-S E\left(b_{k}\right) \cdot t_{\alpha / 2}(n-K)<\bar{\beta}_{k}<b_{k}+S E\left(b_{k}\right) \cdot t_{\alpha / 2}(n-K) .
$$

Therefore, we accept if and only if the hypothesized value $\bar{\beta}_{k}$ falls in the interval:


\begin{equation*}
\left[b_{k}-S E\left(b_{k}\right) \cdot t_{\alpha / 2}(n-K), b_{k}+S E\left(b_{k}\right) \cdot t_{\alpha / 2}(n-K)\right] . \tag{1.4.6}
\end{equation*}


This interval is called the level 1 - $\boldsymbol{\alpha}$ confidence interval. It is narrower the smaller the standard error. Thus, the smallness of the standard error is a measure of the estimator's precision.

\section*{$\boldsymbol{p}$-Value}
The decision rule of the $t$-test can also be stated using the $\boldsymbol{p}$-value.\\
Step 1: Same as above.\\
Step 2: Rather than finding the critical value $t_{\alpha / 2}(n-K)$, calculate

$$
p=\operatorname{Prob}\left(t>\left|t_{k}\right|\right) \times 2
$$

Since the $t$ distribution is symmetric around $0, \operatorname{Prob}\left(t>\left|t_{k}\right|\right)=\operatorname{Prob}(t< -\left|t_{k}\right|$ ), so


\begin{equation*}
\operatorname{Prob}\left(-\left|t_{k}\right|<t<\left|t_{k}\right|\right)=1-p . \tag{1.4.7}
\end{equation*}


Step 3: Accept $\mathrm{H}_{0}$ if $p>\alpha$. Reject otherwise.\\
To see the equivalence of the two decision rules, one based on the critical values such as $t_{\alpha / 2}(n-K)$ and the other based on the $p$-value, refer to Figure 1.3. If $\operatorname{Prob}\left(t>\left|t_{k}\right|\right)$ is greater than $\alpha / 2$ (as in the figure), that is, if the $p$-value is more than $\alpha$, then $\left|t_{k}\right|$ must be to the left of $t_{\alpha / 2}(n-K)$. This means from Step 3 that the null hypothesis is not rejected. Thus, when $p$ is small, the $t$-ratio is surprisingly large for a random variable from the $t$ distribution. The smaller the $p$, the stronger the rejection.

Examples of the $t$-test can be found in Section 1.7.

\section*{Linear Hypotheses}
The null hypothesis we wish to test may not be a restriction about individual regression coefficients of the maintained hypothesis; it is often about linear combinations of them written as a system of linear equations:


\begin{equation*}
\mathrm{H}_{0}: \mathbf{R} \boldsymbol{\beta}=\mathbf{r}, \tag{1.4.8}
\end{equation*}


where values of $\mathbf{R}$ and $\mathbf{r}$ are known and specified by the hypothesis. We denote the number of equations, which is the dimension of $\mathbf{r}$, by $\# \mathbf{r}$. So $\mathbf{R}$ is $\# \mathbf{r} \times K$. These \# $\mathbf{r}$ equations are restrictions on the coefficients in the maintained hypothesis. It is called a linear hypothesis because each equation is linear. To make sure that there are no redundant equations and that the equations are consistent with each other, we require that $\operatorname{rank}(\mathbf{R})=\# \mathbf{r}$ (i.e., $\mathbf{R}$ is of full row rank with its rank equaling the number of rows). But do not be too conscious about the rank condition; in specific applications, it is very easy to spot a failure of the rank condition if there is one.

Example 1.5 (continuation of Example 1.2): Consider the wage equation of Example 1.2 where $K=4$. We might wish to test the hypothesis that education and tenure have equal impact on the wage rate and that there is no experience effect. The hypothesis is two equations (so \# $\mathbf{r}=2$ ):

$$
\beta_{2}=\beta_{3} \text { and } \beta_{4}=0
$$

This can be cast in the format $\mathbf{R} \boldsymbol{\beta}=\mathbf{r}$ if $\mathbf{R}$ and $\mathbf{r}$ are defined as

$$
\mathbf{R}=\left[\begin{array}{cccc}
0 & 1 & -1 & 0 \\
0 & 0 & 0 & 1
\end{array}\right], \quad \mathbf{r}=\left[\begin{array}{l}
0 \\
0
\end{array}\right]
$$

Because the two rows of this $\mathbf{R}$ are linearly independent, the rank condition is satisfied.

But suppose we require additionally that

$$
\beta_{2}-\beta_{3}=\beta_{4}
$$

This is redundant because it holds whenever the first two equations do. With these three equations, $\# \mathbf{r}=3$ and

$$
\mathbf{R}=\left[\begin{array}{cccc}
0 & 1 & -1 & 0 \\
0 & 0 & 0 & 1 \\
0 & 1 & -1 & -1
\end{array}\right], \quad \mathbf{r}=\left[\begin{array}{l}
0 \\
0 \\
0
\end{array}\right] .
$$

Since the third row of $\mathbf{R}$ is the difference between the first two, $\mathbf{R}$ is not of full row rank. The consequence of adding redundant equations is that $\mathbf{R}$ no longer meets the full row rank condition.

As an example of inconsistent equations, consider adding to the first two equations the third equation $\beta_{4}=0.5$. Evidently, $\beta_{4}$ cannot be 0 and 0.5 at the same time. The hypothesis is inconsistent because there is no $\boldsymbol{\beta}$ that satisfies the three equations simultaneously. If we nevertheless included this equation, then $\mathbf{R}$ and $\mathbf{r}$ would become

$$
\mathbf{R}=\left[\begin{array}{cccc}
0 & 1 & -1 & 0 \\
0 & 0 & 0 & 1 \\
0 & 0 & 0 & 1
\end{array}\right], \quad \mathbf{r}=\left[\begin{array}{c}
0 \\
0 \\
0.5
\end{array}\right]
$$

Again, the full row rank condition is not satisfied because the rank of $\mathbf{R}$ is 2 while $\# \mathbf{r}=3$.

\section*{The $\boldsymbol{F}$-Test}
To test linear hypotheses, we look for a test statistic that has a known distribution under the null hypothesis.

Proposition 1.4 (distribution of the $\boldsymbol{F}$-ratio): Suppose Assumptions 1.1-1.5 hold. Under the null hypothesis $\mathrm{H}_{0}: \mathbf{R} \boldsymbol{\beta}=\mathbf{r}$, where $\mathbf{R}$ is $\# \mathbf{r} \times K$ with $\operatorname{rank}(\mathbf{R})= \# \mathbf{r}$, the $\boldsymbol{F}$-ratio defined as


\begin{align*}
F & \equiv \frac{(\mathbf{R} \mathbf{b}-\mathbf{r})^{\prime}\left[\mathbf{R}\left(\mathbf{X}^{\prime} \mathbf{X}\right)^{-1} \mathbf{R}^{\prime}\right]^{-1}(\mathbf{R} \mathbf{b}-\mathbf{r}) / \# \mathbf{r}}{s^{2}} \\
& =(\mathbf{R} \mathbf{b}-\mathbf{r})^{\prime}\left[\mathbf{R} \operatorname{Var}(\mathbf{b} \mid \mathbf{X}) \mathbf{R}^{\prime}\right]^{-1}(\mathbf{R} \mathbf{b}-\mathbf{r}) / \# \mathbf{r} \quad(\text { by (1.3.4)) } \tag{1.4.9}
\end{align*}


is distributed as $F(\# \mathbf{r}, n-K)$ (the $F$ distribution with $\# \mathbf{r}$ and $n-K$ degrees of freedom).

As in Proposition 1.3, it suffices to show that the distribution conditional on $\mathbf{X}$ is $F(\# \mathbf{r}, n-K)$; because the $F$ distribution does not depend on $\mathbf{X}$, it is also the unconditional distribution of the statistic.

Proof. Since $s^{2}=\mathbf{e}^{\prime} \mathbf{e} /(n-K)$, we can write

$$
F=\frac{w / \# \mathbf{r}}{q /(n-K)}
$$

where

$$
w \equiv(\mathbf{R} \mathbf{b}-\mathbf{r})^{\prime}\left[\sigma^{2} \cdot \mathbf{R}\left(\mathbf{X}^{\prime} \mathbf{X}\right)^{-1} \mathbf{R}^{\prime}\right]^{-1}(\mathbf{R} \mathbf{b}-\mathbf{r}) \quad \text { and } \quad q \equiv \frac{\mathbf{e}^{\prime} \mathbf{e}}{\sigma^{2}}
$$

We need to show\\
(1) $w \mid \mathbf{X} \sim \chi^{2}(\# \mathbf{r})$,\\
(2) $q \mid \mathbf{X} \sim \chi^{2}(n-K)$ (this is part (1) in the proof of Proposition 1.3),\\
(3) $w$ and $q$ are independently distributed conditional on $\mathbf{X}$.

Then, by the definition of the $F$ distribution, the $F$-ratio $\sim F(\# \mathbf{r}, n-K)$.\\
(1) Let $\mathbf{v} \equiv \mathbf{R b}-\mathbf{r}$. Under $\mathrm{H}_{0}, \mathbf{R b}-\mathbf{r}=\mathbf{R}(\mathbf{b}-\boldsymbol{\beta})$. So by (1.4.2), conditional on $\mathbf{X}, \mathbf{v}$ is normal with mean $\mathbf{0}$, and its variance is given by

$$
\operatorname{Var}(\mathbf{v} \mid \mathbf{X})=\operatorname{Var}(\mathbf{R}(\mathbf{b}-\boldsymbol{\beta}) \mid \mathbf{X})=\mathbf{R} \operatorname{Var}(\mathbf{b}-\boldsymbol{\beta} \mid \mathbf{X}) \mathbf{R}^{\prime}=\sigma^{2} \cdot \mathbf{R}\left(\mathbf{X}^{\prime} \mathbf{X}\right)^{-1} \mathbf{R}^{\prime}
$$

which is none other than the inverse of the middle matrix in the quadratic form for $w$. Hence, $w$ can be written as $\mathbf{v}^{\prime} \operatorname{Var}(\mathbf{v} \mid \mathbf{X})^{-1} \mathbf{v}$. Since $\mathbf{R}$ is of full row rank and $\mathbf{X}^{\prime} \mathbf{X}$ is nonsingular, $\sigma^{2} \cdot \mathbf{R}\left(\mathbf{X}^{\prime} \mathbf{X}\right)^{-1} \mathbf{R}^{\prime}$ is nonsingular (why? Showing this is a review question). Therefore, by the definition of the $\chi^{2}$ distribution, $w \mid \mathbf{X} \sim \chi^{2}(\# \mathbf{r}) .^{19}$

\footnotetext{${ }^{19}$ Fact: Let $\mathbf{x}$ be an $m$ dimensional random vector. If $\mathbf{x} \sim N(\boldsymbol{\mu}, \boldsymbol{\Sigma})$ with $\boldsymbol{\Sigma}$ nonsingular, then $(\mathbf{x}-\boldsymbol{\mu})^{\prime} \boldsymbol{\Sigma}^{-1}(\mathbf{x}- \boldsymbol{\mu}) \sim \chi^{2}(m)$.
}
(3) $w$ is a function of $\mathbf{b}$ and $q$ is a function of $\mathbf{e}$. But $\mathbf{b}$ and $\mathbf{e}$ are independently distributed conditional on $\mathbf{X}$, as shown in part (2) of the proof of Proposition 1.3. So $w$ and $q$ are independently distributed conditional on $\mathbf{X}$.

If the null hypothesis $\mathbf{R} \boldsymbol{\beta}=\mathbf{r}$ is true, we expect $\mathbf{R} \mathbf{b}-\mathbf{r}$ to be small, so large values of $F$ should be taken as evidence for a failure of the null. This means that we look at only the upper tail of the distribution in the $F$-statistic. The decision rule of the $F$-test at the significance level of $\alpha$ is as follows.

Step 1: Calculate the $F$-ratio by the formula (1.4.9).\\
Step 2: Go to the table of $F$ distribution and look up the entry for \#1 (the numerator degrees of freedom) and $n-K$ (the denominator degrees of freedom). Find the critical value $F_{\alpha}(\# \mathbf{r}, n-K)$ that leaves $\alpha$ for the upper tail of the $F$ distribution, as illustrated in Figure 1.4. For example, when $\# \mathbf{r}=3$, $n-K=30$, and $\alpha=5 \%$, the critical value $F_{.05}(3,30)$ is 2.92 .\\
Step 3: Accept the null if the $F$-ratio from Step 1 is less than $F_{\alpha}(\# \mathbf{r}, n-K)$. Reject otherwise.

This decision rule can also be described in terms of the $p$-value:\\
Step 1: Same as above.\\
Step 2: Calculate\\
$p=$ area of the upper tail of the $F$ distribution to the right of the $F$-ratio.\\
Step 3: Accept the null if $p>\alpha$; reject otherwise.\\
Thus, a small $p$-value is a signal of the failure of the null.

\section*{A More Convenient Expression for $\boldsymbol{F}$}
The above derivation of the $F$-ratio is by the Wald principle, because it is based on the unrestricted estimator, which is not constrained to satisfy the restrictions of the null hypothesis. Calculating the $F$-ratio by the formula (1.4.9) requires matrix inversion and multiplication. Fortunately, there is a convenient alternative formula involving two different sum of squared residuals: one is SSR, the minimized sum of squared residuals obtained from (1.2.1) now denoted as $S S R_{U}$, and the other is the restricted sum of squared residuals, denoted $S S R_{R}$, obtained from


\begin{equation*}
\min _{\widetilde{\boldsymbol{\beta}}} S S R(\widetilde{\boldsymbol{\beta}}) \text { s.t. } \mathbf{R} \widetilde{\boldsymbol{\beta}}=\mathbf{r} \tag{1.4.10}
\end{equation*}


\begin{figure}[h]
\begin{center}
  \includegraphics[alt={},max width=\textwidth]{14fb0795-6d81-4e76-9788-792298f1ddae-060_628_1116_369_481}
\captionsetup{labelformat=empty}
\caption{Figure 1.4: $F$ Distribution}
\end{center}
\end{figure}

Finding the $\widetilde{\boldsymbol{\beta}}$ that achieves this constrained minimization is called the restricted regression or restricted least squares. It is left as an analytical exercise to show that the $F$-ratio equals


\begin{equation*}
F=\frac{\left(S S R_{R}-S S R_{U}\right) / \# \mathbf{r}}{S S R_{U} /(n-K)}, \tag{1.4.11}
\end{equation*}


which is the difference in the objective function deflated by the estimate of the error variance. This derivation of the $F$-ratio is analogous to how the likelihoodratio statistic is derived in maximum likelihood estimation as the difference in log likelihood with and without the imposition of the null hypothesis. For this reason, this second derivation of the $F$-ratio is said to be by the Likelihood-Ratio principle. There is a closed-form expression for the restricted least squares estimator of $\boldsymbol{\beta}$. Deriving the expression is left as an analytical exercise. The computation of restricted least squares will be explained in the context of the empirical example in Section 1.7.

\section*{$\boldsymbol{t}$ versus $\boldsymbol{F}$}
Because hypotheses about individual coefficients are linear hypotheses, the $t$-test of $\mathrm{H}_{0}: \beta_{k}=\bar{\beta}_{k}$ is a special case of the $F$-test. To see this, note that the hypothesis can be written as $\mathbf{R} \boldsymbol{\beta}=\mathbf{r}$ with

$$
\underset{(1 \times K)}{\mathbf{R}}=\left[\begin{array}{lllllll}
0 & \cdots & 0 & 1 & 0 & \cdots & 0
\end{array}\right], \quad \mathbf{r}=\bar{\beta}_{k} .
$$

So by (1.4.9) the $F$-ratio is

$$
F=\left(b_{k}-\bar{\beta}_{k}\right)\left[s^{2} \cdot(k, k) \text { element of }\left(\mathbf{X}^{\prime} \mathbf{X}\right)^{-1}\right]^{-1}\left(b_{k}-\bar{\beta}_{k}\right)
$$

which is the square of the $t$-ratio in (1.4.5). Since a random variable distributed as $F(1, n-K)$ is the square of a random variable distributed as $t(n-K)$, the $t$ - and $F$-tests give the same test result.

Sometimes, the null is that a set of individual regression coefficients equal certain values. For example, assume $K=2$ and consider

$$
\mathrm{H}_{0}: \beta_{1}=1 \text { and } \beta_{2}=0 .
$$

This can be written as a linear hypothesis $\mathbf{R} \boldsymbol{\beta}=\mathbf{r}$ for $\mathbf{R}=\mathbf{I}_{2}$ and $\mathbf{r}=(1,0)^{\prime}$. So the $F$-test can be used. It is tempting, however, to conduct the $t$-test separately for each individual coefficient of the hypothesis. We might accept $\mathrm{H}_{0}$ if both restrictions $\beta_{1}=1$ and $\beta_{2}=0$ pass the $t$-test. This amounts to using the confidence region of

$$
\begin{aligned}
&\left\{\left(\beta_{1}, \beta_{2}\right) \mid\right. b_{1}-S E\left(b_{1}\right) \cdot t_{\alpha / 2}(n-K)<\beta_{1}<b_{1}+S E\left(b_{1}\right) \cdot t_{\alpha / 2}(n-K), \\
&\left.b_{2}-S E\left(b_{2}\right) \cdot t_{\alpha / 2}(n-K)<\beta_{2}<b_{2}+S E\left(b_{2}\right) \cdot t_{\alpha / 2}(n-K)\right\},
\end{aligned}
$$

which is a rectangular region in the ( $\beta_{1}, \beta_{2}$ ) plane, as illustrated in Figure 1.5. If $(1,0)$, the point in the ( $\beta_{1}, \beta_{2}$ ) plane specified by the null, falls in this region, one would accept the null. On the other hand, the confidence region for the $F$-test is

$$
\left.\left\{\left(\beta_{1}, \beta_{2}\right) \mid\left(b_{1}-\beta_{1}, b_{2}-\beta_{2}\right)(\widehat{\operatorname{Var}(\mathbf{b} \mid \mathbf{X}})\right)^{-1}\left[\begin{array}{l}
b_{1}-\beta_{1} \\
b_{2}-\beta_{2}
\end{array}\right]<2 F_{\alpha}(\# \mathbf{r}, n-K)\right\}
$$

Since $\widehat{\operatorname{Var}(\mathbf{b} \mid \mathbf{X})}$ is positive definite, the $F$-test acceptance region is an ellipse in the ( $\beta_{1}, \beta_{2}$ ) plane. The two confidence regions look typically like Figure 1.5.

The $F$-test should be preferred to the test using two $t$-ratios for two reasons. First, if the size (significance level) in each of the two $t$-tests is $\alpha$, then the overall size (the probability that $(1,0)$ is outside the rectangular region) is not $\alpha$. Second, as will be noted in the next section (see (1.5.19)), the $F$-test is a likelihood ratio test and likelihood-ratio tests have certain desirable properties. So even if the significance level in each $t$-test is controlled so that the overall size is $\alpha$, the test is less desirable than the $F$-test. ${ }^{20}$

\footnotetext{${ }^{20}$ For more details on the relationship between the $t$-test and the $F$-tests, see Scheffe (1959, p. 46).
}\begin{figure}[h]
\begin{center}
  \includegraphics[alt={},max width=\textwidth]{14fb0795-6d81-4e76-9788-792298f1ddae-062_717_1127_373_474}
\captionsetup{labelformat=empty}
\caption{Figure 1.5: $t$ - versus $F$-Tests}
\end{center}
\end{figure}

\section*{An Example of a Test Statistic Whose Distribution Depends on X}
To place the discussion of this section in a proper perspective, it may be useful to note that there are some statistics whose conditional distribution depends on $\mathbf{X}$. Consider the celebrated Durbin-Watson statistic:

$$
\frac{\sum_{i=2}^{n}\left(e_{i}-e_{i-1}\right)^{2}}{\sum_{i=1}^{n} e_{i}^{2}} .
$$

The conditional distribution, and hence the critical values, of this statistic depend on X, but J. Durbin and G. S. Watson have shown that the critical values fall between two bounds (which depends on the sample size, the number of regressors, and whether the regressor includes a constant). Therefore, the critical values for the unconditional distribution, too, fall between these bounds.

The statistic is designed for testing whether there is no serial correlation in the error term. Thus, the null hypothesis is Assumption 1.4, while the maintained hypothesis is the other assumptions of the classical regression model (including the strict exogeneity assumption) and the normality assumption. But, as emphasized in Section 1.1, the strict exogeneity assumption is not satisfied in time-series models typically encountered in econometrics, and serial correlation is an issue that arises only in time-series models. Thus, the Durbin-Watson statistic is not useful in econometrics. More useful tests for serial correlation, which are all based on large-sample theory, will be covered in the next chapter.

\section*{QUESTIONS FOR REVIEW}
\begin{enumerate}
  \item (Conditional vs. unconditional distribution) Do we know from Assumptions 1.1-1.5 that the marginal (unconditional) distribution of $\mathbf{b}$ is normal? [Answer: No.] Are the statistics $z_{k}$ (see (1.4.3)), $t_{k}$, and $F$ distributed independently of X? [Answer: Yes, because their distributions conditional on $\mathbf{X}$ don't depend on $\mathbf{X}$.]
  \item (Computation of test statistics) Verify that $S E\left(b_{k}\right)$ as well as $\mathbf{b}, S S R, s^{2}$, and $R^{2}$ can be calculated from the following sample averages: $\mathbf{S}_{\mathbf{x x}}, \mathbf{s}_{\mathbf{x y}}, \mathbf{y}^{\prime} \mathbf{y} / n$, and $\bar{y}$.
  \item For the formula (1.4.9) for the $F$ to be well-defined, the matrix $\mathbf{R}\left(\mathbf{X}^{\prime} \mathbf{X}\right)^{-1} \mathbf{R}^{\prime}$ must be nonsingular. Prove the stronger result that the matrix is positive definite. Hint: $\mathbf{X}^{\prime} \mathbf{X}$ is positive definite. The inverse of a positive definite matrix is positive definite. Since $\mathbf{\mathbf { R }}(\# \mathbf{r} \times K)$ is of full row rank, for any nonzero \#r dimensional vector $\mathbf{z}, \mathbf{R}^{\prime} \mathbf{z} \neq \mathbf{0}$.
  \item (One-tailed $t$-test) The $t$-test described in the text is the two-tailed $t$-test because the significance $\alpha$ is equally distributed between both tails of the $t$ distribution. Suppose the alternative is one-sided and written as $\mathrm{H}_{1}: \beta_{k}>\bar{\beta}_{k}$. Consider the following modification of the decision rule of the $t$-test.
\end{enumerate}

Step 1: Same as above.\\
Step 2: Find the critical value $t_{\alpha}$ such that the area in the $t$ distribution to the right of $t_{\alpha}$ is $\alpha$. Note the difference from the two-tailed test: the left tail is ignored and the area of $\alpha$ is assigned to the upper tail only.\\
Step 3: Accept if $t_{k}<t_{\alpha}$; reject otherwise.\\
Show that the size (significance level) of this one-tailed $\boldsymbol{t}$-test is $\alpha$.\\
5. (Relation between $F(1, n-K)$ and $t(n-K))$ Look up the $t$ and $F$ distribution tables to verify that $F_{\alpha}(1, n-K)=\left(t_{\alpha / 2}(n-K)\right)^{2}$ for degrees of freedom and significance levels of your choice.\\
6. ( $t$ vs. $F$ ) "It is nonsense to test a hypothesis consisting of a large number of equality restrictions, because the $t$-test will most likely reject at least some of the restrictions." Criticize this statement.\\
7. (Variance of $s^{2}$ ) Show that, under Assumptions 1.1-1.5,

$$
\operatorname{Var}\left(s^{2} \mid \mathbf{X}\right)=\frac{2 \sigma^{4}}{n-K}
$$

Hint: If a random variable is distributed as $\chi^{2}(m)$, then its mean is $m$ and variance $2 m$.

\subsection*{1.5 Relation to Maximum Likelihood}
Having specified the distribution of the error vector $\boldsymbol{\varepsilon}$, we can use the maximum likelihood (ML) principle to estimate the model parameters ( $\boldsymbol{\beta}, \sigma^{2}$ ). ${ }^{21}$ In this section, we will show that $\mathbf{b}$, the OLS estimator of $\boldsymbol{\beta}$, is also the ML estimator, and the OLS estimator of $\sigma^{2}$ differs only slightly from the ML counterpart, when the error is normally distributed. We will also show that $\mathbf{b}$ achieves the Cramer-Rao lower bound.

\section*{The Maximum Likelihood Principle}
As you might recall from elementary statistics, the basic idea of the ML principle is to choose the parameter estimates to maximize the probability of obtaining the observed sample. To be more precise, we assume that the probability density of the sample ( $\mathbf{y}, \mathbf{X}$ ) is a member of a family of functions indexed by a finite-dimensional parameter vector $\tilde{\zeta}: f(\mathbf{y}, \mathbf{X} ; \tilde{\zeta})$. (This is described as parameterizing the density function.) This function, viewed as a function of the hypothetical parameter vector $\tilde{\zeta}$, is called the likelihood function. At the true parameter vector $\zeta$, the density of $(\mathbf{y}, \mathbf{X})$ is $f(\mathbf{y}, \mathbf{X} ; \boldsymbol{\zeta})$. The ML estimate of the true parameter vector $\boldsymbol{\zeta}$ is the $\tilde{\zeta}$ that maximizes the likelihood function given the data $(\mathbf{y}, \mathbf{X})$.

\section*{Conditional versus Unconditional Likelihood}
Since a (joint) density is the product of a marginal density and a conditional density, the density of $(\mathbf{y}, \mathbf{X})$ can be written as


\begin{equation*}
f(\mathbf{y}, \mathbf{X} ; \boldsymbol{\zeta})=f(\mathbf{y} \mid \mathbf{X} ; \boldsymbol{\theta}) \cdot f(\mathbf{X} ; \boldsymbol{\psi}) \tag{1.5.1}
\end{equation*}


where $\boldsymbol{\theta}$ is the subset of the parameter vector $\zeta$ that determines the conditional density function and $\boldsymbol{\psi}$ is the subset determining the marginal density function. The parameter vector of interest is $\boldsymbol{\theta}$; for the linear regression model with normal errors, $\boldsymbol{\theta}=\left(\boldsymbol{\beta}^{\prime}, \sigma^{2}\right)^{\prime}$ and $f(\mathbf{y} \mid \mathbf{X} ; \boldsymbol{\theta})$ is given by (1.5.4) below.

Let $\tilde{\zeta} \equiv\left(\tilde{\boldsymbol{\theta}}^{\prime}, \tilde{\boldsymbol{\psi}}^{\prime}\right)^{\prime}$ be a hypothetical value of $\zeta=\left(\boldsymbol{\theta}^{\prime}, \boldsymbol{\psi}^{\prime}\right)^{\prime}$. Then the (unconditional or joint) likelihood function is


\begin{equation*}
f(\mathbf{y}, \mathbf{X} ; \tilde{\zeta})=f(\mathbf{y} \mid \mathbf{X} ; \tilde{\boldsymbol{\theta}}) \cdot f(\mathbf{X} ; \tilde{\boldsymbol{\psi}}) \tag{1.5.2}
\end{equation*}


If we knew the parametric form of $f(\mathbf{X} ; \tilde{\boldsymbol{\psi}})$, then we could maximize this joint

\footnotetext{${ }^{21}$ For a fuller treatment of maximum likelihood, see Chapter 7.
}
likelihood function over the entire hypothetical parameter vector $\tilde{\zeta}$, and the ML estimate of $\boldsymbol{\theta}$ would be the elements of the ML estimate of $\boldsymbol{\zeta}$. We cannot do this for the classical regression model because the model does not specify $f(\mathbf{X} ; \tilde{\boldsymbol{\psi}})$. However, if there is no functional relationship between $\tilde{\boldsymbol{\theta}}$ and $\tilde{\boldsymbol{\psi}}$ (such as a subset of $\tilde{\boldsymbol{\psi}}$ being a function of $\tilde{\boldsymbol{\theta}}$ ), then maximizing (1.5.2) with respect to $\tilde{\zeta}$ is achieved by separately maximizing $f(\mathbf{y} \mid \mathbf{X} ; \tilde{\boldsymbol{\theta}})$ with respect to $\tilde{\boldsymbol{\theta}}$ and maximizing $f(\mathbf{X} ; \tilde{\boldsymbol{\psi}})$ with respect to $\tilde{\boldsymbol{\psi}}$. Thus the ML estimate of $\boldsymbol{\theta}$ also maximizes the conditional likelihood $f(\mathbf{y} \mid \mathbf{X} ; \tilde{\boldsymbol{\theta}})$.

\section*{The Log Likelihood for the Regression Model}
As already observed, Assumption 1.5 (the normality assumption) together with Assumptions 1.2 and 1.4 imply that the distribution of $\boldsymbol{\varepsilon}$ conditional on $\mathbf{X}$ is $N(\mathbf{0}$, $\sigma^{2} \mathbf{I}_{n}$ ) (see (1.4.1)). But since $\mathbf{y}=\mathbf{X} \boldsymbol{\beta}+\boldsymbol{\varepsilon}$ by Assumption 1.1, we have


\begin{equation*}
\mathbf{y} \mid \mathbf{X} \sim N\left(\mathbf{X} \boldsymbol{\beta}, \sigma^{2} \mathbf{I}_{n}\right) \tag{1.5.3}
\end{equation*}


Thus, the conditional density of $\mathbf{y}$ given $\mathbf{X}$ is ${ }^{22}$


\begin{equation*}
f(\mathbf{y} \mid \mathbf{X})=\left(2 \pi \sigma^{2}\right)^{-n / 2} \exp \left[-\frac{1}{2 \sigma^{2}}(\mathbf{y}-\mathbf{X} \boldsymbol{\beta})^{\prime}(\mathbf{y}-\mathbf{X} \boldsymbol{\beta})\right] \tag{1.5.4}
\end{equation*}


Replacing the true parameters ( $\boldsymbol{\beta}, \sigma^{2}$ ) by their hypothetical values ( $\widetilde{\boldsymbol{\beta}}, \tilde{\sigma}^{2}$ ) and taking logs, we obtain the log likelihood function:


\begin{equation*}
\log L\left(\tilde{\boldsymbol{\beta}}, \tilde{\sigma}^{2}\right)=-\frac{n}{2} \log (2 \pi)-\frac{n}{2} \log \left(\tilde{\sigma}^{2}\right)-\frac{1}{2 \tilde{\sigma}^{2}}(\mathbf{y}-\mathbf{X} \tilde{\boldsymbol{\beta}})^{\prime}(\mathbf{y}-\mathbf{X} \tilde{\boldsymbol{\beta}}) \tag{1.5.5}
\end{equation*}


Since the log transformation is a monotone transformation, the ML estimator of ( $\boldsymbol{\beta}, \sigma^{2}$ ) is the ( $\widetilde{\boldsymbol{\beta}}, \tilde{\sigma}^{2}$ ) that maximizes this log likelihood.

\section*{ML via Concentrated Likelihood}
It is instructive to maximize the log likelihood in two stages. First, maximize over $\widetilde{\boldsymbol{\beta}}$ for any given $\tilde{\sigma}^{2}$. The $\tilde{\boldsymbol{\beta}}$ that maximizes the objective function could (but does not, in the present case of Assumptions 1.1-1.5) depend on $\tilde{\sigma}^{2}$. Second, maximize over $\tilde{\sigma}^{2}$ taking into account that the $\tilde{\boldsymbol{\beta}}$ obtained in the first stage could depend on $\tilde{\sigma}^{2}$. The log likelihood function in which $\tilde{\boldsymbol{\beta}}$ is constrained to be the value from\\
${ }^{22}$ Recall from basic probability theory that the density function for an $n$-variate normal distribution with mean $\boldsymbol{\mu}$ and variance matrix $\boldsymbol{\Sigma}$ is

$$
(2 \pi)^{-n / 2}|\boldsymbol{\Sigma}|^{-1 / 2} \exp \left[-\frac{1}{2}(\mathbf{y}-\boldsymbol{\mu})^{\prime} \boldsymbol{\Sigma}^{-1}(\mathbf{y}-\boldsymbol{\mu})\right]
$$

To derive (1.5.4), just set $\mu=\mathbf{X} \boldsymbol{\beta}$ and $\boldsymbol{\Sigma}=\sigma^{2} \mathbf{I}_{n}$.\\
the first stage is called the concentrated log likelihood function (concentrated with respect to $\widetilde{\beta}$ ). For the normal log likelihood (1.5.5), the first stage amounts to minimizing the sum of squares $(\mathbf{y}-\mathbf{X} \widetilde{\boldsymbol{\beta}})^{\prime}(\mathbf{y}-\mathbf{X} \widetilde{\boldsymbol{\beta}})$. The $\widetilde{\boldsymbol{\beta}}$ that does it is none other than the OLS estimator $\mathbf{b}$, and the minimized sum of squares is $\mathbf{e}^{\prime} \mathbf{e}$. Thus, the concentrated log likelihood is


\begin{equation*}
\text { concentrated } \log \text { likelihood }=-\frac{n}{2} \log (2 \pi)-\frac{n}{2} \log \left(\tilde{\sigma}^{2}\right)-\frac{1}{2 \tilde{\sigma}^{2}} \mathbf{e}^{\prime} \mathbf{e} \tag{1.5.6}
\end{equation*}


This is a function of $\tilde{\sigma}^{2}$ alone, and the $\tilde{\sigma}^{2}$ that maximizes the concentrated likelihood is the ML estimate of $\sigma^{2}$. The maximization is straightforward for the present case of the classical regression model, because $\mathbf{e}^{\prime} \mathbf{e}$ is not a function of $\tilde{\sigma}^{2}$ and so can be taken as a constant. Still, taking the derivative with respect to $\tilde{\sigma}^{2}$, rather than with respect to $\tilde{\sigma}$, can be tricky. This can be avoided by denoting $\tilde{\sigma}^{2}$ by $\tilde{\gamma}$. Taking the derivative of (1.5.6) with respect to $\tilde{\gamma}\left(\equiv \tilde{\sigma}^{2}\right)$ and setting it to zero, we obtain the following result.

\section*{Proposition 1.5 (ML Estimator of ( $\boldsymbol{\beta}, \boldsymbol{\sigma}^{\mathbf{2}}$ )): Suppose Assumptions 1.1-1.5 hold.}
Then the ML estimator of $\boldsymbol{\beta}$ is the OLS estimator $\mathbf{b}$ and


\begin{equation*}
M L \text { estimator of } \sigma^{2}=\frac{1}{n} \mathbf{e}^{\prime} \mathbf{e}=\frac{S S R}{n}=\frac{n-K}{n} s^{2} . \tag{1.5.7}
\end{equation*}


We know from Proposition 1.2 that $s^{2}$ is unbiased. Since $s^{2}$ is multiplied by a factor $(n-K) / n$ which is different from 1 , the ML estimator of $\sigma^{2}$ is biased, although the bias becomes arbitrarily small as the sample size $n$ increases for any given fixed $K$.

For later use, we calculate the maximized value of the likelihood function. Substituting (1.5.7) into (1.5.6), we obtain

$$
\text { maximized log likelihood }=-\frac{n}{2} \log \left(\frac{2 \pi}{n}\right)-\frac{n}{2}-\frac{n}{2} \log (S S R),
$$

so that the maximized likelihood is


\begin{equation*}
\max _{\widetilde{\boldsymbol{\beta}}, \tilde{\sigma}^{2}} L\left(\widetilde{\boldsymbol{\beta}}, \tilde{\sigma}^{2}\right)=\left(\frac{2 \pi}{n}\right)^{-n / 2} \cdot \exp \left(-\frac{n}{2}\right) \cdot(S S R)^{-n / 2} \tag{1.5.8}
\end{equation*}


\section*{Cramer-Rao Bound for the Classical Regression Model}
Just to refresh your memory of basic statistics, we temporarily step outside the classical regression model and present without proof the Cramer-Rao inequality for the variance-covariance matrix of any unbiased estimator. For this purpose, define\\
the score vector at a hypothetical parameter value $\tilde{\boldsymbol{\theta}}$ to be the gradient (vector of partial derivatives) of log likelihood:


\begin{equation*}
\text { score: } \mathbf{s}(\tilde{\boldsymbol{\theta}}) \equiv \frac{\partial \log L(\tilde{\boldsymbol{\theta}})}{\partial \tilde{\boldsymbol{\theta}}} \tag{1.5.9}
\end{equation*}


Cramer-Rao Inequality: Let $\mathbf{z}$ be a vector of random variables (not necessarily independent) the joint density of which is given by $f(\mathbf{z} ; \boldsymbol{\theta})$, where $\boldsymbol{\theta}$ is an $m$ dimensional vector of parameters in some parameter space $\Theta$. Let $L(\tilde{\boldsymbol{\theta}}) \equiv f(\mathbf{z} ; \tilde{\boldsymbol{\theta}})$ be the likelihood function, and let $\hat{\boldsymbol{\theta}}(\mathbf{z})$ be an unbiased estimator of $\boldsymbol{\theta}$ with a finite variance-covariance matrix. Then, under some regularity conditions on $f(\mathbf{z} ; \boldsymbol{\theta})$ (not stated here),

$$
\operatorname{Var}[\hat{\boldsymbol{\theta}}(\mathbf{z})] \geq \underset{(m \times m)}{\mathbf{I}(\boldsymbol{\theta})^{-1}}(\equiv \text { Cramer-Rao Lower Bound })
$$

where $\mathbf{I}(\boldsymbol{\theta})$ is the information matrix defined by


\begin{equation*}
\mathbf{I}(\boldsymbol{\theta}) \equiv \mathrm{E}\left[\mathbf{s}(\boldsymbol{\theta}) \mathbf{s}(\boldsymbol{\theta})^{\prime}\right] . \tag{1.5.10}
\end{equation*}


(Note well that the score is evaluated at the true parameter value $\boldsymbol{\theta}$.) Also under the regularity conditions, the information matrix equals the negative of the expected value of the Hessian (matrix of second partial derivatives) of the log likelihood:


\begin{equation*}
\mathbf{I}(\boldsymbol{\theta})=-\mathrm{E}\left[\frac{\partial^{2} \log L(\boldsymbol{\theta})}{\partial \tilde{\boldsymbol{\theta}} \partial \tilde{\boldsymbol{\theta}}^{\prime}}\right] . \tag{1.5.11}
\end{equation*}


This is called the information matrix equality.

See, e.g., Amemiya (1985, Theorem 1.3.1) for a proof and a statement of the regularity conditions. Those conditions guarantee that the operations of differentiation and taking expectations can be interchanged. Thus, for example,

$$
\mathrm{E}[\partial L(\boldsymbol{\theta}) / \partial \tilde{\boldsymbol{\theta}}]=\partial \mathrm{E}[L(\boldsymbol{\theta})] / \partial \tilde{\boldsymbol{\theta}}
$$

Now, for the classical regression model (of Assumptions 1.1-1.5), the likelihood function $L(\tilde{\boldsymbol{\theta}})$ in the Cramer-Rao inequality is the conditional density (1.5.4), so the variance in the inequality is the variance conditional on $\mathbf{X}$. It can be shown that those regularity conditions are satisfied for the normal density (1.5.4) (see, e.g., Amemiya, 1985, Sections 1.3.2 and 1.3.3). In the rest of this subsection, we calculate the information matrix for (1.5.4). The parameter vector $\boldsymbol{\theta}$ is $\left(\boldsymbol{\beta}^{\prime}, \sigma^{2}\right)^{\prime}$.

So $\tilde{\boldsymbol{\theta}}=\left(\tilde{\boldsymbol{\beta}}^{\prime}, \tilde{\gamma}\right)^{\prime}$ and the matrix of second derivatives we seek to calculate is

\[
\frac{\partial^{2} \log L(\boldsymbol{\theta})}{\partial \tilde{\boldsymbol{\theta}} \partial \tilde{\boldsymbol{\theta}}^{\prime}}=\left[\begin{array}{cc}
\frac{\partial^{2} \log L(\boldsymbol{\theta})}{\partial \widetilde{\boldsymbol{\beta}} \partial \widetilde{\boldsymbol{\beta}}^{\prime}} & \frac{\partial^{2} \log L(\boldsymbol{\theta})}{\partial \widetilde{\boldsymbol{\beta}} \partial \tilde{\gamma}}  \tag{1.5.12}\\
((K \times K) \times(K+1)) & (K \times 1) \\
\frac{\partial^{2} \log L(\boldsymbol{\theta})}{\partial \tilde{\gamma} \partial \widetilde{\boldsymbol{\beta}}^{\prime}} & \frac{\partial^{2} \log L(\boldsymbol{\theta})}{\partial^{2} \tilde{\gamma}} \\
(1 \times K) & (1 \times 1)
\end{array}\right] .
\]

The first and second derivatives of the $\log$ likelihood (1.5.5) with respect to $\tilde{\boldsymbol{\theta}}$, evaluated at the true parameter vector $\boldsymbol{\theta}$, are


\begin{align*}
& \frac{\partial \log L(\boldsymbol{\theta})}{\partial \widetilde{\boldsymbol{\beta}}}=\frac{1}{\gamma} \mathbf{X}^{\prime}(\mathbf{y}-\mathbf{X} \boldsymbol{\beta}),  \tag{1.5.13a}\\
& \frac{\partial \log L(\boldsymbol{\theta})}{\partial \tilde{\gamma}}=-\frac{n}{2 \gamma}+\frac{1}{2 \gamma^{2}}(\mathbf{y}-\mathbf{X} \boldsymbol{\beta})^{\prime}(\mathbf{y}-\mathbf{X} \boldsymbol{\beta}) .  \tag{1.5.13b}\\
& \frac{\partial^{2} \log L(\boldsymbol{\theta})}{\partial \widetilde{\boldsymbol{\beta}} \partial \widetilde{\boldsymbol{\beta}}^{\prime}}=-\frac{1}{\gamma} \mathbf{X}^{\prime} \mathbf{X},  \tag{1.5.14a}\\
& \frac{\partial^{2} \log L(\boldsymbol{\theta})}{\partial^{2} \tilde{\gamma}}=\frac{n}{2 \gamma^{2}}-\frac{1}{\gamma^{3}}(\mathbf{y}-\mathbf{X} \boldsymbol{\beta})^{\prime}(\mathbf{y}-\mathbf{X} \boldsymbol{\beta}),  \tag{1.5.14b}\\
& \frac{\partial^{2} \log L(\boldsymbol{\theta})}{\partial \widetilde{\boldsymbol{\beta}} \partial \tilde{\gamma}}=-\frac{1}{\gamma^{2}} \mathbf{X}^{\prime}(\mathbf{y}-\mathbf{X} \boldsymbol{\beta}) . \tag{1.5.14c}
\end{align*}


Since the derivatives are evaluated at the true parameter value, $\mathbf{y}-\mathbf{X} \boldsymbol{\beta}=\boldsymbol{\varepsilon}$ in these expressions. Substituting (1.5.14) into (1.5.12) and using $\mathrm{E}(\boldsymbol{\varepsilon} \mid \mathbf{X})=\mathbf{0}$ (Assumption 1.2), $\mathrm{E}\left(\boldsymbol{\varepsilon}^{\prime} \boldsymbol{\varepsilon} \mid \mathbf{X}\right)=n \sigma^{2}$ (implication of Assumption 1.4), and recalling $\gamma=\sigma^{2}$, we can easily derive

\[
\mathbf{I}(\boldsymbol{\theta})=\left[\begin{array}{cc}
\frac{1}{\sigma^{2}} \mathbf{X}^{\prime} \mathbf{X} & \mathbf{0}  \tag{1.5.15}\\
\mathbf{0}^{\prime} & \frac{n}{2 \sigma^{4}}
\end{array}\right] .
\]

Here, the expectation is conditional on $\mathbf{X}$ because the likelihood function (1.5.4) is a conditional density conditional on $\mathbf{X}$. This block diagonal matrix can be inverted to obtain the Cramer-Rao bound:

\[
\text { Cramer-Rao bound } \equiv \mathbf{I}(\boldsymbol{\theta})^{-1}=\left[\begin{array}{cc}
\sigma^{2} \cdot\left(\mathbf{X}^{\prime} \mathbf{X}\right)^{-1} & \mathbf{0}  \tag{1.5.16}\\
\mathbf{0}^{\prime} & \frac{2 \sigma^{4}}{n}
\end{array}\right] \text {. }
\]

Therefore, the unbiased estimator $\mathbf{b}$, whose variance is $\sigma^{2} \cdot\left(\mathbf{X}^{\prime} \mathbf{X}\right)^{-1}$ by Proposition 1.1, attains the Cramer-Rao bound. We have thus proved

Proposition 1.6 (b is the Best Unbiased Estimator (BUE)): Under Assumptions 1.1-1.5, the OLS estimator $\mathbf{b}$ of $\boldsymbol{\beta}$ is BUE in that any other unbiased (but not necessarily linear) estimator has larger conditional variance in the matrix sense.

This result should be distinguished from the Gauss-Markov Theorem that $\mathbf{b}$ is minimum variance among those estimators that are unbiased and linear in $\mathbf{y}$. Proposition 1.6 says that $\mathbf{b}$ is minimum variance in a larger class of estimators that includes nonlinear unbiased estimators. This stronger statement is obtained under the normality assumption (Assumption 1.5) which is not assumed in the Gauss-Markov Theorem. Put differently, the Gauss-Markov Theorem does not exclude the possibility of some nonlinear estimator beating OLS, but this possibility is ruled out by the normality assumption.

As was already seen, the ML estimator of $\sigma^{2}$ is biased, so the Cramer-Rao bound does not apply. But the OLS estimator $s^{2}$ of $\sigma^{2}$ is unbiased. Does it achieve the bound? We have shown in a review question to the previous section that

$$
\operatorname{Var}\left(s^{2} \mid \mathbf{X}\right)=\frac{2 \sigma^{4}}{n-K}
$$

under the same set of assumptions as in Proposition 1.6. Therefore, $s^{2}$ does not attain the Cramer-Rao bound $2 \sigma^{4} / n$. However, it can be shown that an unbiased estimator of $\sigma^{2}$ with variance lower than $2 \sigma^{4} /(n-K)$ does not exist (see, e.g., Rao, 1973, p. 319).

\section*{The $\boldsymbol{F}$-Test as a Likelihood Ratio Test}
The likelihood ratio test of the null hypothesis compares $L_{U}$, the maximized likelihood without the imposition of the restriction specified in the null hypothesis, with $L_{R}$, the likelihood maximized subject to the restriction. If the likelihood ratio $\lambda \equiv L_{U} / L_{R}$ is too large, it should be a sign that the null is false. The F-test of the null hypothesis $\mathrm{H}_{0}: \mathbf{R} \boldsymbol{\beta}=\mathbf{r}$ considered in the previous section is a likelihood ratio test because the $F$-ratio is a monotone transformation of the likelihood ratio $\lambda$. For the present model, $L_{U}$ is given by (1.5.8) where the SSR, the sum of squared residuals minimized without the constraint $\mathrm{H}_{0}$, is the $S S R_{U}$ in (1.4.11). The restricted likelihood $L_{R}$ is given by replacing this SSR by the restricted sum of squared residuals, $S S R_{R}$. So


\begin{equation*}
L_{R}=\max _{\widetilde{\boldsymbol{\beta}}, \tilde{\sigma}^{2} \text { s.t. } \mathrm{H}_{0}} L\left(\widetilde{\boldsymbol{\beta}}, \tilde{\sigma}^{2}\right)=\left(\frac{2 \pi}{n}\right)^{-n / 2} \cdot \exp \left(-\frac{n}{2}\right) \cdot\left(S S R_{R}\right)^{-n / 2}, \tag{1.5.17}
\end{equation*}


and the likelihood ratio is


\begin{equation*}
\lambda \equiv \frac{L_{U}}{L_{R}}=\left(\frac{S S R_{U}}{S S R_{R}}\right)^{-n / 2} \tag{1.5.18}
\end{equation*}


Comparing this with the formula (1.4.11) for the $F$-ratio, we see that the $F$-ratio is a monotone transformation of the likelihood ratio $\lambda$ :


\begin{equation*}
F=\frac{n-K}{\# \mathbf{r}}\left(\lambda^{2 / n}-1\right), \tag{1.5.19}
\end{equation*}


so that the two tests are the same.

\section*{Quasi-Maximum Likelihood}
All these results assume the normality of the error term. Without normality, there is no guarantee that the ML estimator of $\boldsymbol{\beta}$ is OLS (Proposition 1.5) or that the OLS estimator $\mathbf{b}$ achieves the Cramer-Rao bound (Proposition 1.6). However, Proposition 1.5 does imply that $\mathbf{b}$ is a quasi- (or pseudo-) maximum likelihood estimator, an estimator that maximizes a misspecified likelihood function. The misspecified likelihood function we have considered is the normal likelihood. The results of Section 1.3 can then be interpreted as providing the finite-sample properties of the quasi-ML estimator when the error is incorrectly specified to be normal.

\section*{QUESTIONS FOR REVIEW}
\begin{enumerate}
  \item (Use of regularity conditions) Assuming that taking expectations (i.e., taking integrals) and differentiation can be interchanged, prove that the expected value of the score vector given in (1.5.9), if evaluated at the true parameter value $\boldsymbol{\theta}$, is zero. Hint: What needs to be shown is that
\end{enumerate}

$$
\int \frac{\partial \log f(\mathbf{z} ; \boldsymbol{\theta})}{\partial \tilde{\boldsymbol{\theta}}} f(\mathbf{z} ; \boldsymbol{\theta}) \mathrm{d} \mathbf{z}=\mathbf{0} .
$$

Since $f(\mathbf{z} ; \tilde{\boldsymbol{\theta}})$ is a density, $\int f(\mathbf{z}, \tilde{\boldsymbol{\theta}}) \mathrm{d} \mathbf{z}=1$ for any $\tilde{\boldsymbol{\theta}}$. Differentiate both sides with respect to $\tilde{\boldsymbol{\theta}}$ and use the regularity conditions, which allows us to change the order of integration and differentiation, to obtain $\int[\partial f(\mathbf{z} ; \boldsymbol{\theta}) / \partial \tilde{\boldsymbol{\theta}}] \mathrm{d} \mathbf{z}=\mathbf{0}$. Also, from basic calculus,

$$
\frac{\partial \log f(\mathbf{z} ; \boldsymbol{\theta})}{\partial \tilde{\boldsymbol{\theta}}}=\frac{1}{f(\mathbf{z} ; \boldsymbol{\theta})} \frac{\partial f(\mathbf{z} ; \boldsymbol{\theta})}{\partial \tilde{\boldsymbol{\theta}}} .
$$

\begin{enumerate}
  \setcounter{enumi}{1}
  \item (Maximizing joint log likelihood) Consider maximizing (the log of) the joint likelihood (1.5.2) for the classical regression model, where $\tilde{\boldsymbol{\theta}}=\left(\tilde{\boldsymbol{\beta}}^{\prime}, \tilde{\sigma}^{2}\right)^{\prime}$ and $\log f(\mathbf{y} \mid \mathbf{X} ; \tilde{\boldsymbol{\theta}})$ is given by (1.5.5). You would parameterize the marginal like-\\
lihood $f(\mathbf{X} ; \tilde{\boldsymbol{\psi}})$ and take the $\log$ of (1.5.2) to obtain the objective function to be maximized over $\boldsymbol{\zeta} \equiv\left(\boldsymbol{\theta}^{\prime}, \boldsymbol{\psi}^{\prime}\right)^{\prime}$. What is the ML estimator of $\boldsymbol{\theta} \equiv\left(\boldsymbol{\beta}^{\prime}, \sigma^{2}\right)^{\prime}$ ? [Answer: It should be the same as that in Proposition 1.5.] Derive the CramerRao bound for $\boldsymbol{\beta}$. Hint: By the information matrix equality,
\end{enumerate}

$$
\mathbf{I}(\zeta)=-\mathrm{E}\left[\frac{\partial^{2} \log L(\zeta)}{\partial \tilde{\zeta} \partial \tilde{\zeta}^{\prime}}\right]
$$

Also, $\partial^{2} \log L(\zeta) /\left(\partial \tilde{\boldsymbol{\theta}} \partial \tilde{\boldsymbol{\psi}}^{\prime}\right)=\mathbf{0}$.\\
3. (Concentrated $\log$ likelihood with respect to $\tilde{\sigma}^{2}$ ) Writing $\tilde{\sigma}^{2}$ as $\tilde{\gamma}$, the $\log$ likelihood function for the classical regression model is

$$
\log L(\tilde{\boldsymbol{\beta}}, \tilde{\gamma})=-\frac{n}{2} \log (2 \pi)-\frac{n}{2} \log (\tilde{\gamma})-\frac{1}{2 \tilde{\gamma}}(\mathbf{y}-\mathbf{X} \tilde{\boldsymbol{\beta}})^{\prime}(\mathbf{y}-\mathbf{X} \tilde{\boldsymbol{\beta}}) .
$$

In the two-step maximization procedure described in the text, we first maximized this function with respect to $\widetilde{\boldsymbol{\beta}}$. Instead, first maximize with respect to $\tilde{\gamma}$ given $\widetilde{\boldsymbol{\beta}}$. Show that the concentrated log likelihood (concentrated with respect to $\tilde{\gamma} \equiv \tilde{\sigma}^{2}$ ) is

$$
-\frac{n}{2}[1+\log (2 \pi)]-\frac{n}{2} \log \left(\frac{(\mathbf{y}-\mathbf{X} \tilde{\boldsymbol{\beta}})^{\prime}(\mathbf{y}-\mathbf{X} \tilde{\boldsymbol{\beta}})}{n}\right)
$$

\begin{enumerate}
  \setcounter{enumi}{3}
  \item (Information matrix equality for classical regression model) Verify (1.5.11) for the linear regression model.
  \item (Likelihood equations for classical regression model) We used the two-step procedure to derive the ML estimate for the classical regression model. An alternative way to find the ML estimator is to solve for the first-order conditions that set (1.5.13) equal to zero (the first-order conditions for the log likelihood is called the likelihood equations). Verify that the ML estimator given in Proposition 1.5 solves the likelihood equations.
\end{enumerate}

\subsection*{1.6 Generalized Least Squares (GLS)}
Assumption 1.4 states that the $n \times n$ matrix of conditional second moments $\mathrm{E}\left(\boldsymbol{\varepsilon} \boldsymbol{\varepsilon}^{\prime} \mid\right. \mathbf{X})(=\operatorname{Var}(\boldsymbol{\varepsilon} \mid \mathbf{X}))$ is spherical, that is, proportional to the identity matrix. Without the assumption, each element of the $n \times n$ matrix is in general a nonlinear function\\
of $\mathbf{X}$. If the error is not (conditionally) homoskedastic, the values of the diagonal elements of $\mathrm{E}\left(\boldsymbol{\varepsilon} \boldsymbol{\varepsilon}^{\prime} \mid \mathbf{X}\right)$ are not the same, and if there is correlation in the error term between observations (the case of serial correlation for time-series models), the values of the off-diagonal elements are not zero. For any given positive scalar $\sigma^{2}$, define $\mathbf{V}(\mathbf{X}) \equiv \mathrm{E}\left(\boldsymbol{\varepsilon} \boldsymbol{\varepsilon}^{\prime} \mid \mathbf{X}\right) / \sigma^{2}$ and assume $\mathbf{V}(\mathbf{X})$ is nonsingular and known. That is,


\begin{equation*}
\mathrm{E}\left(\boldsymbol{\varepsilon} \boldsymbol{\varepsilon}^{\prime} \mid \mathbf{X}\right)=\sigma^{2} \underset{(n \times n)}{\mathbf{V}}(\mathbf{X}), \quad \mathbf{V}(\mathbf{X}) \text { nonsingular and known. } \tag{1.6.1}
\end{equation*}


The reason we decompose $\mathrm{E}\left(\boldsymbol{\varepsilon} \boldsymbol{\varepsilon}^{\prime} \mid \mathbf{X}\right)$ into the component $\sigma^{2}$ that is common to all elements of the matrix $\mathrm{E}\left(\boldsymbol{\varepsilon} \boldsymbol{\varepsilon}^{\prime} \mid \mathbf{X}\right)$ and the remaining component $\mathbf{V}(\mathbf{X})$ is that we do not need to know the value of $\sigma^{2}$ for efficient estimation. The model that results when Assumption 1.4 is replaced by (1.6.1), which merely assumes that the conditional second moment $\mathrm{E}\left(\boldsymbol{\varepsilon} \boldsymbol{\varepsilon}^{\prime} \mid \mathbf{X}\right)$ is nonsingular, is called the generalized regression model.

\section*{Consequence of Relaxing Assumption 1.4}
Of the results derived in the previous sections, those that assume Assumption 1.4 are no longer valid for the generalized regression model. More specifically,

\begin{itemize}
  \item The Gauss-Markov Theorem no longer holds for the OLS estimator
\end{itemize}

$$
\mathbf{b} \equiv\left(\mathbf{X}^{\prime} \mathbf{X}\right)^{-1} \mathbf{X}^{\prime} \mathbf{y}
$$

The BLUE is some other estimator.

\begin{itemize}
  \item The $t$-ratio is not distributed as the $t$ distribution. Thus, the $t$-test is no longer valid. The same comments apply to the $F$-test.
  \item However, the OLS estimator is still unbiased, because the unbiasedness result (Proposition 1.1(a)) does not require Assumption 1.4.
\end{itemize}

\section*{Efficient Estimation with Known V}
If the value of the matrix function $\mathbf{V}(\mathbf{X})$ is known, does there exist a BLUE for the generalized regression model? The answer is yes, and the estimator is called the generalized least squares (GLS) estimator, which we now derive. The basic idea of the derivation is to transform the generalized regression model, which consists of Assumptions 1.1-1.3 and (1.6.1), into a model that satisfies all the assumptions, including Assumption 1.4, of the classical regression model.

For economy of notation, we use $\mathbf{V}$ for the value $\mathbf{V}(\mathbf{X})$. Since $\mathbf{V}$ is by construction symmetric and positive definite, there exists a nonsingular $n \times n$ matrix $\mathbf{C}$ such that


\begin{equation*}
\mathbf{V}^{-1}=\mathbf{C}^{\prime} \mathbf{C} \tag{1.6.2}
\end{equation*}


This decomposition is not unique, with more than one choice for $\mathbf{C}$, but, as is clear from the discussion below, the choice of $\mathbf{C}$ doesn't matter. Now consider creating a new regression model by transforming ( $\mathbf{y}, \mathbf{X}, \boldsymbol{\varepsilon}$ ) by $\mathbf{C}$ as


\begin{equation*}
\tilde{\mathbf{y}} \equiv \mathbf{C} \mathbf{y}, \quad \tilde{\mathbf{X}} \equiv \mathbf{C X}, \quad \tilde{\varepsilon} \equiv \mathbf{C} \varepsilon \tag{1.6.3}
\end{equation*}


Then Assumption 1.1 for $(\mathbf{y}, \mathbf{X}, \boldsymbol{\varepsilon})$ implies that $(\tilde{\mathbf{y}}, \tilde{\mathbf{X}}, \tilde{\boldsymbol{\varepsilon}})$ too satisfies linearity:


\begin{equation*}
\tilde{\mathbf{y}}=\widetilde{\mathbf{X}} \boldsymbol{\beta}+\tilde{\boldsymbol{\varepsilon}} \tag{1.6.4}
\end{equation*}


The transformed model satisfies the other assumptions of the classical linear regression model. Strict exogeneity is satisfied because

$$
\begin{aligned}
\mathrm{E}(\tilde{\boldsymbol{\varepsilon}} \mid \widetilde{\mathbf{X}})= & \mathrm{E}(\tilde{\boldsymbol{\varepsilon}} \mid \mathbf{X}) \\
& \quad(\text { since } \mathbf{C} \text { is nonsingular, } \mathbf{X} \text { and } \tilde{\mathbf{X}} \text { contain the same information }) \\
= & \mathrm{E}(\mathbf{C} \boldsymbol{\varepsilon} \mid \mathbf{X}) \\
= & \mathbf{C} \mathrm{E}(\boldsymbol{\varepsilon} \mid \mathbf{X}) \quad \text { (by the linearity of conditional expectations) } \\
= & \mathbf{0} \quad(\text { since } \mathrm{E}(\boldsymbol{\varepsilon} \mid \mathbf{X})=\mathbf{0} \text { by Assumption 1.2) }
\end{aligned}
$$

Because $\mathbf{V}$ is positive definite, the no-multicollinearity assumption is also satisfied (see a review question below for a proof). Assumption 1.4 is satisfied for the transformed model because

$$
\begin{aligned}
\mathrm{E}\left(\tilde{\boldsymbol{\varepsilon}} \tilde{\boldsymbol{\varepsilon}}^{\prime} \mid \widetilde{\mathbf{X}}\right) & =\mathrm{E}\left(\tilde{\boldsymbol{\varepsilon}} \tilde{\boldsymbol{\varepsilon}}^{\prime} \mid \mathbf{X}\right) \quad(\text { since } \tilde{\mathbf{X}} \text { and } \mathbf{X} \text { contain the same information }) \\
& =\mathbf{C} \mathrm{E}\left(\boldsymbol{\varepsilon} \boldsymbol{\varepsilon}^{\prime} \mid \mathbf{X}\right) \mathbf{C}^{\prime} \quad\left(\text { since } \tilde{\boldsymbol{\varepsilon}} \tilde{\boldsymbol{\varepsilon}}^{\prime}=\mathbf{C} \boldsymbol{\varepsilon} \boldsymbol{\varepsilon}^{\prime} \mathbf{C}^{\prime}\right) \\
& =\mathbf{C} \cdot \sigma^{2} \cdot \mathbf{V} \mathbf{C}^{\prime} \quad(\text { by }(1.6 .1)) \\
& =\sigma^{2} \mathbf{C} \mathbf{V} \mathbf{C}^{\prime} \\
& =\sigma^{2} \mathbf{I}_{n} \quad\left(\text { since }\left(\mathbf{C}^{\prime}\right)^{-1} \mathbf{V}^{-1} \mathbf{C}^{-1}=\mathbf{I}_{n} \text { or } \mathbf{C} \mathbf{V C ^ { \prime }}=\mathbf{I}_{n} \text { by }(1.6 .2)\right) .
\end{aligned}
$$

So indeed the variance of the transformed error vector $\tilde{\boldsymbol{\varepsilon}}$ is spherical. Finally, $\tilde{\boldsymbol{\varepsilon}} \mid \widetilde{\mathbf{X}}$ is normal because the distribution of $\tilde{\boldsymbol{\varepsilon}} \mid \tilde{\mathbf{X}}$ is the same as $\tilde{\boldsymbol{\varepsilon}} \mid \mathbf{X}$ and $\tilde{\boldsymbol{\varepsilon}}$ is a linear transformation of $\boldsymbol{\varepsilon}$. This completes the verification of Assumptions 1.1-1.5 for the transformed model.

The Gauss-Markov Theorem for the transformed model implies that the BLUE of $\boldsymbol{\beta}$ for the generalized regression model is the OLS estimator applied to (1.6.4):


\begin{align*}
\widehat{\boldsymbol{\beta}}_{\mathrm{GLS}} & =\left(\tilde{\mathbf{X}}^{\prime} \tilde{\mathbf{X}}\right)^{-1} \tilde{\mathbf{X}}^{\prime} \tilde{\mathbf{y}} \\
& =\left[(\mathbf{C X})^{\prime}(\mathbf{C X})\right]^{-1}(\mathbf{C X})^{\prime} \mathbf{C y} \\
& =\left(\mathbf{X}^{\prime} \mathbf{C}^{\prime} \mathbf{C X}\right)^{-1}\left(\mathbf{X}^{\prime} \mathbf{C}^{\prime} \mathbf{C y}\right) \\
& =\left(\mathbf{X}^{\prime} \mathbf{V}^{-1} \mathbf{X}\right)^{-1} \mathbf{X}^{\prime} \mathbf{V}^{-1} \mathbf{y} \quad(\text { by }(1.6 .2)) . \tag{1.6.5}
\end{align*}


This is the GLS estimator. Its conditional variance is


\begin{align*}
& \operatorname{Var}\left(\widehat{\boldsymbol{\beta}}_{\mathrm{GLS}} \mid \mathbf{X}\right) \\
& =\left(\mathbf{X}^{\prime} \mathbf{V}^{-1} \mathbf{X}\right)^{-1} \mathbf{X}^{\prime} \mathbf{V}^{-1} \operatorname{Var}(\mathbf{y} \mid \mathbf{X}) \mathbf{V}^{-1} \mathbf{X}\left(\mathbf{X}^{\prime} \mathbf{V}^{-1} \mathbf{X}\right)^{-1} \\
& =\left(\mathbf{X}^{\prime} \mathbf{V}^{-1} \mathbf{X}\right)^{-1} \mathbf{X}^{\prime} \mathbf{V}^{-1}\left(\sigma^{2} \mathbf{V}\right) \mathbf{V}^{-1} \mathbf{X}\left(\mathbf{X}^{\prime} \mathbf{V}^{-1} \mathbf{X}\right)^{-1} \quad(\text { since } \operatorname{Var}(\mathbf{y} \mid \mathbf{X})=\operatorname{Var}(\boldsymbol{\varepsilon} \mid \mathbf{X})) \\
& =\sigma^{2} \cdot\left(\mathbf{X}^{\prime} \mathbf{V}^{-1} \mathbf{X}\right)^{-1} \tag{1.6.6}
\end{align*}


Since replacing $\mathbf{V}$ by $\sigma^{2} \cdot \mathbf{V}(=\operatorname{Var}(\boldsymbol{\varepsilon} \mid \mathbf{X}))$ in (1.6.5) does not change the numerical value, the GLS estimator can also be written as

$$
\widehat{\boldsymbol{\beta}}_{\mathrm{GLS}}=\left[\mathbf{X}^{\prime} \operatorname{Var}(\boldsymbol{\varepsilon} \mid \mathbf{X})^{-1} \mathbf{X}\right]^{-1} \mathbf{X}^{\prime} \operatorname{Var}(\boldsymbol{\varepsilon} \mid \mathbf{X})^{-1} \mathbf{y} .
$$

As noted above, the OLS estimator $\left(\mathbf{X}^{\prime} \mathbf{X}\right)^{-1} \mathbf{X}^{\prime} \mathbf{y}$ too is unbiased without Assumption 1.4, but nevertheless the GLS estimator should be preferred (provided $\mathbf{V}$ is known) because the latter is more efficient in that the variance is smaller in the matrix sense. The gain in efficiency is achieved by exploiting the heteroskedasticity and correlation between observations in the error term, which, operationally, is to insert the inverse of (a matrix proportional to) $\operatorname{Var}(\boldsymbol{\varepsilon} \mid \mathbf{X})$ in the OLS formula, as in (1.6.5). The discussion so far can be summarized as

\section*{Proposition 1.7 (finite-sample properties of GLS):}
(a) (unbiasedness) Under Assumption 1.1-1.3, $\mathrm{E}\left(\widehat{\boldsymbol{\beta}}_{\mathrm{GLS}} \mid \mathbf{X}\right)=\boldsymbol{\beta}$.\\
(b) (expression for the variance) Under Assumptions 1.1-1.3 and the assumption (1.6.1) that the conditional second moment is proportional to $\mathbf{V}(\mathbf{X})$,

$$
\operatorname{Var}\left(\widehat{\boldsymbol{\beta}}_{\mathrm{GLS}} \mid \mathbf{X}\right)=\sigma^{2} \cdot\left(\mathbf{X}^{\prime} \mathbf{V}(\mathbf{X})^{-1} \mathbf{X}\right)^{-1} .
$$

(c) (efficiency of GLS) Under the same set of assumptions as in (b), the GLS estimator is efficient in that the conditional variance of any unbiased estimator that is linear in $\mathbf{y}$ is greater than or equal to $\operatorname{Var}\left(\widehat{\boldsymbol{\beta}}_{\mathrm{GLS}} \mid \mathbf{X}\right)$ in the matrix sense.

\section*{A Special Case: Weighted Least Squares (WLS)}
The idea of adjusting for the error variance matrix becomes more transparent when there is no correlation in the error term between observations so that the matrix $\mathbf{V}$ is diagonal. Let $v_{i}(\mathbf{X})$ be the $i$-th diagonal element of $\mathbf{V}(\mathbf{X})$. So

$$
\mathrm{E}\left(\varepsilon_{i}^{2} \mid \mathbf{X}\right)\left(=\operatorname{Var}\left(\varepsilon_{i} \mid \mathbf{X}\right)\right)=\sigma^{2} \cdot v_{i}(\mathbf{X}) .
$$

It is easy to see that $\mathbf{C}$ is also diagonal, with the square root of $1 / v_{i}(\mathbf{X})$ in the $i$-th diagonal. Thus, ( $\tilde{\mathbf{y}}, \tilde{\mathbf{X}}$ ) is given by

$$
\tilde{y}_{i}=\frac{y_{i}}{\sqrt{v_{i}(\mathbf{X})}}, \quad \tilde{\mathbf{x}}_{i}=\frac{\mathbf{x}_{i}}{\sqrt{v_{i}(\mathbf{X})}} \quad(i=1,2, \ldots, n)
$$

Therefore, efficient estimation under a known form of heteroskedasticity is first to weight each observation by the reciprocal of the square root of the variance $v_{i}(\mathbf{X})$ and then apply OLS. This is called the weighted regression (or the weighted least squares (WLS)).

An important further special case is the case of a random sample where $\left\{y_{i}, \mathbf{x}_{i}\right\}$ is i.i.d. across $i$. As was noted in Section 1.1, the error is unconditionally homoskedastic (i.e., $\mathrm{E}\left(\varepsilon_{i}^{2}\right)$ does not depend on $i$ ), but still GLS can be used to increase efficiency because the error can be conditionally heteroskedastic. The conditional second moment $\mathrm{E}\left(\varepsilon_{i}^{2} \mid \mathbf{X}\right)$ for the case of random samples depends only on $\mathbf{x}_{i}$, and the functional form of $\mathrm{E}\left(\varepsilon_{i}^{2} \mid \mathbf{x}_{i}\right)$ is the same across $i$. Thus


\begin{equation*}
v_{i}(\mathbf{X})=v\left(\mathbf{x}_{i}\right) \quad \text { for random samples. } \tag{1.6.7}
\end{equation*}


So the knowledge of $\mathbf{V}(\cdot)$ comes down to a single function of $K$ variables, $v(\cdot)$.

\section*{Limiting Nature of GLS}
All these sanguine conclusions about the finite-sample properties of GLS rest on the assumption that the regressors in the generalized regression model are strictly exogenous $(\mathrm{E}(\tilde{\boldsymbol{\varepsilon}} \mid \tilde{\mathbf{X}})=\mathbf{0})$. This fact limits the usefulness of the GLS procedure. Suppose, as is often the case with time-series models, that the regressors are not strictly exogenous and the error is serially correlated. So neither OLS nor GLS has those good finite-sample properties such as unbiasedness. Nevertheless, as will be shown in the next chapter, the OLS estimator, which ignores serial correlation in the error, will have some good large sample properties (such as "consistency" and "asymptotic normality"), provided that the regressors are "predetermined" (which is weaker than strict exogeneity). The GLS estimator, in contrast, does not have that redeeming feature. That is, if the error is not strictly exogenous\\
but is merely predetermined, the GLS procedure to correct for serial correlation can make the estimator inconsistent (see Section 6.7). A procedure for explicitly taking serial correlation into account while maintaining consistency will be presented in Chapter 6.

If it is not appropriate for correcting for serial correlation, the GLS procedure can still be used to correct for heteroskedasticity when the error is not serially correlated with diagonal $\mathbf{V}(\mathbf{X})$, in the form of WLS. But that is provided that the matrix function $\mathbf{V}(\mathbf{X})$ is known. Very rarely do we have a priori information specifying the values of the diagonal elements of $\mathbf{V}(\mathbf{X})$, which is necessary to weight observations. In the case of a random sample where serial correlation is guaranteed not to arise, the knowledge of $\mathbf{V}(\mathbf{X})$ boils down to a single function of $K$ variables, $v\left(\mathbf{x}_{i}\right)$, as we have just seen, but even for this case the knowledge of such a function is unavailable in most applications.

If we do not know the function $\mathbf{V}(\mathbf{X})$, we can estimate its functional form from the sample. This approach is called the Feasible Generalized Least Squares (FGLS). But if the function $\mathbf{V}(\mathbf{X})$ is estimated from the sample, its value $\mathbf{V}$ becomes a random variable, which affects the distribution of the GLS estimator. Very little is known about the finite-sample properties of the FGLS estimator. We will cover the large-sample properties of the FGLS estimator in the context of heteroskedasticity correction in the next chapter.

Before closing, one positive side of GLS should be noted: most linear estimation techniques - including the 2SLS, 3SLS, and the random effects estimators to be introduced later - can be expressed as a GLS estimator, with some liberal definition of data matrices. However, those estimators and OLS can also be interpreted as a GMM (generalized method of moments) estimator, and the GMM interpretation is more useful for developing large-sample results.

\section*{QUESTIONS FOR REVIEW}
\begin{enumerate}
  \item (The no-multicollinearity assumption for the transformed model) Assumption 1.3 for the transformed model is that $\operatorname{rank}(\mathbf{C X})=K$. This is satisfied since $\mathbf{C}$ is nonsingular and $\mathbf{X}$ is of full column rank. Show this. Hint: Since $\mathbf{X}$ is of full column rank, for any $K$-dimensional vector $\mathbf{c} \neq \mathbf{0}, \mathbf{X c} \neq \mathbf{0}$.
  \item (Generalized $S S R$ ) Show that $\widehat{\boldsymbol{\beta}}_{\mathrm{GLS}}$ minimizes $(\mathbf{y}-\mathbf{X} \widetilde{\boldsymbol{\beta}})^{\prime} \mathbf{V}^{-1}(\mathbf{y}-\mathbf{X} \widetilde{\boldsymbol{\beta}})$.
  \item Derive the expression for $\operatorname{Var}(\mathbf{b} \mid \mathbf{X})$ for the generalized regression model. What is the relation of it to $\operatorname{Var}\left(\widehat{\boldsymbol{\beta}}_{\text {GLS }} \mid \mathbf{X}\right)$ ? Verify that Proposition 1.7(c) implies
\end{enumerate}

$$
\left(\mathbf{X}^{\prime} \mathbf{X}\right)^{-1} \mathbf{X}^{\prime} \mathbf{V} \mathbf{X}\left(\mathbf{X}^{\prime} \mathbf{X}\right)^{-1} \geq\left(\mathbf{X}^{\prime} \mathbf{V}^{-1} \mathbf{X}\right)^{-1}
$$

\begin{enumerate}
  \setcounter{enumi}{3}
  \item (Sampling error of GLS) Show: $\widehat{\boldsymbol{\beta}}_{\mathrm{GLS}}-\boldsymbol{\beta}=\left(\mathbf{X}^{\prime} \mathbf{V}^{-1} \mathbf{X}\right)^{-1} \mathbf{X}^{\prime} \mathbf{V}^{-1} \boldsymbol{\varepsilon}$.
\end{enumerate}

\subsection*{1.7 Application: Returns to Scale in Electricity Supply}
Nerlove's 1963 paper is a classic study of returns to scale in a regulated industry. It also is excellent material for illustrating the techniques of this chapter and presenting a few more not yet covered.

\section*{The Electricity Supply Industry}
At the time of Nerlove's writing, the U.S. electric power supply industry had the following features:\\
(1) Privately owned local monopolies supply power on demand.\\
(2) Rates (electricity prices) are set by the utility commission.\\
(3) Factor prices (e.g., the wage rate) are given to the firm, either because of perfect competition in the market for factor inputs or through long-term contracts with labor unions.

These institutional features will be relevant when we examine whether the OLS is an appropriate estimation procedure. ${ }^{23}$

\section*{The Data}
Nerlove assembled a cross-section data set on 145 firms in 44 states in the year 1955 for which data on all the relevant variables were available. The variables in the data are total costs, factor prices (the wage rate, the price of fuel, and the rental price of capital), and output. Although firms own capital (such as power plants, equipment, and structures), the standard investment theory of Jorgenson (1963) tells us that (as long as there are no costs in changing the capital stock) the firm should behave as if it rents capital on a period-to-period basis from itself at a rental price called the "user cost of capital," which is defined as $(r+\delta) \cdot p_{I}$, where $r$ here is the real interest rate (below we will use $r$ for the degree of returns to scale), $\delta$ is the depreciation rate, and $p_{I}$ is the price of capital goods. For this reason capital

\footnotetext{${ }^{23}$ Thanks to the deregulation of the industry since the time of Nerlove's writing, multiple firms are now allowed to compete in the same local market, and the strict price control has been lifted in many states. So the first two features no longer characterize the industry.
}
input can be treated as if it is a variable factor of production, just like labor and fuel inputs.

Appendix B of Nerlove (1963) contains a careful and honest discussion of how the data were constructed. Data on output, fuel, and labor costs (which, along with capital costs, make up total costs) were obtained from the Federal Power Commission (1956). For the wage rate, Nerlove used statewide average wages for utility workers. Ideally, one would calculate capital costs as the reproduction cost of capital times the user cost of capital. Due to data limitation, Nerlove instead used interest and depreciation charges available from the firm’s books.

\section*{Why Do We Need Econometrics?}
Why do we need a fancy econometric technique like OLS to determine returns to scale? Why can't we be simple-minded and plot the average cost (which can be easily calculated from the data as the ratio of total costs to output) against output and see whether the AC (average cost) curve is downward sloping? The reason is that each firm can have a different AC curve. If firms face different factor prices, then the average cost is less for firms facing lower factor prices. That cross-section units at a given moment face the same prices is usually a good assumption to make, but not for the U.S. electricity industry with substantial regional differences in factor prices. The effect of factor prices on the AC curve has to be isolated somehow. The approach taken by Nerlove, which became a standard econometric practice, is to estimate a parameterized cost function.

Another factor that shifts the individual AC curve is the level of production efficiency. If more efficient firms produce more output, then it is possible that the individual AC curve is upward sloping but the line connecting the observed combination of the average cost and output is downward sloping. To illustrate, consider a competitive industry described in Figure 1.6, where the AC and MC (marginal cost) curves are drawn for two firms competing in the same market. To focus on the connection between production efficiency and output, assume that all firms face the same factor prices so that the only reason the AC and MC curves differ between firms is the difference in production efficiency. The AC and MC curves are upward sloping to reflect decreasing returns to scale. The AC and MC curves for firm A lie above those for firm B because firm A is less efficient than B. Because the industry is competitive, both firms face the same price $p$. Since output is determined at the intersection of the MC curve and the market price, the combinations of output and the average cost for two firms are points A and B in the figure. The curve obtained from connecting these two points can be downward sloping, giving a false impression of increasing returns to scale.

\begin{figure}[h]
\begin{center}
  \includegraphics[alt={},max width=\textwidth]{14fb0795-6d81-4e76-9788-792298f1ddae-079_569_1159_373_468}
\captionsetup{labelformat=empty}
\caption{Figure 1.6: Output Determination}
\end{center}
\end{figure}

\section*{The Cobb-Douglas Technology}
To derive a parameterized cost function, we start with the Cobb-Douglas production function


\begin{equation*}
Q_{i}=A_{i} x_{i 1}^{\alpha_{1}} x_{i 2}^{\alpha_{2}} x_{i 3}^{\alpha_{3}}, \tag{1.7.1}
\end{equation*}


where $Q_{i}$ is firm $i$ 's output, $x_{i 1}$ is labor input for firm $i, x_{i 2}$ is capital input, and $x_{i 3}$ is fuel. $A_{i}$ captures unobservable differences in production efficiency (this term is often called firm heterogeneity). The sum $\alpha_{1}+\alpha_{2}+\alpha_{3} \equiv r$ is the degree of returns to scale. Thus, it is assumed a priori that the degree of returns to scale is constant (this should not be confused with constant returns to scale, which is that $r=1$ ). Since the electric utilities in the sample are privately owned, it is reasonable to suppose that they are engaged in cost minimization (see, however, the discussion at the end of this section). We know from microeconomics that the cost function associated with the Cobb-Douglas production function is Cobb-Douglas:


\begin{equation*}
T C_{i}=r \cdot\left(A_{i} \alpha_{1}^{\alpha_{1}} \alpha_{2}^{\alpha_{2}} \alpha_{3}^{\alpha_{3}}\right)^{-1 / r} Q_{i}^{1 / r} p_{i 1}^{\alpha_{1} / r} p_{i 2}^{\alpha_{2} / r} p_{i 3}^{\alpha_{3} / r}, \tag{1.7.2}
\end{equation*}


where $T C_{i}$ is total costs for firm $i$. Taking logs, we obtain the following log-linear relationship:


\begin{equation*}
\log \left(T C_{i}\right)=\mu_{i}+\frac{1}{r} \log \left(Q_{i}\right)+\frac{\alpha_{1}}{r} \log \left(p_{i 1}\right)+\frac{\alpha_{2}}{r} \log \left(p_{i 2}\right)+\frac{\alpha_{3}}{r} \log \left(p_{i 3}\right), \tag{1.7.3}
\end{equation*}


where $\mu_{i}=\log \left[r \cdot\left(A_{i} \alpha_{1}^{\alpha_{1}} \alpha_{2}^{\alpha_{2}} \alpha_{3}^{\alpha_{3}}\right)^{-1 / r}\right]$. The equation is said to be log-linear because both the dependent variable and the regressors are logs. Coefficients in log-linear equations are elasticities. The $\log \left(p_{i 1}\right)$ coefficient, for example, is the elasticity of total costs with respect to the wage rate, i.e., the percentage change in total costs when the wage rate changes by 1 percent. The degree of returns to scale, which in (1.7.3) is the reciprocal of the output elasticity of total costs, is independent of the level of output.

Now let $\mu \equiv \mathrm{E}\left(\mu_{i}\right)$ and define $\varepsilon_{i} \equiv \mu_{i}-\mu$ so that $\mathrm{E}\left(\varepsilon_{i}\right)=0$. This $\varepsilon_{i}$ represents the inverse of the firm's production efficiency relative to the industry's average efficiency; firms with positive $\varepsilon_{i}$ are high-cost firms. With this notation, (1.7.3) becomes


\begin{equation*}
\log \left(T C_{i}\right)=\beta_{1}+\beta_{2} \log \left(Q_{i}\right)+\beta_{3} \log \left(p_{i 1}\right)+\beta_{4} \log \left(p_{i 2}\right)+\beta_{5} \log \left(p_{i 3}\right)+\varepsilon_{i}, \tag{1.7.4}
\end{equation*}


where


\begin{equation*}
\beta_{1}=\mu, \quad \beta_{2}=\frac{1}{r}, \quad \beta_{3}=\frac{\alpha_{1}}{r}, \quad \beta_{4}=\frac{\alpha_{2}}{r}, \quad \text { and } \quad \beta_{5}=\frac{\alpha_{3}}{r} . \tag{1.7.5}
\end{equation*}


Thus, the cost function has been cast in the regression format of Assumption 1.1 with $K=5$. We noted a moment ago that the simple-minded approach of plotting the average cost against output cannot account for the factor price effect. What we have shown is that under the Cobb-Douglas technology the factor price effect is controlled for by the inclusion in the cost function of the logs of factor prices. Because the equation is derived from an explicit description of the firm's technology, the error term as well as the regression coefficients have clear interpretations.

\section*{How Do We Know Things Are Cobb-Douglas?}
The Cobb-Douglas functional form is certainly a very convenient parameterization of technology. But how do we know that the true production function is CobbDouglas? The Cobb-Douglas form satisfies the properties, such as diminishing marginal productivities, that we normally require for the production function, but the Cobb-Douglas form is certainly not the only functional form with those desirable properties. A number of more general functional forms have been proposed in the literature, but the Cobb-Douglas form, despite its simplicity, has proved to be a surprisingly good description of technology. Nerlove's paper is one of the relatively few studies in which the Cobb-Douglas (log-linear) form is found to be inadequate, but it only underscores the importance of the Cobb-Douglas functional form as the benchmark from which one can usefully contemplate generalizations.

\section*{Are the OLS Assumptions Satisfied?}
To justify the use of least squares, we need to make sure that Assumptions 1.11.4 are satisfied for the equation (1.7.4). Evidently, Assumption 1.1 (linearity) is satisfied with

$$
y_{i}=\log \left(T C_{i}\right), \mathbf{x}_{i}=\left(1, \log \left(Q_{i}\right), \log \left(p_{i 1}\right), \log \left(p_{i 2}\right), \log \left(p_{i 3}\right)\right)^{\prime} .
$$

There is no reason to expect that the regressors in (1.7.4) are perfectly multicollinear. Indeed, in Nerlove's data set, $\operatorname{rank}(\mathbf{X})=5$ and $n=145$, so Assumption 1.3 (no multicollinearity) is satisfied as well.

In verifying the strict exogeneity assumption (Assumption 1.2), the features of the electricity industry mentioned above are relevant. It is reasonable to assume, as in most cross-section data, that $\mathbf{x}_{i}$ is independent of $\varepsilon_{j}$ for $i \neq j$. So the question is whether $\mathbf{x}_{i}$ is independent of $\varepsilon_{i}$. If it is, then $\mathrm{E}(\boldsymbol{\varepsilon} \mid \mathbf{X})=\mathbf{0}$. According to the third feature of the industry, factor prices are given to the firm with no regard for the firm's efficiency, so it is eminently reasonable to assume that factor prices are independent of $\varepsilon_{i}$.

What about output? Since the firm's output is supplied on demand (the first feature of the industry), output depends on the price of electricity set by the utility commission (the second feature). If the regulatory scheme is such that the price is determined regardless of the firm's efficiency, then $\log \left(Q_{i}\right)$ and $\varepsilon_{i}$ are independently distributed. On the other hand, if the price is set to cover the average cost, then the firm's efficiency affects output through the effect of the electricity price on demand and output in this case is endogenous, being correlated with the error term. We will very briefly come back to this point at the end, but until then we will ignore the possible endogeneity of output. This certainly would not do if we were dealing with a competitive industry. Since high-cost firms tend to produce less, there would be a negative correlation between $\log \left(Q_{i}\right)$ and $\varepsilon_{i}$, making OLS an inappropriate estimation procedure.

Regarding Assumption 1.4, the assumption of no correlation in the error term between firms (observations) would be suspect if, for example, there were technology spillovers running from one firm to other closely located firms. For the industry under study, this is probably not the case.

There is no a priori reason to suppose that homoskedasticity is satisfied. Indeed, the plot of residuals to be shown shortly suggests a failure of this condition. The main part of Nerlove's paper is exploring ways to deal with this problem.

\section*{Restricted Least Squares}
The equation (1.7.4) is overidentified in that its five coefficients, being functions of the four technology parameters (which are $\alpha_{1}, \alpha_{2}, \alpha_{3}$, and $\mu$ ), are not free parameters. We can easily see that from (1.7.5): $\beta_{3}+\beta_{4}+\beta_{5}=1$ (recall: $r \equiv \alpha_{1}+\alpha_{2}+\alpha_{3}$ ). This is a reflection of the generic property of the cost function that it is linearly homogeneous in factor prices. Indeed, multiplying total costs $T C_{i}$ and all factor prices ( $p_{i 1}, p_{i 2}, p_{i 3}$ ) by a common factor leaves the cost function (1.7.4) intact if and only if $\beta_{3}+\beta_{4}+\beta_{5}=1$.

Estimating the equation by least squares while imposing a priori restrictions on the coefficient vector is the restricted least squares. It can be done easily by deriving from the original regression a separate regression that embodies the restrictions. In the present example, to impose the homogeneity restriction $\beta_{3}+\beta_{4}+\beta_{5}=1$ on the cost function, we take any one of the factor prices, say $p_{i 3}$, and subtract $\log \left(p_{i 3}\right)$ from both sides of (1.7.4) to obtain


\begin{equation*}
\log \left(\frac{T C_{i}}{p_{i 3}}\right)=\beta_{1}+\beta_{2} \log \left(Q_{i}\right)+\beta_{3} \log \left(\frac{p_{i 1}}{p_{i 3}}\right)+\beta_{4} \log \left(\frac{p_{i 2}}{p_{i 3}}\right)+\varepsilon_{i} . \tag{1.7.6}
\end{equation*}


There are now four coefficients in the regression, from which unique values of the four technology parameters can be determined. The restricted least squares estimate of ( $\beta_{1}, \ldots, \beta_{4}$ ) is simply the OLS estimate of the coefficients in (1.7.6). The restricted least squares estimate of $\beta_{5}$ is the value implied by the estimate of ( $\beta_{1}, \ldots, \beta_{4}$ ) and the restriction.

\section*{Testing the Homogeneity of the Cost Function}
Before proceeding to the estimation of the restricted model (1.7.6), in order to test the homogeneity restriction $\beta_{3}+\beta_{4}+\beta_{5}=1$, we will first estimate the unrestricted model (1.7.4). If one uses the data available in printed form in Nerlove's paper, the OLS estimate of the equation is:


\begin{gather*}
\log \left(T C_{i}\right)=\underset{(1.8)}{-3.5}+\underset{(0.017)}{0.72} \log \left(Q_{i}\right)+\underset{(0.29)}{0.44 \log \left(p_{i 1}\right)} \\
-\underset{(0.34)}{0.22 \log \left(p_{i 2}\right)}+\underset{(0.10)}{0.43 \log \left(p_{i 3}\right)} \\
R^{2}=0.926, \text { mean of dep. variable }=1.72, \\
S E R=0.392, S S R=21.552, n=145 .
\end{gather*}


Here, numbers in parentheses are the standard errors of the OLS coefficient estimates. Since $\beta_{2}=1 / r$, the estimate of the degree of returns to scale implied\\
by the OLS coefficient estimates is about $1.4(=1 / 0.72)$. The OLS estimate of $\beta_{4}=\alpha_{2} / r$ has the wrong sign. As noted by Nerlove, there are reasons to believe that $p_{i 2}$, the rental price of capital, is poorly measured. This may explain why $b_{4}$ is so imprecisely determined (i.e., the standard error is large relative to the size of the coefficient estimate) that one cannot reject the hypothesis that $\beta_{4}=0$ with a $t$-ratio of $-0.65(=-0.22 / 0.34) .^{24}$

To test the homogeneity restriction $\mathrm{H}_{0}: \beta_{3}+\beta_{4}+\beta_{5}=1$, we could write the hypothesis in the form $\mathbf{R} \boldsymbol{\beta}=\mathbf{r}$ with $\mathbf{R}=(0,0,1,1,1)$ and $\mathbf{r}=1$ and use the formula (1.4.9) to calculate the $F$-ratio. The maintained hypothesis is the unrestricted model (1.7.4) (that is, Assumptions 1.1-1.5 where the equation in Assumption 1.1 is (1.7.4)), so the $\mathbf{b}$ and the estimated variance of $\mathbf{b}$ in the $F$-ratio formula should come from the OLS estimation of (1.7.4). Alternatively, we can use the $F$-ratio formula (1.4.11). The unrestricted model producing $S S R_{U}$ is (1.7.4) and the restricted model producing $S S R_{R}$ is (1.7.6), which superimposes the null hypothesis on the unrestricted model. The OLS estimate of (1.7.6) is


\begin{gather*}
\log \left(\frac{T C_{i}}{p_{i 3}}\right)=\underset{(0.88)}{-4.7}+\underset{(0.017)}{0.72 \log \left(Q_{i}\right)} \\
+\underset{(0.20)}{0.59 \log \left(p_{i 1} / p_{i 3}\right)-} \underset{(0.19)}{0.007} \log \left(p_{i 2} / p_{i 3}\right) \\
R^{2}=0.932, \text { mean of dep. var. }=-1.48 \\
S E R=0.39, S S R=21.640, n=145 \tag{1.7.8}
\end{gather*}


The $F$ test of the homogeneity restriction proceeds as follows.\\
Step 1: Using (1.4.11), the $F$-ratio can be calculated as

$$
\frac{(21.640-21.552) / 1}{21.552 /(145-5)}=0.57
$$

Step 2: Find the critical value. The number of restrictions (equations) in the null hypothesis is 1 , and $K$ (the number of coefficients) in the unrestricted model (which is the maintained hypothesis) is 5 . So the degrees of freedom are 1 and $140(=145-5)$. From the table of $F$ distributions, the critical value is about 3.9.\\
Step 3: Thus, we can easily accept the homogeneity restriction, a very comforting conclusion for those who take microeconomics seriously (like us).

\footnotetext{${ }^{24}$ The consequence of measurement error is not just that the coefficient of the variable measured with error is poorly determined; it could also contaminate the coefficient estimates for all other regressors. The appropriate context to address this problem is the large sample theory for endogeneous regressors in Chapter 3.
}\section*{Detour: A Cautionary Note on $\boldsymbol{R}^{\mathbf{2}}$}
The $R^{2}$ of 0.926 is surprisingly high for cross-section estimates, but some of the explanatory power of the regression comes from the scale effect that total costs increase with firm size. To gauge the contribution of the scale effect on the $R^{2}$, subtract $\log \left(Q_{i}\right)$ from both sides of (1.7.4) to obtain an equivalent cost function:


\begin{align*}
\log \left(\frac{T C_{i}}{Q_{i}}\right)= & \beta_{1}+\left(\beta_{2}-1\right) \log \left(Q_{i}\right) \\
& +\beta_{3} \log \left(p_{i 1}\right)+\beta_{4} \log \left(p_{i 2}\right)+\beta_{5} \log \left(p_{i 3}\right)+\varepsilon_{i} \tag{1.7.4'}
\end{align*}


Here, the dependent variable is the average cost rather than total costs. Application of the OLS to (1.7.4') using the same data yields


\begin{align*}
\log \left(\frac{T C_{i}}{Q_{i}}\right)= & \underset{(1.8)}{-3.5}-\underset{(0)}{0.28} \log \left(Q_{i}\right) \\
& +\underset{(0.29)}{0.44} \log \left(p_{i 1}\right)-\underset{(0.34)}{0.22} \log \left(p_{i 2}\right)+\underset{(0.10)}{0.43} \log \left(p_{i 3}\right) \\
& R^{2}=0.695, \text { mean of dep. var. }=-4.83 \\
& S E R=0.392, S S R=21.552, n=145
\end{align*}


As you no doubt have anticipated, the output coefficient is now $-0.28(=0.72-1)$ with the standard errors and the other coefficient estimates unchanged. The $R^{2}$ changes only because the dependent variable is different. It is nonsense to say that the higher $R^{2}$ makes (1.7.4) preferable to (1.7.4'), because the two equations represent the same model. The point is: when comparing equations on the basis of the fit, the equations must share the same dependent variable.

\section*{Testing Constant Returns to Scale}
As an application of the $t$-test, consider testing whether returns to scale are constant ( $r=1$ ). We take the maintained hypothesis to be the restricted model (1.7.6). Because $\beta_{2}$ (the log output coefficient) equals 1 if and only if $r=1$, the null hypothesis is that $\mathrm{H}_{0}: \beta_{2}=1$. The $t$-test of constant returns to scale proceeds as follows.

Step 1: Calculate the $t$-ratio for the hypothesis. From the estimation of the restricted model, we have $b_{2}=0.72$ with a standard error of 0.017 , so

$$
t \text {-ratio }=\frac{0.72-1}{0.017}=-16
$$

Because the maintained hypothesis here is the restricted model (1.7.6), $K$ (the number of coefficients) $=4$.\\
Step 2: Look for the critical value in the $t$ (141) distribution. If the size of the test is 5 percent, the critical value is 1.98 .\\
Step 3: Since the absolute value of the $t$-ratio is far greater than the critical value, we reject the hypothesis of constant returns to scale.

\section*{Importance of Plotting Residuals}
The regression has a problem that cannot be seen from the estimated coefficients and their standard errors. Figure 1.7 plots the residuals against $\log \left(Q_{i}\right)$. Notice two things from the plot. First, as output increases, the residuals first tend to be positive, then negative, and again positive. This strongly suggests that the degree of returns to scale ( $r$ ) is not constant as assumed in the log-linear specification. Second, the residuals are more widely scattered for lower outputs, which is a sign of a failure of the homoskedasticity assumption that the error variance does not depend on the regressors. To deal with these problems, Nerlove divided the sample of 145 firms into five groups of 29 , ordered by output, and estimated the model (1.7.6) separately for each group. This amounts to allowing all the coefficients (including $\beta_{2}=1 / r$ ) and the error variance to differ across the five groups differing in size. Nerlove finds that returns to scale diminish steadily, from a high of well over 2 to a low of slightly below 1 , over the output range of the data. In the empirical exercise of this chapter, the reader is asked to replicate this finding and do some further analysis using dummy variables and the weighted least squares.

\section*{Subsequent Developments}
One strand of the subsequent literature is concerned about generalizing the CobbDouglas technology while maintaining the assumption of cost minimization. An obvious alternative to Cobb-Douglas is the Constant Elasticity of Substitution (CES) production function, but it has two problems. First, the cost function implied by the CES production function is highly nonlinear (which, though, could be overcome by the use of nonlinear least squares to be covered in Chapter 7). Second, the CES technology implies a constant degree of returns to scale. One of Nerlove's main findings is that the degree varies with output. Christensen and Greene (1976) are probably the first to estimate the technology parameters allowing for variable degrees of returns to scale. Using the translog cost function introduced by Christensen, Jorgenson, and Lau (1973), they find that the significant scale economies evident in the 1955 data were mostly exhausted by 1970, with most firms operating at much higher output levels where the AC curve is essentially flat. Their work will be examined in detail in Chapter 4.

\begin{figure}[h]
\begin{center}
  \includegraphics[alt={},max width=\textwidth]{14fb0795-6d81-4e76-9788-792298f1ddae-086_753_1135_373_470}
\captionsetup{labelformat=empty}
\caption{Figure 1.7: Plot of Residuals against Log Output}
\end{center}
\end{figure}

Another issue is whether regulated firms minimize costs. The influential paper by Averch and Johnson (1962) argues that the practice by regulators to guarantee utilities a "fair rate of return" on their capital stock distorts the choice of input levels. Since the fair rate of return is usually higher than the interest rate, utilities have an incentive to overinvest. That is, they minimize costs, but the relevant rate of return in the definition of the user cost of capital is the fair rate of return. Consequently, unless the fair rate of return is used in the calculation of $p_{i 2}$, the true technology parameters cannot be estimated from the cost function. The fair-rate-of-return regulation creates another econometric problem: to guarantee utilities a fair rate of return, the price of electricity must be kept relatively high in markets served by high-cost utilities. Thus output will be endogenous.

A more recent issue is whether the regulator has enough information to bring about cost minimization. If the utility has more information about costs, it has an incentive to misreport to the regulator the true value of the efficiency parameter. Schemes to be adopted by the regulator to take into account this incentive problem may not lead to cost minimization. Wolak’s (1994) empirical results for California's water utility industry indicate that the observed level of costs and output is better modeled as the outcome of a regulator-utility interaction under asymmetric information. Wolak resolves the problem of the endogeneity of output by estimating the demand function along with the cost function. Doing so, however, requires an estimation technique more sophisticated than the OLS.

\section*{QUESTIONS FOR REVIEW}
\begin{enumerate}
  \item (Review of duality theory) Consult your favorite microeconomic textbook to remember how to derive the Cobb-Douglas cost function from the CobbDouglas production function.
  \item (Change of units) In Nerlove's data, output is measured in kilowatt hours. If output were measured in megawatt hours, how would the estimated restricted regression change?
  \item (Recovering technology parameters from regression coefficients) Show that the technology parameters ( $\mu, \alpha_{1}, \alpha_{2}, \alpha_{3}$ ) can be determined uniquely from the first four equations in (1.7.5) and the definition $r \equiv \alpha_{1}+\alpha_{2}+\alpha_{3}$. (Do not use the fifth equation $\beta_{5}=\alpha_{3} / r$.)
  \item (Recovering left-out coefficients from restricted OLS) Calculate the restricted OLS estimate of $\beta_{5}$ from (1.7.8). How do you calculate the standard error of $b_{5}$ from the printout of the restricted OLS? Hint: Write $b_{5}=a+\mathbf{c}^{\prime} \mathbf{b}$ for suitably chosen $a$ and $\mathbf{c}$ where $\mathbf{b}$ here is $\left(b_{1}, \ldots, b_{4}\right)^{\prime}$. So $\operatorname{Var}\left(b_{5} \mid \mathbf{X}\right)=\mathbf{c}^{\prime} \operatorname{Var}(\mathbf{b} \mid \mathbf{X}) \mathbf{c}$. The printout from the restricted OLS should include $\widehat{\operatorname{Var}(\mathbf{b} \mid \mathbf{X}) \text {. }}$
  \item If you take $p_{i 2}$ instead of $p_{i 3}$ and subtract $\log \left(p_{i 2}\right)$ from both sides of (1.7.4), how does the restricted regression look? Without actually estimating it on Nerlove's data, can you tell from the estimated restricted regression in the text what the restricted OLS estimate of $\left(\beta_{1}, \ldots, \beta_{5}\right)$ will be? Their standard errors? The $S S R$ ? What about the $R^{2}$ ?
  \item Why is the $R^{2}$ of 0.926 from the unrestricted model (1.7.7) lower than the $R^{2}$ of 0.932 from the restricted model (1.7.8)?
  \item A more realistic assumption about the rental price of capital may be that there is an economy-wide capital market so $p_{i 2}$ is the same across firms. In this case,\\
(a) Can we estimate the technology parameters? Hint: The answer is yes, but why? When $p_{i 2}$ is constant, (1.7.4) will have the perfect multicollinearity problem. But recall that ( $\beta_{1}, \ldots, \beta_{5}$ ) are not free parameters.\\
(b) Can we test homogeneity of the cost function in factor prices?
  \item Taking logs of both sides of the production function (1.7.1), one can derive the log-linear relationship:
\end{enumerate}

$$
\log \left(Q_{i}\right)=\alpha_{0}+\alpha_{1} \log \left(x_{i 1}\right)+\alpha_{2} \log \left(x_{i 2}\right)+\alpha_{3} \log \left(x_{i 3}\right)+\varepsilon_{i},
$$

where $\varepsilon_{i}$ here is defined as $\log \left(A_{i}\right)-\mathrm{E}\left[\log \left(A_{i}\right)\right]$ and $\alpha_{0}=\mathrm{E}\left[\log \left(A_{i}\right)\right]$. Suppose, in addition to total costs, output, and factor prices, we had data on factor inputs. Can we estimate $\alpha$ 's by applying OLS to this log-linear relationship? Why or why not? Hint: Do input levels depend on $\varepsilon_{i}$ ? Suggest a different way to estimate $\alpha$ 's. Hint: Look at input shares.

\section*{PROBLEM SET FOR CHAPTER 1}
$\_\_\_\_$

\section*{ANALYTICAL EXERCISES}
\begin{enumerate}
  \item (Proof that $\mathbf{b}$ minimizes $S S R$ ) Let $\mathbf{b}$ be the OLS estimator of $\boldsymbol{\beta}$. Prove that, for any hypothetical estimate, $\widetilde{\boldsymbol{\beta}}$, of $\boldsymbol{\beta}$,
\end{enumerate}

$$
(\mathbf{y}-\mathbf{X} \tilde{\boldsymbol{\beta}})^{\prime}(\mathbf{y}-\mathbf{X} \tilde{\boldsymbol{\beta}}) \geq(\mathbf{y}-\mathbf{X} \mathbf{b})^{\prime}(\mathbf{y}-\mathbf{X} \mathbf{b})
$$

In your proof, use the add-and-subtract strategy: take $\mathbf{y}-\mathbf{X} \boldsymbol{\boldsymbol { \beta }}$, add $\mathbf{X b}$ to it and then subtract the same from it. It produces the decomposition of $\mathbf{y}-\mathbf{X} \widetilde{\boldsymbol{\beta}}$ :

$$
\mathbf{y}-\mathbf{X} \tilde{\boldsymbol{\beta}}=(\mathbf{y}-\mathbf{X} \mathbf{b})+(\mathbf{X} \mathbf{b}-\mathbf{X} \tilde{\boldsymbol{\beta}})
$$

Hint: $(\mathbf{y}-\mathbf{X} \widetilde{\boldsymbol{\beta}})^{\prime}(\mathbf{y}-\mathbf{X} \widetilde{\boldsymbol{\beta}})=[(\mathbf{y}-\mathbf{X b})+\mathbf{X}(\mathbf{b}-\widetilde{\boldsymbol{\beta}})]^{\prime}[(\mathbf{y}-\mathbf{X b})+\mathbf{X}(\mathbf{b}-\widetilde{\boldsymbol{\beta}})]$. Using the normal equations, show that this equals

$$
(\mathbf{y}-\mathbf{X b})^{\prime}(\mathbf{y}-\mathbf{X b})+(\mathbf{b}-\widetilde{\boldsymbol{\beta}})^{\prime} \mathbf{X}^{\prime} \mathbf{X}(\mathbf{b}-\widetilde{\boldsymbol{\beta}})
$$

\begin{enumerate}
  \setcounter{enumi}{1}
  \item (The annihilator associated with the vector of ones) Let $\mathbf{1}$ be the $n$-dimensional column vector of ones, and let $\mathbf{M}_{1} \equiv \mathbf{I}_{n}-\mathbf{1}\left(\mathbf{1}^{\prime} \mathbf{1}\right)^{-1} \mathbf{1}^{\prime}$. That is, $\mathbf{M}_{1}$ is the annihilator associated with 1. Prove the following:\\
(a) $\mathbf{M}_{1}$ is symmetric and idempotent.\\
(b) $\mathbf{M}_{1} \mathbf{1}=\mathbf{0}$.\\
(c) $\mathbf{M}_{1} \mathbf{y}=\mathbf{y}-\bar{y} \cdot \mathbf{1}$ where
\end{enumerate}

$$
\bar{y}=\frac{1}{n} \sum_{i=1}^{n} y_{i}
$$

$\mathbf{M}_{1} \mathbf{y}$ is the vector of deviations from the mean.\\
(d) $\mathbf{M}_{1} \mathbf{X}=\mathbf{X}-\mathbf{1}^{\prime}$ where $\overline{\mathbf{x}}=\mathbf{X}^{\prime} \mathbf{1} / n$. The $k$-th element of the $K \times 1$ vector $\overline{\mathbf{x}}$ is $\frac{1}{n} \sum_{i=1}^{n} X_{i k}$.\\
3. (Deviation-from-the-mean regression) Consider a regression model with a constant. Let $\mathbf{X}$ be partitioned as

$$
\underset{(n \times K)}{\mathbf{X}}=\left[\underset{n \times 1}{\mathbf{1}} \vdots \underset{n \times(K-1)}{\mathbf{X}_{2}}\right]
$$

so the first regressor is a constant. Partition $\boldsymbol{\beta}$ and $\mathbf{b}$ accordingly:

$$
\boldsymbol{\beta}=\left[\begin{array}{l}
\beta_{1} \\
\boldsymbol{\beta}_{2}
\end{array}\right] \begin{aligned}
& \leftarrow \text { scalar } \\
& \leftarrow(K-1) \times 1
\end{aligned}, \quad \mathbf{b}=\left[\begin{array}{l}
b_{1} \\
\mathbf{b}_{2}
\end{array}\right] .
$$

Also let $\widetilde{\mathbf{X}}_{2} \equiv \mathbf{M}_{1} \mathbf{X}_{2}$ and $\tilde{\mathbf{y}} \equiv \mathbf{M}_{1} \mathbf{y}$. They are the deviations from the mean for the nonconstant regressors and the dependent variable. Prove the following:\\
(a) The $K$ normal equations are

$$
\bar{y}-b_{1}-\overline{\mathbf{x}}_{2}^{\prime} \mathbf{b}_{2}=0
$$

where $\overline{\mathbf{x}}_{2}=\mathbf{X}_{2}^{\prime} \mathbf{1} / n$,

$$
\mathbf{X}_{2}^{\prime} \mathbf{y}-n \cdot b_{1} \cdot \overline{\mathbf{x}}_{2}-\mathbf{X}_{2}^{\prime} \mathbf{X}_{2} \mathbf{b}_{2}=\underset{((K-1) \times 1)}{\mathbf{0}} .
$$

(b) $\mathbf{b}_{2}=\left(\widetilde{\mathbf{X}}_{2}^{\prime} \widetilde{\mathbf{X}}_{2}\right)^{-1} \widetilde{\mathbf{X}}_{2}^{\prime} \tilde{\mathbf{y}}$. Hint: Substitute the first normal equation into the other $K-1$ equations to eliminate $b_{1}$ and solve for $\mathbf{b}_{2}$. This is a generalization of the result you proved in Review Question 3 in Section 1.2.\\
4. (Partitioned regression, generalization of Exercise 3) Let $\mathbf{X}$ be partitioned as

$$
\underset{(n \times K)}{\mathbf{X}}=\left[\underset{\left(n \times K_{1}\right)}{\mathbf{X}_{1}}: \underset{\left(n \times K_{2}\right)}{\mathbf{X}_{2}}\right] .
$$

Partition $\boldsymbol{\beta}$ accordingly:

$$
\boldsymbol{\beta}=\left[\begin{array}{l}
\boldsymbol{\beta}_{1} \\
\boldsymbol{\beta}_{2}
\end{array}\right] \begin{aligned}
& \leftarrow K_{1} \times 1 \\
& \leftarrow K_{2} \times 1
\end{aligned} .
$$

Thus, the regression can be written as

$$
\mathbf{y}=\mathbf{X}_{1} \boldsymbol{\beta}_{1}+\mathbf{X}_{2} \boldsymbol{\beta}_{2}+\boldsymbol{\varepsilon}
$$

Let $\mathbf{P}_{1} \equiv \mathbf{X}_{1}\left(\mathbf{X}_{1}^{\prime} \mathbf{X}_{1}\right)^{-1} \mathbf{X}_{1}^{\prime}, \mathbf{M}_{1} \equiv \mathbf{I}-\mathbf{P}_{1}, \widetilde{\mathbf{X}}_{2} \equiv \mathbf{M}_{1} \mathbf{X}_{2}$ and $\tilde{\mathbf{y}} \equiv \mathbf{M}_{1} \mathbf{y}$. Thus, $\tilde{\mathbf{y}}$ is the residual vector from the regression of $\mathbf{y}$ on $\mathbf{X}_{1}$, and the $k$-th column of $\widetilde{\mathbf{X}}_{2}$ is the residual vector from the regression of the corresponding $k$-th column of\\
$\mathbf{X}_{2}$ on $\mathbf{X}_{1}$. Prove the following:\\
(a) The normal equations are


\begin{align*}
& \mathbf{X}_{1}^{\prime} \mathbf{X}_{1} \mathbf{b}_{1}+\mathbf{X}_{1}^{\prime} \mathbf{X}_{2} \mathbf{b}_{2}=\mathbf{X}_{1}^{\prime} \mathbf{y}  \tag{*}\\
& \mathbf{X}_{2}^{\prime} \mathbf{X}_{1} \mathbf{b}_{1}+\mathbf{X}_{2}^{\prime} \mathbf{X}_{2} \mathbf{b}_{2}=\mathbf{X}_{2}^{\prime} \mathbf{y} \tag{**}
\end{align*}


(b) $\mathbf{b}_{2}=\left(\widetilde{\mathbf{X}}_{2}^{\prime} \widetilde{\mathbf{X}}_{2}\right)^{-1} \widetilde{\mathbf{X}}_{2}^{\prime} \tilde{\mathbf{y}}$. That is, $\mathbf{b}_{2}$ can be obtained by regressing the residuals $\tilde{\mathbf{y}}$ on the matrix of residuals $\tilde{\mathbf{X}}_{2}$. Hint: Derive $\mathbf{X}_{1} \boldsymbol{\beta}_{1}=-\mathbf{P}_{1} \mathbf{X}_{2} \boldsymbol{\beta}_{2}+\mathbf{P}_{1} \mathbf{y}$ from ( $*$ ). Substitute this into ( $* *$ ) to obtain $\mathbf{X}_{2}^{\prime} \mathbf{M}_{1} \mathbf{X}_{2} \boldsymbol{\beta}_{2}=\mathbf{X}_{2}^{\prime} \mathbf{M}_{1} \mathbf{y}$. Then use the fact that $\mathbf{M}_{1}$ is symmetric and idempotent. Or, if you wish, you can apply the brute force of the partitioned inverse formula (A.10) of Appendix A to the coefficient matrix

$$
\mathbf{X}^{\prime} \mathbf{X}=\left[\begin{array}{ll}
\mathbf{X}_{1}^{\prime} \mathbf{X}_{1} & \mathbf{X}_{1}^{\prime} \mathbf{X}_{2} \\
\mathbf{X}_{2}^{\prime} \mathbf{X}_{1} & \mathbf{X}_{2}^{\prime} \mathbf{X}_{2}
\end{array}\right]
$$

Show that the second diagonal block of $\left(\mathbf{X}^{\prime} \mathbf{X}\right)^{-1}$ is $\left(\widetilde{\mathbf{X}}_{2}^{\prime} \widetilde{\mathbf{X}}_{2}\right)^{-1}$.\\
(c) The residuals from the regression of $\tilde{\mathbf{y}}$ on $\tilde{\mathbf{X}}_{2}$ numerically equals $\mathbf{e}$, the residuals from the regression of $\mathbf{y}$ on $\mathbf{X}\left(\equiv\left(\mathbf{X}_{1} \vdots \mathbf{X}_{2}\right)\right)$. Hint: If $\mathbf{e}$ is the residual from the regression of $\mathbf{y}$ on $\mathbf{X}$,

$$
\mathbf{y}=\mathbf{X}_{1} \mathbf{b}_{1}+\mathbf{X}_{2} \mathbf{b}_{2}+\mathbf{e}
$$

Premultiplying both sides by $\mathbf{M}_{1}$ and using $\mathbf{M}_{1} \mathbf{X}_{1}=\mathbf{0}$, we obtain

$$
\tilde{\mathbf{y}}=\widetilde{\mathbf{X}}_{2} \mathbf{b}_{2}+\mathbf{M}_{1} \mathbf{e}
$$

Show that $\mathbf{M}_{1} \mathbf{e}=\mathbf{e}$ and observe that $\mathbf{b}_{2}$ equals the OLS coefficient estimate in the regression of $\tilde{\mathbf{y}}$ on $\tilde{\mathbf{X}}_{2}$.\\
(d) $\mathbf{b}_{2}=\left(\widetilde{\mathbf{X}}_{2}^{\prime} \widetilde{\mathbf{X}}_{2}\right)^{-1} \widetilde{\mathbf{X}}_{2}^{\prime} \mathbf{y}$. Note the difference from (b). Here, the vector of dependent variable is $\mathbf{y}$, not $\tilde{\mathbf{y}}$. Are the residuals from the regression of $\mathbf{y}$ on $\widetilde{\mathbf{X}}_{2}$ numerically the same as $\mathbf{e}$ ? [Answer: No.] Is the SSR from the regression of $\mathbf{y}$ on $\widetilde{\mathbf{X}}_{2}$ the same as the SSR from the regression of $\tilde{\mathbf{y}}$ on $\widetilde{\mathbf{X}}_{2}$ ? [Answer: No.]\\
The results in (b)-(d) are known as the Frisch-Waugh Theorem.\\
(e) Show:

$$
\tilde{\mathbf{y}}^{\prime} \tilde{\mathbf{y}}-\mathbf{e}^{\prime} \mathbf{e}=\tilde{\mathbf{y}}^{\prime} \mathbf{X}_{2}\left(\mathbf{X}_{2}^{\prime} \mathbf{M}_{1} \mathbf{X}_{2}\right)^{-1} \mathbf{X}_{2}^{\prime} \tilde{\mathbf{y}}
$$

Hint: Apply the general decomposition formula (1.2.15) to the regression in (c) to derive

$$
\tilde{\mathbf{y}}^{\prime} \tilde{\mathbf{y}}=\mathbf{b}_{2}^{\prime} \tilde{\mathbf{X}}_{2}^{\prime} \tilde{\mathbf{X}}_{2} \mathbf{b}_{2}+\mathbf{e}^{\prime} \mathbf{e}
$$

Then use (b).\\
(f) Consider the following four regressions:\\
(1) regress $\tilde{\mathbf{y}}$ on $\mathbf{X}_{1}$.\\
(2) regress $\tilde{\mathbf{y}}$ on $\widetilde{\mathbf{X}}_{2}$.\\
(3) regress $\tilde{\mathbf{y}}$ on $\mathbf{X}_{1}$ and $\mathbf{X}_{2}$.\\
(4) regress $\tilde{\mathbf{y}}$ on $\mathbf{X}_{2}$.

Let $S S R_{j}$ be the sum of squared residuals from regression $j$. Show:\\
(i) $S S R_{1}=\tilde{\mathbf{y}}^{\prime} \tilde{\mathbf{y}}$. Hint: $\tilde{\mathbf{y}}$ is constructed so that $\mathbf{X}_{1}^{\prime} \tilde{\mathbf{y}}=\mathbf{0}$, so $\mathbf{X}_{1}$ should have no explanatory power.\\
(ii) $S S R_{2}=\mathbf{e}^{\prime} \mathbf{e}$. Hint: Use (c).\\
(iii) $S S R_{3}=\mathbf{e}^{\prime} \mathbf{e}$. Hint: Apply the Frisch-Waugh Theorem on regression (3). $\mathbf{M}_{1} \tilde{\mathbf{y}}=\tilde{\mathbf{y}}$.\\
(iv) Verify by numerical example that $S S R_{4}$ is not necessarily equal to $\mathbf{e}^{\prime} \mathbf{e}$.\\
5. (Restricted regression and $F$ ) In the restricted least squares, the sum of squared residuals is minimized subject to the constraint implied by the null hypothesis $\mathbf{R} \boldsymbol{\beta}=\mathbf{r}$. Form the Lagrangian as

$$
\mathcal{L}=\frac{1}{2}(\mathbf{y}-\mathbf{X} \tilde{\boldsymbol{\beta}})^{\prime}(\mathbf{y}-\mathbf{X} \tilde{\boldsymbol{\beta}})+\lambda^{\prime}(\mathbf{R} \tilde{\boldsymbol{\beta}}-\mathbf{r}),
$$

where $\boldsymbol{\lambda}$ here is the \#r-dimensional vector of Lagrange multipliers (recall: $\mathbf{R}$ is $\# \mathbf{r} \times K, \widetilde{\boldsymbol{\beta}}$ is $K \times 1$, and $\mathbf{r}$ is $\# \mathbf{r} \times 1$ ). Let $\widehat{\boldsymbol{\beta}}$ be the restricted least squares estimator of $\boldsymbol{\beta}$. It is the solution to the constrained minimization problem.\\
(a) Let $\mathbf{b}$ be the unrestricted OLS estimator. Show:

$$
\begin{aligned}
\widehat{\boldsymbol{\beta}} & =\mathbf{b}-\left(\mathbf{X}^{\prime} \mathbf{X}\right)^{-1} \mathbf{R}^{\prime}\left[\mathbf{R}\left(\mathbf{X}^{\prime} \mathbf{X}\right)^{-1} \mathbf{R}^{\prime}\right]^{-1}(\mathbf{R} \mathbf{b}-\mathbf{r}) \\
\lambda & =\left[\mathbf{R}\left(\mathbf{X}^{\prime} \mathbf{X}\right)^{-1} \mathbf{R}^{\prime}\right]^{-1}(\mathbf{R} \mathbf{b}-\mathbf{r})
\end{aligned}
$$

Hint: The first-order conditions are $\mathbf{X}^{\prime} \mathbf{y}-\left(\mathbf{X}^{\prime} \mathbf{X}\right) \widehat{\boldsymbol{\beta}}=\mathbf{R}^{\prime} \boldsymbol{\lambda}$ or $\mathbf{X}^{\prime}(\mathbf{y}-\mathbf{X} \widehat{\boldsymbol{\beta}})= \mathbf{R}^{\prime} \boldsymbol{\lambda}$. Combine this with the constraint $\mathbf{R} \widehat{\boldsymbol{\beta}}=\mathbf{r}$ to solve for $\boldsymbol{\lambda}$ and $\widehat{\boldsymbol{\beta}}$.\\
(b) Let $\hat{\boldsymbol{\varepsilon}} \equiv \mathbf{y}-\mathbf{X} \widehat{\boldsymbol{\beta}}$, the residuals from the restricted regression. Show:

$$
\begin{aligned}
S S R_{R}-S S R_{U} & =(\mathbf{b}-\widehat{\boldsymbol{\beta}})^{\prime}\left(\mathbf{X}^{\prime} \mathbf{X}\right)(\mathbf{b}-\widehat{\boldsymbol{\beta}}) \\
& =(\mathbf{R} \mathbf{b}-\mathbf{r})^{\prime}\left[\mathbf{R}\left(\mathbf{X}^{\prime} \mathbf{X}\right)^{-1} \mathbf{R}^{\prime}\right]^{-1}(\mathbf{R} \mathbf{b}-\mathbf{r}) \\
& =\lambda^{\prime} \mathbf{R}\left(\mathbf{X}^{\prime} \mathbf{X}\right)^{-1} \mathbf{R}^{\prime} \lambda \\
& =\hat{\boldsymbol{\varepsilon}}^{\prime} \mathbf{P} \hat{\boldsymbol{\varepsilon}}
\end{aligned}
$$

where $\mathbf{P}$ is the projection matrix. Hint: For the first equality, use the add-and-subtract strategy:

$$
\begin{aligned}
S S R_{R} & =(\mathbf{y}-\mathbf{X} \widehat{\boldsymbol{\beta}})^{\prime}(\mathbf{y}-\mathbf{X} \widehat{\boldsymbol{\beta}}) \\
& =[(\mathbf{y}-\mathbf{X} \mathbf{b})+\mathbf{X}(\mathbf{b}-\widehat{\boldsymbol{\beta}})]^{\prime}[(\mathbf{y}-\mathbf{X} \mathbf{b})+\mathbf{X}(\mathbf{b}-\widehat{\boldsymbol{\beta}})]
\end{aligned}
$$

Use the normal equations $\mathbf{X}^{\prime}(\mathbf{y}-\mathbf{X b})=\mathbf{0}$. For the second and third equalities, use (a). To prove the fourth equality, the easiest way is to use the first-order condition mentioned in (a) that $\mathbf{R}^{\prime} \boldsymbol{\lambda}=\mathbf{X}^{\prime} \boldsymbol{\hat { \varepsilon }}$.\\
(c) Verify that you have proved in (b) that (1.4.9) $=(1.4 .11)$.\\
6. (Proof of the decomposition (1.2.17)) Take the unrestricted model to be a regression where one of the regressors is a constant, and the restricted model to be a regression where the only regressor is a constant.\\
(a) Show that (b) in the previous exercise is the decomposition (1.2.17) for this case. Hint: What is $\widehat{\boldsymbol{\beta}}$ for this case? Show that $\operatorname{SSR}_{R}=\sum_{i}\left(y_{i}-\bar{y}\right)^{2}$ and $(\mathbf{b}-\widehat{\boldsymbol{\beta}})^{\prime}\left(\mathbf{X}^{\prime} \mathbf{X}\right)(\mathbf{b}-\widehat{\boldsymbol{\beta}})=\sum_{i}(\hat{y}-\bar{y})^{2}$.\\
(b) ( $R^{2}$ as an $F$-ratio) For a regression where one of the regressors is a constant, prove that

$$
F=\frac{R^{2} /(K-1)}{\left(1-R^{2}\right) /(n-K)}
$$

\begin{enumerate}
  \setcounter{enumi}{6}
  \item (Hausman principle in finite samples) For the generalized regression model, prove the following. Here, it is understood that the expectations, variances, and covariances are all conditional on $\mathbf{X}$.\\
(a) $\operatorname{Cov}\left(\widehat{\boldsymbol{\beta}}_{\mathrm{GLS}}, \mathbf{b}-\widehat{\boldsymbol{\beta}}_{\mathrm{GLS}}\right)=\mathbf{0}$. Hint: Recall that, for any two random vectors $\mathbf{x}$ and $\mathbf{y}$,
\end{enumerate}

$$
\operatorname{Cov}(\mathbf{x}, \mathbf{y}) \equiv \mathrm{E}\left[(\mathbf{x}-\mathrm{E}(\mathbf{x}))(\mathbf{y}-\mathrm{E}(\mathbf{y}))^{\prime}\right]
$$

So

$$
\operatorname{Cov}(\mathbf{A x}, \mathbf{B y})=\mathbf{A} \operatorname{Cov}(\mathbf{x}, \mathbf{y}) \mathbf{B}^{\prime}
$$

Also, since $\boldsymbol{\beta}$ is nonrandom,

$$
\operatorname{Cov}\left(\widehat{\boldsymbol{\beta}}_{\mathrm{GLS}}, \mathbf{b}-\widehat{\boldsymbol{\beta}}_{\mathrm{GLS}}\right)=\operatorname{Cov}\left(\widehat{\boldsymbol{\beta}}_{\mathrm{GLS}}-\boldsymbol{\beta}, \mathbf{b}-\widehat{\boldsymbol{\beta}}_{\mathrm{GLS}}\right)
$$

(b) Let $\widetilde{\boldsymbol{\beta}}$ be any unbiased estimator and define $\mathbf{q} \equiv \widetilde{\boldsymbol{\beta}}-\widehat{\boldsymbol{\beta}}_{\mathrm{GLS}}$. Assume $\widetilde{\boldsymbol{\beta}}$ is such that $\mathbf{V}_{\mathbf{q}} \equiv \operatorname{Var}(\mathbf{q})$ is nonsingular. Prove: $\operatorname{Cov}\left(\widehat{\boldsymbol{\beta}}_{\mathrm{GLS}}, \mathbf{q}\right)=\mathbf{0}$. (If we set $\widetilde{\boldsymbol{\beta}}=\mathbf{b}$, we are back to (a).) Hint: Define: $\widehat{\boldsymbol{\beta}} \equiv \widehat{\boldsymbol{\beta}}_{\mathrm{GLS}}+\mathbf{H q}$ for some $\mathbf{H}$. Show:

$$
\operatorname{Var}(\widehat{\boldsymbol{\beta}})=\operatorname{Var}\left(\widehat{\boldsymbol{\beta}}_{\mathrm{GLS}}\right)+\mathbf{C H}^{\prime}+\mathbf{H C}^{\prime}+\mathbf{H V}_{\mathbf{q}} \mathbf{H}^{\prime}
$$

where $\mathbf{C} \equiv \operatorname{Cov}\left(\widehat{\boldsymbol{\beta}}_{\mathrm{GLS}}, \mathbf{q}\right)$. Show that, if $\mathbf{C} \neq \mathbf{0}$ then $\operatorname{Var}(\widehat{\boldsymbol{\beta}})$ can be made smaller than $\operatorname{Var}\left(\widehat{\boldsymbol{\beta}}_{\text {GLS }}\right)$ by setting $\mathbf{H}=-\mathbf{C V}_{\mathbf{q}}^{-1}$. Argue that this is in contradiction to Proposition 1.7(c).\\
(c) (Optional, only for those who are proficient in linear algebra) Prove: if the $K$ columns of $\mathbf{X}$ are characteristic vectors of $\mathbf{V}$, then $\mathbf{b}=\widehat{\boldsymbol{\beta}}_{\mathrm{GLS}}$, where $\mathbf{V}$ is the $n \times n$ variance-covariance matrix of the $n$-dimensional error vector $\boldsymbol{\varepsilon}$. (So not all unbiased estimators satisfy the requirement in (b) that $\operatorname{Var}(\widetilde{\boldsymbol{\beta}}- \widehat{\boldsymbol{\beta}}_{\text {GLS }}$ ) be nonsingular.) Hint: For any $n \times n$ symmetric matrix $\mathbf{V}$, there exists an $n \times n$ matrix $\mathbf{H}$ such that $\mathbf{H}^{\prime} \mathbf{H}=\mathbf{I}_{n}$ (so $\mathbf{H}$ is an orthogonal matrix) and $\mathbf{H}^{\prime} \mathbf{V H}=\boldsymbol{\Lambda}$, where $\boldsymbol{\Lambda}$ is a diagonal matrix with the characteristic roots (which are real since $\mathbf{V}$ is symmetric) of $\mathbf{V}$ in the diagonal. The columns of $\mathbf{H}$ are called the characteristic vectors of $\mathbf{V}$. Show that

$$
\mathbf{H}^{-1}=\mathbf{H}^{\prime}, \quad \mathbf{H}^{\prime} \mathbf{V}^{-1} \mathbf{H}=\mathbf{\Lambda}^{-1}, \quad \mathbf{H}^{\prime} \mathbf{V}^{-1}=\mathbf{\Lambda}^{-1} \mathbf{H}^{\prime}
$$

Without loss of generality, $\mathbf{X}$ can be taken to be the first $K$ columns of $\mathbf{H}$. So $\mathbf{X}=\mathbf{H F}$, where

$$
\underset{(n \times K)}{\mathbf{F}}=\left[\begin{array}{c}
\mathbf{I}_{K} \\
\mathbf{0}
\end{array}\right]
$$

\section*{EMPIRICAL EXERCISES}
Read Marc Nerlove, "Returns to Scale in Electricity Supply" (except paragraphs of equations (6)-(9), the part of section 2 from p. 184 on, and Appendix A and\\
C) before doing this exercise. For 145 electric utility companies in 1955, the file NERLOVE.ASC has data on the following:

Column 1: total costs (call it TC) in millions of dollars\\
Column 2: output ( $Q$ ) in billions of kilowatt hours\\
Column 3: price of labor ( $P L$ )\\
Column 4: price of fuels ( $P F$ )\\
Column 5: price of capital (PK).\\
They are from the data appendix of his article. There are 145 observations, and the observations are ordered in size, observation 1 being the smallest company and observation 145 the largest. Using the data transformation facilities of your computer software, generate for each of the 145 firms the variables required for estimation. To estimate (1.7.4), for example, you need to generate $\log (T C)$, a constant, $\log (Q), \log (P L), \log (P K)$, and $\log (P F)$, for each of the 145 firms.\\
(a) (Data question) Does Nerlove's construction of the price of capital conform to the definition of the user cost of capital? Hint: Read Nerlove's Appendix B.4.\\
(b) Estimate the unrestricted model (1.7.4) by OLS. Can you replicate the estimates in the text?\\
(c) (Restricted least squares) Estimate the restricted model (1.7.6) by OLS. To do this, you need to generate a new set of variables for each of the 145 firms. For example, the dependent variable is $\log (T C / P F)$, not $\log (T C)$. Can you replicate the estimates in the text? Can you replicate Nerlove's results? Nerlove's estimate of $\beta_{2}$, for example, is 0.721 with a standard error of 0.0174 (the standard error in his paper is 0.175 , but it is probably a typographical error). Where in Nerlove's paper can you find this estimate? What about the other coefficients? (Warning: You will not be able to replicate Nerlove's results precisely. One reason is that he used common rather than natural logarithms; however, this should affect only the estimated intercept term. The other reason: the data set used for his results is a corrected version of the data set published with his article.)

As mentioned in the text, the plot of residuals suggests a nonlinear relationship between $\log (T C)$ and $\log (Q)$. Nerlove hypothesized that estimated returns to scale varied with the level of output. Following Nerlove, divide the sample of 145 firms into five subsamples or groups, each having 29 firms. (Recall that since the data are ordered by level of output, the first 29 observations will have the smallest output levels, whereas the last 29 observations will have the largest output levels.) Consider the following three generalizations of the model (1.7.6):

Model 1: Both the coefficients ( $\beta$ 's) and the error variance in (1.7.6) differ across groups.\\
Model 2: The coefficients are different, but the error variance is the same across groups.\\
Model 3: While each group has common coefficients for $\beta_{3}$ and $\beta_{4}$ (price elasticities) and common error variance, it has a different intercept term and a different $\beta_{2}$. Model 3 is what Nerlove called the hypothesis of neutral variations in returns to scale.

For Model 1, the coefficients and error variances specific to groups can be estimated from

$$
\mathbf{y}^{(j)}=\mathbf{X}^{(j)} \boldsymbol{\beta}^{(j)}+\boldsymbol{\varepsilon}^{(j)} \quad(j=1, \ldots, 5),
$$

where $\mathbf{y}^{(j)}(29 \times 1)$ is the vector of the values of the dependent variable for group $j, \mathbf{X}^{(j)}(29 \times 4)$ is the matrix of the values of the four regressors for group $j, \boldsymbol{\beta}^{(j)} (4 \times 1)$ is the coefficient vector for group $j$, and $\boldsymbol{\varepsilon}^{(j)}(29 \times 1)$ is the error vector. The second column of $\mathbf{X}^{(5)}$, for example, is $\log (Q)$ for $i=117, \ldots, 145$. Model 1 assumes conditional homoskedasticity $\left.\mathrm{E}^{\left(\boldsymbol{\varepsilon}^{(j)}\right.} \boldsymbol{\varepsilon}^{(j) \prime} \mid \mathbf{X}^{(j)}\right)=\sigma_{j}^{2} \mathbf{I}_{29}$ within (but not necessarily across) groups.\\
(d) Estimate Model 1 by OLS. How well can you replicate Nerlove's reported results? On the basis of your estimates of $\beta_{2}$, compute the point estimates of returns to scale in each of the five groups. What is the general pattern of estimated scale economies as the level of output increases? What is the general pattern of the estimated error variance as output increases?

Model 2 assumes for Model 1 that $\sigma_{j}^{2}=\sigma^{2}$ for all $j$. This equivariance restriction can be incorporated by stacking vectors and matrices as follows:

$$
\mathbf{y}=\mathbf{X} \boldsymbol{\beta}+\boldsymbol{\varepsilon}
$$

where

\[
\underset{(145 \times 1)}{\mathbf{y}}=\left[\begin{array}{c}
\mathbf{y}^{(1)}  \tag{*}\\
\vdots \\
\mathbf{y}^{(5)}
\end{array}\right], \quad \underset{(145 \times 20)}{\mathbf{X}}=\left[\begin{array}{lll}
\mathbf{X}^{(1)} & & \\
& \ddots & \\
& & \mathbf{X}^{(5)}
\end{array}\right], \quad \underset{(145 \times 1)}{\boldsymbol{\varepsilon}}=\left[\begin{array}{c}
\boldsymbol{\varepsilon}^{(1)} \\
\vdots \\
\boldsymbol{\varepsilon}^{(5)}
\end{array}\right] .
\]

In particular, $\mathbf{X}$ is now a block-diagonal matrix. The equivariance restriction can be expressed as $\mathrm{E}\left(\boldsymbol{\varepsilon} \boldsymbol{\varepsilon}^{\prime} \mid \mathbf{X}\right)=\sigma^{2} \mathbf{I}_{145}$. There are now 20 variables derived from the original four regressors. The 145 dimensional vector corresponding to the second variable, for example, has $\log \left(Q_{1}\right), \ldots, \log \left(Q_{29}\right)$ as the first 29 elements and\\
zeros elsewhere. The vector corresponding to the 6th variable, which represents log output for the second group of firms, has $\log \left(Q_{30}\right), \ldots, \log \left(Q_{58}\right)$ for the 30th through 58th elements and zeros elsewhere, and so on.

The stacking operation needed to form the $\mathbf{y}$ and $\mathbf{X}$ in (*) can be done easily if your computer software is matrix-based. Otherwise, you trick your software into accomplishing the same thing by the use of dummy variables. Define the $j$-th dummy variable as

$$
D_{j i}=\left\{\begin{array}{ll}
1 & \text { if firm } i \text { belongs to the } j \text {-th group, } \\
0 & \text { otherwise, }
\end{array} \quad(i=1, \ldots, 145) .\right.
$$

Then the second regressor is $D_{1 i} \cdot \log \left(Q_{i}\right)$. The 6th variable is $D_{2 i} \cdot \log \left(Q_{i}\right)$, and so forth.\\
(e) Estimate Model 2 by OLS. Verify that the OLS coefficient estimates here are the same as those in (d). Also verify that

$$
\sum_{j=1}^{5} S S R_{j}=S S R,
$$

where $S S R_{j}$ is the $S S R$ from the $j$-th group in your estimation of Model 1 in (d) and SSR is the SSR from Model 2. This agreement is not by accident, i.e., not specific to the present data set. Prove that this agreement for the coefficients and the SSR holds in general, temporarily assuming just two groups without loss of generality. Hint: First show that the coefficient estimate is the same between Model 1 and Model 2. Use formulas (A.4), (A.5), and (A.9) of Appendix A.\\
(f) (Chow test) Model 2 is more general than Model (1.7.6) because the coefficients can differ across groups. Test the null hypothesis that the coefficients are the same across groups. How many equations (restrictions) are in the null hypothesis? This test is sometimes called the Chow test for structural change. Calculate the $p$-value of the $F$-ratio. Hint: This is a linear hypothesis about the coefficients of Model 2. So take Model 2 to be the maintained hypothesis and (1.7.6) to be the restricted model. Use the formula (1.4.11) for the $F$-ratio.

Gauss Tip: If $x$ is the $F$-ratio, the Gauss command cdffc $(x, d f 1, d f 2)$ gives the area to the right of $F$ for the $F$ distribution with $d f 1$ and $d f 2$ degrees of freedom.

TSP Tip: The TSP command to do the same is cdf(f, df1=df1, df2= df2) $x$. An output of TSP's OLS command, OLSQ, is @SSR, which is the SSR for the regression.

RATS Tip: The RATS command is cdf ftest $x d f 1 d f 2$. An output of RATS's OLS command, LINREG, is \%RSS, which is the SSR for the regression.

The restriction in Model 3 that the price elasticities are the same across firm groups can be imposed on Model 2 by applying the dummy variable transformation only to the constant and log output. Thus, there are $12(=2 \times 5+2)$ variables in X. Now X looks like


\begin{align*}
& \mathbf{X}= \\
& {\left[\begin{array}{ccccccc}
1 & \log \left(Q_{1}\right) & & 0 & 0 & \log \left(P L_{1} / P F_{1}\right) & \log \left(P K_{1} / P F_{1}\right) \\
\vdots & \vdots & & \vdots & \vdots & \vdots & \vdots \\
1 & \log \left(Q_{29}\right) & & 0 & 0 & \log \left(P L_{29} / P F_{29}\right) & \log \left(P K_{29} / P F_{29}\right) \\
& & \ddots & & & \vdots & \vdots \\
0 & 0 & & 1 & \log \left(Q_{117}\right) & \log \left(P L_{117} / P F_{117}\right) & \log \left(P K_{117} / P F_{117}\right) \\
\vdots & \vdots & & \vdots & \vdots & \vdots & \vdots \\
0 & 0 & & 1 & \log \left(Q_{145}\right) & \log \left(P L_{145} / P F_{145}\right) & \log \left(P K_{145} / P F_{145}\right)
\end{array}\right]} \tag{**}
\end{align*}


(g) Estimate Model 3. The model is a special case of Model 2, with the hypothesis that the two price elasticities are the same across the five groups. Test the hypothesis at a significance level of 5 percent, assuming normality. (Note: Nerlove's $F$-ratio on p. 183 is wrong.)

As has become clear from the plot of residuals in Figure 1.7, the conditional second moment $\mathrm{E}\left(\varepsilon_{i}^{2} \mid \mathbf{X}\right)$ is likely to depend on log output, which is a violation of the conditional homoskedasticity assumption. This time we do not attempt to test conditional homoskedasticity, because to do so requires large sample theory and is postponed until the next chapter. Instead, we pretend to know the form of the function linking the conditional second moment to log output. The function, specified below, implies that the conditional second moment varies continuously with output, contrary to the three models we have considered above. Also contrary to those models, we assume that the degree of returns to scale varies continuously\\
with output by including the square of log output. ${ }^{25}$ Model 4 is

\section*{Model 4:}
$$
\begin{aligned}
& \log \left(\frac{T C_{i}}{p_{i 3}}\right)= \beta_{1}+\beta_{2} \log \left(Q_{i}\right)+\beta_{3}\left[\log \left(Q_{i}\right)\right]^{2} \\
&+\beta_{4} \log \left(\frac{p_{i 1}}{p_{i 3}}\right)+\beta_{5} \log \left(\frac{p_{i 2}}{p_{i 3}}\right)+\varepsilon_{i} \\
& \mathrm{E}\left(\varepsilon_{i}^{2} \mid \mathbf{X}\right)=\sigma^{2} \cdot\left(0.0565+\frac{2.1377}{Q_{i}}\right) \quad(i=1,2, \ldots, 145)
\end{aligned}
$$

for some unknown $\sigma^{2}$.\\
(h) Estimate Model 4 by weighted least squares on the whole sample of 145 firms. (Be careful about the treatment of the intercept; in the equation after weighting, none of the regressors is a constant.) Plot the residuals. Is there still evidence for conditional homoskedasticity or further nonlinearities?

\section*{MONTE CARLO EXERCISES}
Monte Carlo analysis simulates a large number of samples from the model to study the finite-sample distribution of estimators. In this exercise, we use the technique to confirm the two finite-sample results of the text: the unbiasedness of the OLS coefficient estimator and the distribution of the $t$-ratio. The model is the following simple regression model satisfying Assumptions 1.1-1.5 with $n=32$. The regression equation is


\begin{gather*}
y_{i}=\beta_{1}+\beta_{2} x_{i}+\varepsilon_{i} \quad(i=1,2, \ldots, n) \\
\text { or } \mathbf{y}=\mathbf{1} \cdot \beta_{1}+\mathbf{x} \cdot \beta_{2}+\boldsymbol{\varepsilon}=\mathbf{X} \boldsymbol{\beta}+\boldsymbol{\varepsilon}, \tag{*}
\end{gather*}


where $\mathbf{X}=(\mathbf{1} \vdots)$ and $\boldsymbol{\beta}=\left(\beta_{1}, \beta_{2}\right)^{\prime}$. The model parameters are ( $\beta_{1}, \beta_{2}, \sigma^{2}$ ).\\
As mentioned in the text, a model is a set of joint distributions of $(\mathbf{y}, \mathbf{X})$. We pick a particular joint distribution by specifying the regression model as follows. Set $\beta_{1}=1, \beta_{2}=0.5$, and $\sigma^{2}=1$. The distribution of $\mathbf{x}=\left(x_{1}, x_{2}, \ldots, x_{n}\right)^{\prime}$ is specified by the following AR(1) process:


\begin{equation*}
x_{i}=c+\phi x_{i-1}+\eta_{i} \quad(i=1,2, \ldots, n) \tag{**}
\end{equation*}


\footnotetext{${ }^{25}$ We have derived the log-linear cost function from the Cobb-Douglas production function. Does there exist a production function from which this generalized cost function with a quadratic term in log output can be derived? This is a question of the "integrability" of cost functions and is discussed in detail in Christensen et al. (1973).
}
where $\left\{\eta_{i}\right\}$ is i.i.d. $N(0,1)$ and
$$
x_{0} \sim N\left(\frac{c}{1-\phi}, \frac{1}{1-\phi^{2}}\right), \quad c=2, \quad \phi=0.6
$$

This fixes the joint distribution of $(\mathbf{y}, \mathbf{X})$. From this distribution, a large number of samples will be drawn.

In programming the simulation, the following expression for $\mathbf{x}$ will be useful. Solve the first-order difference equation (**) to obtain

$$
\begin{aligned}
x_{i}= & \phi^{i} x_{0}+\left(1+\phi+\phi^{2}+\cdots+\phi^{i-1}\right) c \\
& +\left(\eta_{i}+\phi \eta_{i-1}+\phi^{2} \eta_{i-2}+\cdots+\phi^{i-1} \eta_{1}\right),
\end{aligned}
$$

or, in matrix notation,


\begin{equation*}
\underset{(n \times 1)}{\mathbf{x}}=\underset{(n \times 1)}{\mathbf{r}} \cdot x_{0}+\underset{(n \times 1)}{\mathbf{d}}+\underset{(n \times n)}{\mathbf{A}} \underset{(n \times 1)}{\boldsymbol{\eta}}, \tag{***}
\end{equation*}


where $\mathbf{d}=\left(d_{1}, d_{2}, \ldots, d_{n}\right)^{\prime}$ and

$$
\begin{gathered}
d_{1}=c, d_{2}=(1+\phi) c, \ldots, d_{i}=\left(1+\phi+\phi^{2}+\cdots+\phi^{i-1}\right) c, \ldots, \\
\mathbf{r}=\left[\begin{array}{c}
\phi \\
\phi^{2} \\
\vdots \\
\phi^{n}
\end{array}\right], \quad \mathbf{A}=\left[\begin{array}{ccccc}
1 & 0 & \ldots & \ldots & 0 \\
\phi & 1 & 0 & \ldots & 0 \\
\phi^{2} & \phi & 1 & \ddots & 0 \\
\vdots & \vdots & \ddots & \ddots & \vdots \\
\phi^{n-1} & \phi^{n-2} & \ldots & \phi & 1
\end{array}\right], \quad \eta=\left[\begin{array}{c}
\eta_{1} \\
\eta_{2} \\
\vdots \\
\eta_{n}
\end{array}\right] .
\end{gathered}
$$

Gauss Tip: To form the $\mathbf{r}$ matrix, use seqm. To form the A matrix, use toeplitz and lowmat.\\
(a) Run two Monte Carlo simulations. The first simulation calculates E(b|x) and the distribution of the $t$-ratio as a distribution conditional on $\mathbf{x}$. A computer program for the first simulation should consist of the following steps.\\
(1) (Generate $\mathbf{x}$ just once) Using the random number generator, draw a vector $\boldsymbol{\eta}$ of $n$ i.i.d. random variables from $N(0,1)$ and $x_{0}$ from $N(c /(1-\phi), 1 / \left(1-\phi^{2}\right)$ ), and calculate $\mathbf{x}$ by (\textit{\textbf{). (Calculation of $\mathbf{x}$ can also be accomplished recursively by (}) with a do loop, but vector operations such as (}**) consume less CPU time than do loops. This becomes a consideration in the second simulation, where $\mathbf{x}$ has to be generated in each replication.)\\
(2) Set a counter to zero. The counter will record the incidence that $|t|> t_{0.025}(n-2)$. Also, set a two-dimensional vector at zero; this vector will be used for calculating the mean of the OLS estimator $\mathbf{b}$ of $\left(\beta_{1}, \beta_{2}\right)^{\prime}$.\\
(3) Start a do loop of a large number of replications (1 million, say). In each replication, do the following.\\
(i) (Generate $\mathbf{y}$ ) Draw an $n$ dimensional vector $\boldsymbol{\varepsilon}$ of $n$ i.i.d. random variables from $N(0,1)$, and calculate $\mathbf{y}=\left(y_{1}, \ldots, y_{n}\right)^{\prime}$ by ( $*$ ). This $\mathbf{y}$ is paired with the same $\mathbf{x}$ from step (1) to form a sample $(\mathbf{y}, \mathbf{x})$.\\
(ii) From the sample, calculate the OLS estimator $\mathbf{b}$ and the $t$-value for $\mathrm{H}_{0}: \beta_{2}=0.5$.\\
(iii) Increase the counter by one if $|t|>t_{0.025}(n-2)$. Also, add $\mathbf{b}$ to the two-dimensional vector.\\
(4) After the do loop, divide the counter by the number of replications to calculate the frequency of rejecting the null. Also, divide the two-dimensional vector that has accumulated $\mathbf{b}$ by the number of replications. It should equal $E(\mathbf{b} \mid \mathbf{x})$ if the number of replications is infinite.

Note that in this first simulation, $\mathbf{x}$ is fixed throughout the do loop for $\mathbf{y}$. The second simulation calculates the unconditional distribution of the $t$-ratio. It should consist of the following steps.\\
(1) Set the counter to zero.\\
(2) Start a do loop of a large number of replications. In each replication, do the following.\\
(i) (Generate $\mathbf{x}$ ) Draw a vector $\boldsymbol{\eta}$ of $n$ i.i.d. random variables from $N(0,1)$ and $x_{0}$ from $N\left(c /(1-\phi), 1 /\left(1-\phi^{2}\right)\right)$, and calculate $\mathbf{x}$ by $(* * *)$.\\
(ii) (Generate $\mathbf{y}$ ) Draw a vector $\boldsymbol{\varepsilon}$ of $n$ i.i.d. random variables from $N(0,1)$, and calculate $\mathbf{y}=\left(y_{1}, \ldots, y_{n}\right)^{\prime}$ by $(*)$.\\
(iii) From a sample ( $\mathbf{y}, \mathbf{x}$ ) thus generated, calculate the $t$-value for $\mathrm{H}_{0}: \beta=$ 0.5 from the sample $(\mathbf{y}, \mathbf{x})$.\\
(iv) Increase the counter by one if $|t|>t_{0.025}(n-2)$.\\
(3) After the do loop, divide the counter by the number of replications.

For the two simulations, verify that, for a sufficiently large number of replications,

\begin{enumerate}
  \item the mean of $\mathbf{b}$ from the first simulation is arbitrarily close to the true value (1, 0.5);
  \item the frequency of rejecting the true hypothesis $\mathrm{H}_{0}$ (the type I error) is arbitrarily close to 5 percent in either simulation.\\
(b) In those two simulations, is the (nonconstant) regressor strictly exogenous? Is the error conditionally homoskedastic?
\end{enumerate}

\section*{ANSWERS TO SELECTED QUESTIONS}
ANALYTICAL EXERCISES

\begin{enumerate}
  \item $(\mathbf{y}-\mathbf{X} \widetilde{\boldsymbol{\beta}})^{\prime}(\mathbf{y}-\mathbf{X} \widetilde{\boldsymbol{\beta}})$\\
$=[(\mathbf{y}-\mathbf{X b})+\mathbf{X}(\mathbf{b}-\widetilde{\boldsymbol{\beta}})]^{\prime}[(\mathbf{y}-\mathbf{X b})+\mathbf{X}(\mathbf{b}-\widetilde{\boldsymbol{\beta}})]$\\
(by the add-and-subtract strategy)\\
$=\left[(\mathbf{y}-\mathbf{X b})^{\prime}+(\mathbf{b}-\widetilde{\boldsymbol{\beta}})^{\prime} \mathbf{X}^{\prime}\right][(\mathbf{y}-\mathbf{X} \mathbf{b})+\mathbf{X}(\mathbf{b}-\widetilde{\boldsymbol{\beta}})]$\\
$=(\mathbf{y}-\mathbf{X b})^{\prime}(\mathbf{y}-\mathbf{X b})+(\mathbf{b}-\widetilde{\boldsymbol{\beta}})^{\prime} \mathbf{X}^{\prime}(\mathbf{y}-\mathbf{X b})$\\
$+(\mathbf{y}-\mathbf{X b})^{\prime} \mathbf{X}(\mathbf{b}-\widetilde{\boldsymbol{\beta}})+(\mathbf{b}-\widetilde{\boldsymbol{\beta}})^{\prime} \mathbf{X}^{\prime} \mathbf{X}(\mathbf{b}-\widetilde{\boldsymbol{\beta}})$\\
$=(\mathbf{y}-\mathbf{X b})^{\prime}(\mathbf{y}-\mathbf{X b})+2(\mathbf{b}-\widetilde{\boldsymbol{\beta}})^{\prime} \mathbf{X}^{\prime}(\mathbf{y}-\mathbf{X b})+(\mathbf{b}-\widetilde{\boldsymbol{\beta}})^{\prime} \mathbf{X}^{\prime} \mathbf{X}(\mathbf{b}-\widetilde{\boldsymbol{\beta}})$\\
(since $(\mathbf{b}-\widetilde{\boldsymbol{\beta}})^{\prime} \mathbf{X}^{\prime}(\mathbf{y}-\mathbf{X b})=(\mathbf{y}-\mathbf{X b})^{\prime} \mathbf{X}(\mathbf{b}-\widetilde{\boldsymbol{\beta}})$ )\\
$=(\mathbf{y}-\mathbf{X b})^{\prime}(\mathbf{y}-\mathbf{X b})+(\mathbf{b}-\widetilde{\boldsymbol{\beta}})^{\prime} \mathbf{X}^{\prime} \mathbf{X}(\mathbf{b}-\widetilde{\boldsymbol{\beta}})$\\
(since $\mathbf{X}^{\prime}(\mathbf{y}-\mathbf{X b})=\mathbf{0}$ by the normal equations)\\
$\geq(\mathbf{y}-\mathbf{X b})^{\prime}(\mathbf{y}-\mathbf{X b})$
\end{enumerate}

$$
\text { (since }(\mathbf{b}-\widetilde{\boldsymbol{\beta}})^{\prime} \mathbf{X}^{\prime} \mathbf{X}(\mathbf{b}-\widetilde{\boldsymbol{\beta}})=\mathbf{z}^{\prime} \mathbf{z}=\sum_{i=1}^{n} z_{i}^{2} \geq 0 \text { where } \mathbf{z} \equiv \mathbf{X}(\mathbf{b}-\widetilde{\boldsymbol{\beta}}) \text { ). }
$$

7a. $\widehat{\boldsymbol{\beta}}_{\mathrm{GLS}}-\boldsymbol{\beta}=\mathbf{A} \boldsymbol{\varepsilon}$ where $\mathbf{A} \equiv\left(\mathbf{X}^{\prime} \mathbf{V}^{-1} \mathbf{X}\right)^{-1} \mathbf{X}^{\prime} \mathbf{V}^{-1}$ and $\mathbf{b}-\widehat{\boldsymbol{\beta}}_{\mathrm{GLS}}=\mathbf{B} \boldsymbol{\varepsilon}$ where $\mathbf{B} \equiv\left(\mathbf{X}^{\prime} \mathbf{X}\right)^{-1} \mathbf{X}^{\prime}-\left(\mathbf{X}^{\prime} \mathbf{V}^{-1} \mathbf{X}\right)^{-1} \mathbf{X}^{\prime} \mathbf{V}^{-1}$. So

$$
\begin{aligned}
& \operatorname{Cov}\left(\widehat{\boldsymbol{\beta}}_{\mathrm{GLS}}-\boldsymbol{\beta}, \mathbf{b}-\widehat{\boldsymbol{\beta}}_{\mathrm{GLS}}\right) \\
& =\operatorname{Cov}(\mathbf{A} \boldsymbol{\varepsilon}, \mathbf{B} \boldsymbol{\varepsilon}) \\
& =\mathbf{A} \operatorname{Var}(\boldsymbol{\varepsilon}) \mathbf{B}^{\prime} \\
& =\sigma^{2} \mathbf{A V B}
\end{aligned}
$$

It is straightforward to show that $\mathbf{A V B}^{\prime}=\mathbf{0}$.

7b. For the choice of $\mathbf{H}$ indicated in the hint,

$$
\operatorname{Var}(\widehat{\boldsymbol{\beta}})-\operatorname{Var}\left(\widehat{\boldsymbol{\beta}}_{\mathrm{GLS}}\right)=-\mathbf{C V}_{q}^{-1} \mathbf{C}^{\prime}
$$

If $\mathbf{C} \neq \mathbf{0}$, then there exists a nonzero vector $\mathbf{z}$ such that $\mathbf{C}^{\prime} \mathbf{z} \equiv \mathbf{v} \neq \mathbf{0}$. For such $\mathbf{z}$,

$$
\mathbf{z}^{\prime}\left[\operatorname{Var}(\widehat{\boldsymbol{\beta}})-\operatorname{Var}\left(\widehat{\boldsymbol{\beta}}_{\mathrm{GLS}}\right)\right] \mathbf{z}=-\mathbf{v}^{\prime} \mathbf{V}_{q}^{-1} \mathbf{v}<0 \quad \text { (since } \mathbf{V}_{q} \text { is positive definite), }
$$

which is a contradiction because $\widehat{\boldsymbol{\beta}}_{\text {GLS }}$ is efficient.

\section*{EMPIRICAL EXERCISES}
(a) Nerlove's description in Appendix B. 4 leads one to believe that he did not include the depreciation rate $\delta$ in his construction of the price of capital.\\
(b) Your estimates should agree with (1.7.7).\\
(c) Our estimates differ from Nerlove's slightly. This would happen even if the data used by Nerlove were the same as those provided to you, because computers in his age were much less precise and had more frequent rounding errors.\\
(d) How well can you replicate Nerlove's reported results? Fairly well. The point estimates of returns to scale in each of the five subsamples are 2.5, 1.5, 1.1, 1.1, and .96. As the level of output increases, the returns to scale decline.\\
(e) Model 2 can be written as $\mathbf{y}=\mathbf{X} \boldsymbol{\beta}+\boldsymbol{\varepsilon}$, where $\mathbf{y}, \mathbf{X}$, and $\boldsymbol{\varepsilon}$ are as in (*). So (setting $j=2$ ),

$$
\mathbf{X}^{\prime} \mathbf{X}=\left[\begin{array}{cc}
\mathbf{X}^{(1) \prime} \mathbf{X}^{(1)} & \mathbf{0} \\
\mathbf{0} & \mathbf{X}^{(2) \prime} \mathbf{X}^{(2)}
\end{array}\right],
$$

which means

$$
\left(\mathbf{X}^{\prime} \mathbf{X}\right)^{-1}=\left[\begin{array}{cc}
\left(\mathbf{X}^{(1) \prime} \mathbf{X}^{(1)}\right)^{-1} & \mathbf{0} \\
\mathbf{0} & \left(\mathbf{X}^{(2) \prime} \mathbf{X}^{(2)}\right)^{-1}
\end{array}\right] .
$$

And

$$
\mathbf{X}^{\prime} \mathbf{y}=\left[\begin{array}{l}
\mathbf{X}^{(1) \prime} \mathbf{y}^{(1)} \\
\mathbf{X}^{(2) \prime} \mathbf{y}^{(2)}
\end{array}\right] .
$$

Therefore,

$$
\left(\mathbf{X}^{\prime} \mathbf{X}\right)^{-1} \mathbf{X}^{\prime} \mathbf{y}=\left[\begin{array}{l}
\left(\mathbf{X}^{(1) \prime} \mathbf{X}^{(1)}\right)^{-1} \mathbf{X}^{(1) \prime} \mathbf{y}^{(1)} \\
\left(\mathbf{X}^{(2) \prime} \mathbf{X}^{(2)}\right)^{-1} \mathbf{X}^{(2) \prime} \mathbf{y}^{(2)}
\end{array}\right]
$$

Thus, the OLS estimate of the coefficient vector for Model 2 is the same as that for Model 1. Since the estimate of the coefficient vector is the same, the sum of squared residuals, too, is the same.\\
(f) The number of restrictions is $16 . K=$ \#coefficients in Model $2=20$. So the two degrees of freedom should be $(16,125) . ~ S S R_{U}=12.262$ and $S S R_{R}=$ 21.640. $F$-ratio $=5.97$ with a $p$-value of 0.0000 . So this can be rejected at any reasonable significance level.\\
(g) $S S R_{U}=12.262$ and $S S R_{R}=12.577$. So $F=.40$ with 8 and 125 degrees of freedom. Its $p$-value is 0.92 . So the restrictions can be accepted at any reasonable significance level. Nerlove's $F$-ratio (see p. 183, 8th line from bottom) is 1.576 .\\
(h) The plot still shows that the conditional second moment is somewhat larger for smaller firms, but now there is no evidence for possible nonlinearities.

\section*{References}
Amemiya, T., 1985, Advanced Econometrics, Cambridge: Harvard University Press.\\
Averch, H., and L. Johnson, 1962, "Behavior of the Firm under Regulatory Constraint," American Economic Review, 52, 1052-1069.\\
Christensen, L., and W. Greene, 1976, "Economies of Scale in US Electric Power Generation," Journal of Political Economy, 84, 655-676.\\
Christensen, L., D. Jorgenson, and L. Lau, 1973, "Transcendental Logarithmic Production Frontiers," Review of Economics and Statistics, 55, 28-45.\\
Davidson, R., and J. MacKinnon, 1993, Estimation and Inference in Econometrics, Oxford: Oxford University Press.\\
DeLong, B., and L. Summers, 1991, "Equipment Investment and Growth," Quarterly Journal of Economics, 99, 28-45.\\
Engle, R., D. Hendry, and J.-F. Richards, 1983, "Exogeneity," Econometrica, 51, 277-304.\\
Federal Power Commission, 1956, Statistics of Electric Utilities in the United States, 1955, Class A and B Privately Owned Companies, Washington, D.C.\\
Jorgenson, D., 1963, "Capital Theory and Investment Behavior," American Economic Review, 53, 247-259.\\
Koopmans, T., and W. Hood, 1953, "The Estimation of Simultaneous Linear Economic Relationships," in W. Hood, and T. Koopmans (eds.), Studies in Econometric Method, New Haven: Yale University Press.

Krasker, W., E. Kuh, and R. Welsch, 1983, "Estimation for Dirty Data and Flawed Models," Chapter 11 in Z. Griliches, and M. Intriligator (eds.), Handbook of Econometrics, Volume 1, Amsterdam: North-Holland.\\
Nerlove, M., 1963, "Returns to Scale in Electricity Supply," in C. Christ (ed.), Measurement in Economics: Studies in Mathematical Economics and Econometrics in Memory of Yehuda Grunfeld, Stanford: Stanford University Press.\\
Rao, C. R., 1973, Linear Satistical Inference and Its Applications (2d ed.), New York: Wiley.\\
Scheffe, H., 1959, The Analysis of Variance, New York: Wiley.\\
Wolak, F., 1994, "An Econometric Analysis of the Asymmetric Information, Regulator-Utility Interaction," Annales D'Economie et de Statistique, 34, 13-69.


\end{document}
